\section{A. The untwisted limit}
\subsection{A.1. The nodal plumbing}

We write $L=\mathcal O_S(P-O)$, $T=t_{6P}$,
$J=L\otimes T^*L^{-1}=\pi^*\mathcal O_B(1)$, and
$R=\mathbb P_S(\mathcal O\oplus L)$. The convention for this projective
bundle is that its zero and infinity sections have normal bundles
$L$ and $L^{-1}$. The identification defining $Q$ is
$s_0(x)\sim s_\infty(Tx)$. The equality defining $J$ is the equality
in Engel's Lemma 2.1, with this direction of $T$; its intersection
calculation uses $P\cdot O=0$, $P^2=O^2=-1$,
$6P\cdot O=6P\cdot P=2$. Thus the degree on $O$ is
$(O-P)\cdot6P+(P-O)\cdot O=1$. We use the section
$\phi\in H^0(S,J)$ whose divisor is $S_\infty$, as in his
Corollary 2.2. The source being checked is the attached version of
\cite{E26}, not a later inference from its conjecture.

\begin{proposition}[A.1.1; Proved]\label{push:A11}
The cones $P_{k,l}$ in Lead 2 and their faces form a fan. Its two
period automorphisms on the cocharacter lattice are
\[
 G_1(y_1,y_2,h_q,h_u)=(y_1+h_q,y_2,h_q,h_u),\qquad
 G_2(y_1,y_2,h_q,h_u)=(y_1,y_2+h_q+h_u,h_q,h_u).
\]
The affine chart of $P_{0,0}$ is
\[
 U_{0,0}=\operatorname{Spec}\mathbb C[x,y,z,v,u]/(xy-uzv),
 \qquad q=zv.
\]
In particular the unsubdivided chart is not a smooth simultaneous
model of the cusp.
\end{proposition}
\begin{proof}
At fixed $h_q,h_u\geq0$, the cone $P_{k,l}$ cuts out the rectangle
\[
 kh_q\leq y_1\leq(k+1)h_q,\qquad
 l(h_q+h_u)\leq y_2\leq(l+1)(h_q+h_u).
\]
When $h_q=0$ the first interval is a point. These rectangles tile
their support, and their intersections are common faces, including
this collapsed boundary. This proves the fan assertion. The displayed
integral transformations preserve the fan and have integral inverses.

For a character $(a,b,c,d)$, the six inequalities of the dual cone
are $c,a+c,b+c,a+b+c,d,b+d\geq0$. The characters
\[
 x=(0,1,0,0),\quad y=(0,-1,1,1),\quad
 z=(1,0,0,0),\quad v=(-1,0,1,0),\quad u=(0,0,0,1)
\]
generate its integral monoid. Indeed if $b\geq0$ we subtract $b x$;
if $b<0$ we subtract $(-b)y$. The remaining character has second
coordinate zero and is a nonnegative integral combination of $z,v,u$.
The sole relation is $x+y=u+z+v$, and the character of the first
base coordinate is $z+v$. This proves the coordinate-ring assertion,
including normality. The calculation is checked exactly in
\texttt{code/ab\_push\_local\_checks.py}.
\end{proof}

\begin{theorem}[A.1.2; Proved]\label{push:A12}
There is a reduced complex space $\mathcal X^\flat$ and a proper flat
holomorphic map
\[
 (f,\varpi):\mathcal X^\flat\longrightarrow B\times\Delta_u
\]
whose central threefold is $Q$. For $u\ne0$ its fibre is the
contraction of Engel's chosen $X(u)$ at one rational curve of the
Mumford fibre. This curve is the double curve corresponding to the
diagonal of the square in the chosen $A_2$ triangulation; in the
return's coordinates its direction is $e_1-e_2$. Its normal bundle
in $X(u)$ is $\mathcal O_{\mathbb P^1}(-1)^{\oplus2}$. Its image is
the unique ordinary double point of $\mathcal X^\flat_u$.
\end{theorem}
\begin{proof}
Here is a construction on the whole compact surface $S$. Choose a
Hermitian metric on $L$, use the dual metric on $L^{-1}$, and choose
$0<r<1/4$. For sufficiently small $\epsilon>0$ we have
$|u|\|\phi_x\|<r^2$ for all $x\in S$ and $|u|<\epsilon$.
Take the core
\[
 C=\{(x,a,u):a\in L_x,\ r<\|a\|<r^{-1},\ |u|<\epsilon\}
\]
and the node space
\[
 N=\{(x,a,b,u):a\in L_x,\ b\in L^{-1}_{Tx},\
 a\otimes b=u\phi_x,\ \|a\|,\|b\|<2r,\ |u|<\epsilon\}.
\]
On $r<\|a\|<2r$ identify a point of $N$ with $(x,a,u)\in C$.
On $r<\|b\|<2r$ identify it with $(Tx,b^{-1},u)\in C$.
Both identifications are biholomorphisms between open complex spaces;
in the first collar $b=u\phi_x/a$, and in the second
$a=u\phi_xb^{-1}$. Their domains are disjoint because
$\|a\|\|b\|<r^2$. Changes of frames multiply the equation by
units and give the stated bundle-valued equation, so no choice of
frame enters the gluing.

This gluing is Hausdorff. At either end of a collar, exactly one
of its two charts contains the limiting circle: the core omits
$\|a\|=r$, and the node chart omits $\|a\|=2r$; the infinity
collar has the identical property in $b$. Thus the graph of the
identification has no unpaired pair of limiting points. The two
collars in each chart are separated, so there are no additional
chains of identifications. These observations give disjoint
saturated neighbourhoods of any two inequivalent points. Properness
is equally concrete. Over a compact subset of $B\times\Delta_\epsilon$,
the closed core $3r/2\leq\|a\|\leq2/(3r)$ together with the
closed node subspace $\|a\|,\|b\|\leq7r/4$ is compact and maps
onto the glued space. To see surjectivity, a point outside the smaller
core lies on a node collar, and a node point outside the smaller
node subspace lies on the smaller core. Compactness of $S$ supplies
the uniform bound used above. This proves properness of the joint
map to $(t,u)$.

The total node rings are hypersurface rings over the regular
two-dimensional local rings of $S$. They are Cohen--Macaulay.
Every fibre of the joint map has pure dimension two: first fix
$t$, and then impose $xy=u\phi$ on the two line coordinates over
the one-dimensional fibre of $S$. Consequently $t-t_0,u-u_0$ is a
regular sequence in every total local ring. The regular-sequence
flatness criterion over the smooth two-dimensional base proves
flatness. On the core the same conclusion is immediate from
flatness of $S\to B$. This includes the nonreduced Kodaira fibres
at $0,1$; smoothness of $S$, rather than smoothness of those fibres,
is what the argument uses. At $u=0$ the node equation is $ab=0$,
and the two collar maps recover precisely the specified pinching
of the two sections of $R$.

For $u\ne0$ and $t\ne\infty$, the equation gives
$a_{Tx}=a_x/(u\phi_x)$. Hence the resulting quotient is exactly
$L^\times/\langle u\phi\rangle$, where the contracting generator
covers $T^{-1}$. Every orbit meets a core or a collar, by the
uniform contraction estimate; the displayed comparison identifies
exactly its orbit relation. It is therefore a biholomorphism,
not just a meromorphic comparison.

At the nodal point of $S_\infty$, write $q=zv$. Absorbing a unit
in one line coordinate gives $xy=uzv$. At a smooth point of
$S_\infty$ the equation is $xy=uq$. For $u\ne0$ the latter
threefold is smooth; the former has one nondegenerate quadratic
singularity at the origin. There are no other singularities of a
nonzero $u$ fibre.

For completeness, the relation to the infinite fan quotient is
also explicit. In a Tate covering of a neighbourhood of the $I_1$
fibre use the torus coordinates $(Z,W)$ of Engel's Proposition 2.7.
The two periods act as
\[
 (Z,W,q,u)\longmapsto(c_1qZ,c_2W,q,u),\qquad
 (Z,W,q,u)\longmapsto(c_3Z,c'_4uqW,q,u),
\]
with holomorphic units $c_i$. Their monomial parts are $G_1,G_2$.
The chart indexed by $(k,l)$ has
$z=Zq^{-k}$, $v=q^{k+1}/Z$, $x=W(qu)^{-l}$,
$y=qu/x$, with the harmless unit changes supplied by the line
frames. It is the chart of Proposition~\ref{push:A11}.
The Tate cover unfolds the node of $S_\infty$ into its infinite
chain, and unfolding the two plumbing collars gives the additional
integer index $l$. The foregoing collar identifications are exactly
the two deck identifications. Thus, on an invariant neighbourhood
over a sufficiently small bidisc, the fan space maps by a
holomorphic covering to the plumbing, with deck group $\mathbb Z^2$.
Each core or node coordinate neighbourhood has disjoint lifts
indexed by the two unfolded integers; this proves freeness and
proper discontinuity, including the $\rho_l$ orbits. The first
period stabilizes such a ray but on its orbit translates the Tate
coordinate by $c_1q$, with $q\ne0$ and $|c_1q|<1$; it does not
fix a point. Properness of its quotient was proved by the compact
core and node representatives. This supplies the analytic
comparison needed here. Construction 3.8 of \cite{EGFS25} is a
related multivariable polarized construction; its symmetry
hypothesis on period units is not being silently imposed on these
generally unpolarized tori.

Subdividing the square at $h_u=0$ by the chosen diagonal gives
the usual small resolution of $xy=uzv$ for $u\ne0$. Its exceptional
curve has the two affine charts of the total space of
$\mathcal O_{\mathbb P^1}(-1)\oplus
\mathcal O_{\mathbb P^1}(-1)$, as follows by setting $x=za$ on
one chart and $z=xb$ on the other. The two normal coordinates
transform by the reciprocal affine coordinate. There is one
diagonal edge modulo the periods, so there is exactly one such
curve. Removing that edge contracts the paired opposite
$(-1)$-curves in the normalization $dP_6$ to the square model
$\mathbb P^1\times\mathbb P^1$. This identifies the contracted
curve intrinsically in the selected toric model. It is not the
$e_2$-direction curve fixed by $\operatorname{Aut}^0$.
\end{proof}

\subsection{A.2. Two minimal toroidal fillings and further modifications}

\begin{proposition}[A.2.1; Proved]\label{push:A21}
Among subdivisions with no new rays, the fan in A.1.1 has exactly
two smooth period-invariant subdivisions restricting to the chosen
$A_2$ triangulation at $h_u=0$. These are also all subdivisions
which are flat over the original $(q,u)$ bidisc, have reduced
$u$-central divisor, and retain that boundary triangulation.
Without the joint-flatness restriction there are infinitely many
smooth invariant subdivisions retaining that triangulation and
having reduced $u$-central divisor.
\end{proposition}
\begin{proof}
Write the six rays of $P_{0,0}$ as
$A=(0,0,1,0)$, $B=(1,0,1,0)$, $C=(0,1,1,0)$,
$D=(1,1,1,0)$, $E=(0,0,0,1)$, $F=(0,1,0,1)$.
The return's $A_2$ diagonal is $BC$. A triangular prism has six
triangulations without new vertices. With this bottom diagonal
the three possibilities are
\[
\begin{array}{c|c|c}
\text{maximal tetrahedra}&\text{left diagonal}&\text{right diagonal}\\\hline
ABCE,\ BCDE,\ CDEF&CE&DE\\
ABCE,\ BCDF,\ BCEF&CE&BF\\
ABCF,\ ABEF,\ BCDF&AF&BF.
\end{array}
\]
Periodicity identifies the left and right faces by $A\mapsto B$,
$C\mapsto D$, $E\mapsto E$, $F\mapsto F$; it excludes exactly
the middle row. Every listed maximal cone has determinant $\pm1$.
There are no missed triangulations: the prism has normalized volume
three, all its nondegenerate four-vertex tetrahedra have normalized
volume one, and enumeration of their triples with pairwise
face intersections gives exactly these six possibilities before
fixing the bottom. The program \texttt{ab\_push\_fan.py} performs
this enumeration in rational arithmetic; intersections are tested
by all vertices of the combined barycentric halfspaces.

For the stronger assertion, a new ray with $h_q,h_u>0$ would
give a divisor mapping to the origin of the two-dimensional base.
Its fibre dimension would be three, contradicting joint flatness.
A ray with $h_u=0$ is excluded by the retained boundary
triangulation. On $h_q=0$, the fan support has $y_1=0$.
Reducedness requires the primitive ray to have $h_u=1$, so it
is already one of the $\rho_l$. Thus there are no new rays under
the stated three restrictions. Conversely the two triangulations
have all heights zero or one, and every cone maps onto a face of
the quadrant with saturated image. Their local monomial maps
are flat and their $u$ divisors reduced.

If joint flatness is omitted, star-subdivide the orbit of the
maximal cone $ABCE$ by its interior ray
$v_1=A+B+C+E=(1,1,3,1)$. Distinct translates have disjoint
interiors and distinct closed torus orbits. These are simultaneous
ordinary blow-ups at disjoint points, so the operation is
unambiguously period-invariant. Repeat in the cone $ABCv_{m-1}$,
where $v_0=E$. The new ray is
$v_m=v_{m-1}+A+B+C=(m,m,3m,1)$. Replacing any basis vector
by their sum preserves determinant one, and every new ray has
$h_u=1$. The boundary triangulations are unchanged. On the
compact quotient this blows up one smooth central point away
from the conductor at each step. There are therefore infinitely
many invariant smooth refinements, all reduced over $u=0$, with
increasing numbers of rays modulo the period group. They are
not jointly flat. An earlier attempted witness using an orbit
of two-faces did not justify compatibility of intersecting
centres; its correction is retained in
\path{evidence/ab_push_correction_orbit_stars.md}.
\end{proof}

\begin{theorem}[A.2.2; Proved]\label{push:A22}
The two minimal subdivisions of A.2.1 give proper flat holomorphic
families $\mathcal X^+,\mathcal X^-\to\Delta_u$ with smooth total
spaces. Their joint maps to $B\times\Delta_u$ are flat. Over
$\Delta_u^*$ they are isomorphic as families to the specified
untwisted family $X(u)$. Each central fibre $Q^\pm$ is reduced,
irreducible, and has normal crossings, with smooth double locus
$S$ and no triple points.

Its normalization $R^\pm$ is obtained by blowing up $R$ along
$S_\infty$ in, respectively, one of its two sections and then
taking the small resolution at the unique ordinary double point
of that blow-up prescribed by the fixed bottom diagonal.
The two proper transforms of the sections are still copies of
$S$, glued by $T$. Their normal-bundle product is trivial;
in particular $\mathcal T^1_{Q^\pm}\simeq\mathcal O_S$.
\end{theorem}
\begin{proof}
Apply either fan subdivision over the cusp neighbourhood in
Theorem~\ref{push:A12}, and use the identity modification elsewhere.
Toric modifications are proper, and a subdivision is an
isomorphism off the strata it subdivides. Its restriction over
$q\ne0$ is the identity, so this gluing has no additional
transition choice. The period actions preserve the subdivisions,
and the covering description in A.1.2 descends the proper
modifications to the quotient. All maximal cones are regular.
Every local ring of the resulting fourfold is therefore regular.
The coordinate $u$ is the product of those affine coordinates
whose rays have $h_u=1$; there are at most two such rays in a
maximal cone. This proves reduced normal crossings and the absence
of triple points for the $u$-central fibre. The regular-sequence
argument, or the saturated face images in A.2.1, proves joint
flatness. Properness follows from properness of the toric
modification and of $\mathcal X^\flat$.

We give the normalization calculation, including the singular
centre. At a node of $S_\infty$, write the original normalization
as coordinates $(X,z,v)$, with section $X=0$ and centre
$(X,zv)$. Blowing up that ideal has charts
\[
 X=zvA;\qquad zv=XB.
\]
The first is smooth, and the second has one conifold at its
origin. A small resolution of the latter is, for one of the
bottom-diagonal conventions,
\[
 X=za,\quad v=aB;\qquad z=Xd,\quad B=dv.
\]
Interchanging $z,v$ gives the other diagonal convention. These
are holomorphic charts of a proper small modification: they are
the two usual charts of the blow-up of the appropriate rank-one
ideal on the conifold. The strict transform of $X=0$ is $A=0$
in the first chart and never passes through the conifold. It is
therefore unchanged as a surface. The other section is disjoint
from the centre and is unchanged as well. Over a smooth point
of $S_\infty$ the exceptional fibre is one $\mathbb P^1$;
over its node it is a chain of two $\mathbb P^1$'s, one being
the exceptional curve of the small resolution. This identifies
the normalization morphism, including its order of operations.

For example the corresponding total fourfold modification of
$xy=uzv$ can be written by first taking $x=ua$, $ay=zv$ and
then resolving the conifold in $(a,y,z,v)$, together with the
other chart $u=xb$, $y=bzv$. It has three smooth affine charts.
The assignments in \texttt{ab\_push\_local\_checks.py} verify
the equations and the map to $q=zv$ exactly. Exchanging the
chosen zero and infinity branches gives the other model.

The normal bundle of the section containing the centre changes
by tensoring with $\mathcal O_S(-S_\infty)$; the second normal
bundle is unchanged. Since the original product is $J$, the
new product is $J(-S_\infty)=\mathcal O_S$. The local
formula for $\mathcal T^1$ is the product of the two normal
bundles, with $T$ used on the infinity branch. This proves
d-semistability in this precise sense. The $\rho_l$ form one
orbit, so the central fibre is irreducible and its conductor is
exactly the two stated sections. Over $u\ne0$ both subdivisions
restrict to the fixed bottom triangulation, giving precisely
Engel's Mumford filling, simultaneously in $u$.

The birational comparison $\mathcal X^+\dashrightarrow
\mathcal X^-$ is the identity over the punctured $u$ disc and
off the exceptional loci over $(q,u)=(0,0)$. Locally it replaces
the two compatible side-face diagonals by the other two.
All its rays lie in the Gorenstein hyperplane $h_q+h_u=1$,
so both small modifications are crepant over
$\mathcal X^\flat$. Over a smooth point of the nodal curve
this is the ordinary Atiyah flop times a curve; at the node its
extension is exactly the displayed pair of prism subdivisions.
This specifies the birational map without asserting a uniqueness
theorem for arbitrary non-toroidal fillings.
\end{proof}

\begin{remark}[A.2.3; Refuted; the literal blow-up shape]\label{push:A23}
The unrestricted request that every normalized component
be a blow-up of $R$ is false. The first maximal-cone
star already contributes a $\mathbb P^3$ component,
which admits no proper bimeromorphic map to $R$, since
such a map would inject a rank-eleven rational $H^2$
into a rank-one group. A.2.4 gives every normalization
by explicit star charts; A.2.5 supplies the terminating
blow-up and blow-down algorithm with centres, order,
exceptional bundles and gluing maps. Thus the literal
shape is refuted and its general algorithmic
replacement is proved.
\end{remark}

\begin{theorem}[A.2.4; Proved; exhaustive finite fan data]\label{push:A24}
The following finite-data procedure parametrizes all smooth
period-invariant subdivisions of the cusp fan after
$u=r^n$, retaining the prescribed bottom triangulation and
having finitely many cones modulo the periods. Here $n=1$
gives Part A, and $n=2,12$ give the two cases of B.2.
Each output defines a proper flat family over $\Delta_r$
with smooth total space, whose punctured restriction is the
given untwisted family pulled back by $u=r^n$.

For a bound $H\geq1$, retain the four bottom rays
$A,B,C,D$ and consider all primitive vectors
\[
 \begin{gathered}
 v=(a,b,k,\ell),\qquad
 \ell\geq1,\quad k\geq0,\quad k+\ell\leq H,\\
 0\leq a\leq k,\qquad0\leq b\leq k+n\ell.
 \end{gathered}
 \tag{A.24a}
\]
Take any collection of four-element subsets of this finite
set satisfying the following exact tests:
\begin{enumerate}
\item Each four-column determinant has absolute value one.
\item After dividing each ray by $k+\ell$, the resulting
tetrahedra intersect pairwise in common faces, and their
volumes sum to $(n+2)/6$.
\item Their bottom boundary is the two triangles $ABC,BCD$.
Their boundary subdivisions on $a=0$ and $a=k$ agree under
$G_1$, and those on $b=0$ and $b=k+n\ell$ agree under $G_2$.
\end{enumerate}
All these are finite rational-arithmetic tests. Extend a
passing collection by the two periods, and discard duplicates
as $H$ increases. This is exhaustive. For a reduced central
divisor, retain only $\ell=1$ in (A.24a).

The central prime components are indexed by period orbits of
the selected rays with $\ell>0$, and the multiplicity of the
component indexed by $v$ is exactly $\ell$. The normalization
charts of that component are the star fan of $v$ in
$N/\mathbb Zv$, with the inherited bundle units and the
stabilizer action. For $k>0$ the stabilizer is trivial and this
normalization is a smooth complete toric threefold. All
multiple loci, including possible triple and quadruple loci,
are specified by the positive-$\ell$ rays belonging to each
cone. In the reduced case
$\mathcal T^1\simeq\mathcal O_{\operatorname{Sing}D}$,
so every output is d-semistable.
\end{theorem}
\begin{proof}
Write the period shears as
\begin{align*}
 G_1(a,b,k,\ell)&=(a+k,b,k,\ell),\\
 G_2(a,b,k,\ell)&=(a,b+k+n\ell,k,\ell).
\end{align*}
One closed fundamental cone is
\[
 P_n=\{k,\ell\geq0,\ 0\leq a\leq k,
                    \ 0\leq b\leq k+n\ell\}.
\]
Its bottom is the original square; its top is the edge from
$(0,0,0,1)$ to $(0,n,0,1)$. The slice $k+\ell=1$, in
coordinates $(a,b,k)$, is the compact polytope
\[
 0\leq k\leq1,\qquad0\leq a\leq k,
 \qquad0\leq b\leq n-(n-1)k.
\]
Its volume is
$\int_0^1 k(n-(n-1)k)\,dk=(n+2)/6$.
A tetrahedron with ray columns $v_i$ has slice volume
\[
 \frac{|\det(v_1,v_2,v_3,v_4)|}
      {6\prod_{i=1}^4(k_i+\ell_i)}.
 \tag{A.24b}
\]
This follows by applying the determinant-one coordinate
change from $(a,b,k,\ell)$ to $(a,b,k,k+\ell)$ and then
normalizing the last coordinate to one.

A common-face test is also completely explicit: write each
tetrahedron in barycentric halfspaces, intersect the two
finite lists of halfspaces, and enumerate the possible
vertices by solving triples of bounding equations. The
intersection is a common face precisely when all its vertices
are common tetrahedron vertices. Since the candidate
tetrahedra lie in $P_n$, pairwise face intersections and the
volume equality imply that they cover the slice. Otherwise
their finite closed union would omit a nonempty relatively
open set, which has positive volume. The boundary tests make
the translates agree on every shared face. This proves that
every passing collection gives the claimed invariant fan.

Conversely, intersect any invariant smooth refinement with
$P_n$. A smooth maximal cone has four primitive generators
with determinant $\pm1$, and its slice is one of the stated
tetrahedra. Finitely many cone orbits give finitely many rays
in this prism, hence some maximum height $H$. No additional
bottom rays are permitted by the boundary condition; every
other ray has the form (A.24a). The resulting collection
passes all three tests. Thus the procedure is a finite,
exhaustive search at every height, rather than a restriction
to regular triangulations or to new rays in a chosen orbit.
There is no projectivity condition in the tests.

The finite-orbit condition is necessary for a proper quotient
of such a refinement, not an additional restriction on the
requested proper models. Each full-dimensional cone has a
closed torus orbit consisting of one point with $q=r=0$.
These points form a closed discrete set: an affine chart
of a full-dimensional cone contains exactly its own closed
torus orbit. After the properly discontinuous period
quotient the points are indexed by cone orbits and remain
closed and discrete. Infinitely many such orbits would
therefore give an infinite closed discrete subset of a
compact central fibre, which is impossible.

The fundamental orbit representatives of mixed rays are
\[
 (a,b,k,\ell),\quad k,\ell>0,
 \quad0\leq a<k,\quad0\leq b<k+n\ell,
 \quad\gcd(a,b,k,\ell)=1.
 \tag{A.24c}
\]
For pure top rays they are $(0,b,0,\ell)$ with
$0\leq b<n\ell$ and $\gcd(b,\ell)=1$. A smooth top
subdivision necessarily includes each $(0,j,0,1)$,
$0\leq j\leq n$. Indeed an omitted such vector would lie
inside a regular two-cone; writing it in that cone's integral
basis would express its height one as the sum of two
positive integral heights, which is impossible.
In the reduced search there are exactly $n$ pure top orbits
and $\sum_{k=1}^{H-1}k(k+n)$ available mixed orbits.

The program \path{code/ab_push_all_refinements.py} enumerates
all candidate rays for $n=1,2,12$ through $H=3$, and all
four-column determinant tests at the bounds reported here:
\[
\begin{array}{c|c|c|r|r}
n&H&\text{reduced}&\text{rays in }P_n&
       \text{unimodular four-subsets}\\\hline
1&1&\text{yes}&6&12\\
1&2&\text{yes}&12&267\\
1&2&\text{no}&13&357\\
2&2&\text{yes}&15&544\\
2&2&\text{no}&17&784\\
12&1&\text{yes}&17&100
\end{array}
\]
These are candidate counts, not a claim that every candidate
four-subset or collection passes the other tests. At height
one, the complete collections are the two and $n+1$ fans
already proved in A.2.1 and B.2.1. The saved program also
certifies the first eight members of the explicit infinite
family in A.2.1.

For existence of the corresponding compact family, start
with the proper nodal plumbing of A.1.2 and base-change it
by $u=r^n$. On the cusp cover apply the proper toric
modification defined by the refinement. Its support is
unchanged, which is the toric properness criterion
\cite[Section 2, p.6]{DMM16}, and each
original affine chart is replaced by a finite fan. On
$q\ne0$ use the subdivision of the top two-dimensional fan
to modify the chart $ab=r^n\varphi$ fibrewise over $S$.
Changes of trivialization multiply $a$ and $b$ by the given
nonvanishing line-bundle units; a torus action extends to
every toric fan chart, so these modifications glue. The same
top fan occurs on the overlap with the cusp construction.
The core outside the plumbing charts is unchanged. This defines
an actual proper modification of the base-changed compact
nodal family, not a gluing inferred from a vanishing theorem.

The period maps lift because the fan is invariant. Their
lifted action is free: a fixed point would project to a fixed
point in the nodal covering. It is properly discontinuous:
compact sets project to compact sets, and any translating
element producing an intersection upstairs produces one
downstairs, where A.1.2 gives finiteness. Thus the quotient
exists as a Hausdorff complex space, and the equivariant
proper modification descends to a proper modification of
the compact nodal quotient. In a unimodular chart its
coordinates $z_i$ satisfy
\[
 r=\text{unit}\cdot\prod_{i=1}^4 z_i^{\ell_i},
 \qquad q=\text{unit}\cdot\prod_{i=1}^4 z_i^{k_i}.
 \tag{A.24d}
\]
The total space is smooth. The nonconstant coordinate $r$
is a nonzerodivisor, so the morphism to the disc is flat.
Its central multiplicities are the indicated $\ell_i$;
it is reduced exactly when all selected positive heights
are one. The bottom fan gives the specified smooth
Mumford model over $r\ne0$, including its original period
identification. Thus the punctured-family assertion holds
simultaneously in the parameter.

For completeness the normalization description is an
explicit chart prescription. For a maximal cone with basis
$(v,v_2,v_3,v_4)$, let $(m_1,m_2,m_3,m_4)$ be the dual
integral basis. The chart of the normalized component is
$z_1=0$ with coordinates $z_i=\chi^{m_i}$, $i=2,3,4$.
On overlaps the transition matrix is the product of the two
integer basis matrices, and the original bundle units are
raised to those same integer exponents. This is precisely
the star fan in $N/\mathbb Zv$. When $k>0$, (A.24c) shows
that the period stabilizer of the ray is trivial. The ray
lies in the interior of the full support, so its star is
complete. Its normalized component is therefore the
asserted complete smooth toric threefold, lying over the
cusp. When $k=0$, the stabilizer is generated by $G_1$.
Away from $S_\infty$ the resulting component is the ruled
bundle $\mathbb P_S(\mathcal O\oplus L^\ell)$: a character
on its open ruled fibre is $W^\ell/r^b$, which transforms
by the $\ell$th power of the transition of $L$. The stated
star charts describe its replacement over $S_\infty$.
This also specifies all gluing maps, rather than only their
combinatorial incidence: they are the original period maps
on those same integer characters.

Intersections are read from common cones, with separate
local branches retained before the period quotient. In a
chart (A.24d), the number of branches of the reduced
central divisor is the number of positive $\ell_i$.
When all these heights are one, the hypersurface
calculation gives
$\mathcal T^1_D=\mathcal O_D(D)|_{\operatorname{Sing}D}$.
The divisor is globally principal, cut out by $r$, so this
line on the singular locus is trivial. Finally, any two
output fans have an explicit common refinement by all
intersections of their cones. Its normal toric modification
gives a proper common model, and the resulting birational
map is the identity on the fixed lattice and on the
punctured family. No meromorphic extension of an arbitrary
punctured isomorphism is needed here.
\end{proof}

\begin{proposition}[A.2.5; Proved; explicit normalization algorithm and its necessary scope]\label{push:A25}
The requirement that every normalized component of every
smooth refinement be a blow-up of
$R=\mathbb P_S(\mathcal O\oplus L)$ is false, already
for a reduced refinement with $n=1$.
The following is a uniform replacement for any finite
fan certificate $\Sigma$ in A.2.4.

The star charts and integer transition matrices of A.2.4
give every normalized component, its gluing maps, and
all multiple loci directly. There is also a terminating
finite-data algorithm that expresses its comparison
with a reference model as an explicit sequence of
blow-ups and blow-downs along specified smooth strata.
At each blow-up the algorithm gives the normalized
strict transforms, the exceptional projective bundle,
its multiplicity, and every new intersection.
This algorithm allows blow-downs; no toric
strong-factorization assertion is used.

For a component of multiplicity $\ell$, the local
nonreduced central equation is the monomial in A.24d.
More generally, in a unimodular chart with central
equation $F=\prod_i z_i^{\ell_i}$, the cotangent-complex
first-order deformation sheaf is
\[
 \mathcal T^1_D
  =\mathcal O_D\Big/
    \left(\ell_i z_i^{\ell_i-1}
          \prod_{j\ne i}z_j^{\ell_j}:\ell_i>0\right).
 \tag{A.25a}
\]
Here $\mathcal T^1_D$ means
$\mathcal Ext^1(\mathbb L_D,\mathcal O_D)$; in the
reduced case it also equals
$\mathcal Ext^1(\Omega_D,\mathcal O_D)$.
The reduced case is the trivial line on the singular
scheme, as in A.2.4. For a nonreduced centre the displayed
Jacobian quotient is the appropriate deformation sheaf;
the reduced normal-crossings notion of d-semistability
is not asserted for it.
\end{proposition}
\begin{proof}
The first maximal-cone star in A.4.6 produces a normalized
component $\mathbb P^3$. There cannot be a proper
bimeromorphic holomorphic map $\mathbb P^3\to R$:
between smooth compact oriented manifolds such a map
has degree one, and proper pushforward gives a left
inverse to pullback on rational cohomology. But
$b_2(R)=11$ and $b_2(\mathbb P^3)=1$.
Thus this component is not a blow-up of $R$, regardless
of the proposed centres or their order. The actual
normalization in this example is
$\operatorname{Bl}_pR^+\sqcup\mathbb P^3$,
with its specified plane intersection. This is an
explicit counterexample to the universal shape, not
a failure to find a blow-up sequence.

Here is the promised algorithm. Retain all marking
data from A.2.4: the lattice, the two period matrices,
the original line-bundle units, and the common
unsubdivided fan. Strictify the finite quotient cone
complex by the decreasing-dimensional barycentric
stars used in A.4.8. Record these operations and do
the same for the reference model. Form their common
intersection complex. A regular common subdivision
can be found either by toric desingularization or by
enumerating rational subdivisions with increasing
primitive-ray bounds and testing their determinants
and common faces. A regular subdivision exists, so
the latter enumeration terminates.

On a finite strict regular complex there are finitely
many possible ordinary star blow-ups: choose a cone
of dimension at least two and insert the sum of its
primitive basis rays. There are also finitely many
possible inverse moves: choose a ray and test whether
its star is exactly the star created by such a
subdivision of a regular cone. Both tests are exact
integer calculations. Require that every intermediate
complex remain a subdivision of the marked original
complex. Record face identifications; equality of two
marked complexes is a finite rational comparison,
using the period orbit representatives in A.24c and
their pure top and pure bottom analogues.

Enumerate this locally finite graph of moves in
increasing path length, discarding repeated marked
complexes. Applied separately to the two proper
subdivision morphisms from the common regular
complex, the finite-complex weak-factorization theorem
used in A.4.8 guarantees a path, all of whose
intermediates are over the corresponding target.
Thus the search terminates. Concatenate the two paths
and the recorded strictification paths. This gives a
fully specified finite sequence once the input fan is
given. Temporary pure bottom subdivisions may occur;
they are allowed in a presentation of the birational
comparison, and are removed before the prescribed
endpoint is reached.

We describe explicitly what the recorded move means
geometrically. If a cone has primitive basis vectors
$v_1,\ldots,v_c$, its centre is the closed stratum
\[
 C=\{z_1=\cdots=z_c=0\}
\]
in the marked toroidal chart. Strictness makes its
global closure smooth; the period translates defining
one centre are disjoint. Blow-up chart $i$ has
coordinates
\[
 z_i=t_i,\qquad z_j=t_i w_j\quad
 (1\leq j\leq c,\ j\ne i),
 \qquad z_j\quad(j>c).
\]
The changes between charts are these integral
monomials, with the original units raised to the
same powers. The exceptional divisor is
$\mathbb P_C(N_{C/Y})$, where $Y$ is the smooth
intermediate total space. If $C$ is central, its
central multiplicity is
\[
 \ell_E=\sum_{j=1}^c\ell_j.
\]
If all $\ell_j$ vanish, it is a horizontal exceptional
divisor and is not a new central component.

For an old normalized component, take the inverse
image of this centre on its marked branch and the
ordinary strict transform in these charts. It is
the blow-up along that smooth inverse-image centre;
when that centre is Cartier the blow-up is the
identity. Its intersections with the new exceptional
component are the appropriate coordinate-hyperplane
projective subbundles. Intersections of several old
branches are given by imposing the corresponding
coordinate equations simultaneously. Thus the same
chart prescription records double, triple and
quadruple loci, their normal lines, and their actual
gluing maps. An inverse star is the inverse
contraction of this recorded configuration. This
describes centres and order, not merely a claim that
some birational map exists.

The algorithm starts with the explicit reference
normalizations of A.2.2 and the $n$-cycle reference
in A.4.8. It is therefore a uniform geometric
presentation of every output by specified blow-ups,
blow-downs and projective bundles. Independently,
the direct star charts avoid running this search
when only a normalization chart or an intersection
is required. The mixed-ray components are complete
smooth toric threefolds; a pure top component
$(0,b,0,\ell)$ is bimeromorphic to
$\mathbb P_S(\mathcal O\oplus L^\ell)$.
There is no assertion that all these distinct
components are modifications of a single copy of $R$.

Finally $D$ is globally the principal Cartier
divisor of the parameter $r$ in a smooth total
space. Its two-term cotangent complex is
$[\mathcal O_D\xrightarrow{dF}\Omega_Y|_D]$,
so its first-order deformation module is the
cokernel of the gradient. Equivalently, a lifted
hypersurface equation $F+\epsilon G$ is changed by
an infinitesimal coordinate change precisely by a
Jacobian term. This argument applies to nonreduced
monomials as well; an ordinary conormal sequence
need not be left-exact there. Absorbing the local unit, or reducing its
derivative modulo $F$, gives (A.25a). A global
parameter makes the normal line trivial, so these
Jacobian quotients glue by the canonical
hypersurface presentation. In the reduced case
the Jacobian ideal defines the singular scheme,
giving the asserted trivial deformation line there.

The exact certificate
\path{code/ab_push_normalization_algorithm.py}
executes the star search on the first two
maximal-cone examples, verifies both inverse
contractions, and checks all local multiplicity
substitutions with exponents zero, one and two.
The finite-complex theorem supplies termination for
arbitrary input; these tests check the explicit
local implementation.
\end{proof}


\subsection{A.3. The virtual Euler obstruction and the conjecture}

\begin{lemma}[A.3.1; Proved]\label{push:A31}
Let $Z$ be a compact reduced local-complete-intersection complex
space of pure dimension $n$. In any proper flat holomorphic
family over a disc with central fibre $Z$ and smooth nearby
fibres, their Euler number is
\[
 e_{\mathrm{vir}}(Z)=\langle c_n(T_Z^{\mathrm{vir}}),[Z]\rangle.
\]
This statement needs no projectivity or K\"ahler hypothesis.
\end{lemma}
\begin{proof}
Flat deformation preserves the local-complete-intersection
property relative to the disc near $Z$. Locally this follows by
lifting an embedding and its regular sequence of equations;
Nakayama's lemma and flatness show that their lifts generate
the ideal and remain a regular sequence. Properness allows one
to shrink the disc so that these neighbourhoods cover the total
space. The relative cotangent complex is consequently perfect
of amplitude $[-1,0]$. Its dual defines a topological complex
$K$-class $V$ on that total space; local finite vector-bundle
complexes define the class by descent, or by stabilization and
partitions of unity; the analytic-to-topological comparison is
given in \cite[Lemmas 3.1--3.2 and Section 3.1]{Levy08}.
Its restriction to a smooth fibre is its
tangent bundle, and its restriction to $Z$ is its intrinsic
virtual tangent class.

The compact fundamental cycles of the central and a nearby
fibre are homologous in the total space. More explicitly,
the inverse image $K$ of a short closed radial segment
$[0,\epsilon]$ is compact and semianalytic. Away from
$\operatorname{Sing}Z$ it is an oriented real
$(2n+1)$-manifold with boundary the nearby fibre and
$-Z_{\mathrm{reg}}$, because flatness with smooth fibre
implies smoothness of the morphism at those points.
The exceptional set $\operatorname{Sing}Z$ has real
dimension at most $2n-2$. Triangulate compatibly with
this set and the two endpoints, using
\cite[Section 3, Theorems 1--2]{Loj64}. The sum of the
oriented top simplices has boundary exactly
$[X_\epsilon]-[Z]$: no real $2n$-simplex can be supported
on that exceptional set. Reducedness and flatness give
multiplicity one on every central component. Pairing the globally defined
$c_n(V)$ with these homologous cycles gives the equality of
Chern numbers. Gauss--Bonnet on the smooth compact complex
fibre gives its Euler number. This proves the analytic assertion
directly. The algebraic counterpart is Theorem 5.8 and
Remark 5.9 of \cite{FG10}; those scheme statements alone would
not justify omitting the analytic comparison above.
\end{proof}

\begin{theorem}[A.3.2; Proved; Conjecture 8.2 refuted]\label{push:A32}
For the specified pinching $Q$ one has $e_{\mathrm{vir}}(Q)=0$.
Every smooth proper-flat smoothing of $Q$ therefore has Euler
number zero. In particular there is no proper flat holomorphic
filling of Engel's untwisted family with central fibre exactly
$Q$, even if no map to $B$ is required. Thus Conjecture 8.2
of the attached \cite{E26} is false as worded. The proper nodal
filling of A.1.2 and the two natural semistable fillings of
A.2.2 are precise replacements, with different central fibres.
\end{theorem}
\begin{proof}
Let $D_0,D_\infty\subset R$ be the conductor sections and
$\nu:R\to Q$ the normalization. On either branch of the
hypersurface $xy=0$, the derivative comparison with its
normalization gives in $K^0(R)$, where pullback of the
perfect tangent class means derived pullback,
\[
 \nu^*T_Q^{\mathrm{vir}}
 =T_R-i_{0*}(J|_{D_0})-i_{\infty*}(J|_{D_\infty}).
\]
Indeed in a smooth local ambient space the normal line $A$ of
the branch and the normal line of the whole hypersurface
differ by $D$: the latter is $A(D)$. Their difference in
$K$-theory is $-[i_*(A(D)|_D)]$, and $A(D)|_D$ is exactly
the product of the two branch normals, namely $\mathcal T^1_Q$.
These local comparisons are intrinsic and agree on overlaps.
On the infinity branch $J$ means its $T^{-1}$-pullback
from the double locus, canonically the same line bundle
because it comes from $B$. Thus $J$ extends from both sections by the pullback from $B$,
so $[i_*J]=[J]-[J(-D)]$.

Put $j=c_1(J)=\pi_R^*F$, $d_i=[D_i]$. Since the sections are
disjoint, $d_0d_\infty=0$, and since $F$ is a fibre, $j^2=0$.
Consequently
\[
 c(\nu^*T_Q^{\mathrm{vir}})
 =c(T_R)\frac{(1+j-d_0)(1+j-d_\infty)}{(1+j)^2}
 =c(T_R)\{1-d_0-d_\infty+j(d_0+d_\infty)\}.
\]
The degree-three part has integral
\[
 e(R)-\sum_i\int_{D_i}c_2(T_R|_{D_i})
       +\sum_i\int_{D_i}j\,c_1(T_R|_{D_i}).
\]
We have $e(R)=2e(S)=24$. On a section with normal class
$n_i$, $c_2(T_R|_{D_i})=c_2(S)+c_1(S)n_i$ and
$c_1(T_R|_{D_i})=c_1(S)+n_i$. The two normal classes on
the unglued normalization are $m,-m$, where $m=c_1(L)$.
Their sum is zero; $c_1(S)=F$ and $F^2=0$. The displayed
number is therefore $24-24+0=0$. Because
$\nu_*[R]=[Q]$, this is the intrinsic virtual number on $Q$.
The projective-bundle calculation is independently checked
in \texttt{ab\_push\_local\_checks.py}, using
$H(H+m)=0$, $d_0=H+m$, $d_\infty=H$.

The untwisted quotient away from its cusp has Euler number zero:
scalar multiplication supplies a free circle action, including
the two additive places before logarithmic transformation. Indeed
a circle scalar identified with an iterate of $u\phi$ must be
the identity, since every positive iterate strictly contracts
the Hermitian norm and every negative iterate expands it.
Compact fibres with a free circle action have Euler number zero.
The chosen Mumford fibre has two zero-dimensional torus orbits
and all other torus strata have Euler number zero, so
$e(X(u))=2$. The same number follows independently from the
semistable nearby-cycle count in A.4.3 below. Lemma~\ref{push:A31}
therefore excludes such a smoothing of $Q$. No requirement
on its leading smoothing section, and no extension of $f$,
entered this argument.

There is also a direct check that zero is the correct smooth
number. Replace $\phi$ in the global construction of A.1.2
by a section of $J$ vanishing on a smooth elliptic fibre of
$S$. The compact collar construction still applies. At its
zero divisor the fixed nonzero-parameter equation is $xy=uz$
with $dz\ne0$, so every nearby threefold is smooth. The
double-locus nearby Euler function is one on that elliptic
curve and zero off it; its integral is zero. This is an actual
proper-flat smoothing of $Q$, but is not Engel's family.
\end{proof}

\begin{proposition}[A.3.3; Refuted and corrected]\label{push:A33}
The assertion in Lead 3, or A.3(iii), that the two natural
fillings exhaust smooth-total-space reduced-normal-crossings
fillings is false, even when joint flatness over
$B\times\Delta$ is imposed. There are infinitely many such
fillings of the same punctured untwisted family. The
two-model classification in A.2.1 is a classification of the
specified minimal toroidal refinements, not of all fillings.
\end{proposition}
\begin{proof}
We give centres that preserve joint flatness, rather than
blowing up a point over a point of the two-dimensional base.
The section $2P$ is disjoint from $P$ and $O$. In the
Weierstrass equation its coordinates are
\[
 x(2P)=\frac{9t^2}{8},\qquad
 y(2P)=\frac{\sqrt2\,t^2(5t-32)}{32}.
\]
At $0$ it meets the nonidentity component different from
that met by $P$; at $1$ the displayed coordinates are finite
and different from those of $P,O$; at infinity the minimal
coordinates are finite and different as well. This also
checks directly the zero intersection numbers
$P\cdot2P=O\cdot2P=0$. The rational section of
$\mathcal O(P-O)$ therefore restricts to a nowhere-vanishing
holomorphic section on $2P\simeq B$. Its graph is a compact
curve $C\subset L^\times\subset R^\pm$ disjoint from both
conductor sections and from the modification centres.
It gives a curve in the smooth locus of $Q^\pm$ mapping
isomorphically to $B$.

Blow up $C$ in the smooth total fourfold $\mathcal X^\pm$.
Locally the coordinates are $(t,u,a,b)$ with
$C=(u,a,b)=0$. The blow-up charts replace $u$ by $u$ itself
or by $aU$ or $bU$. Thus the new central divisor is reduced
normal crossings, and the total space is smooth. All fibres
over $(t,u)$ still have pure dimension two; the regular base
criterion proves joint flatness. The morphism is proper and
is an isomorphism for $u\ne0$. It adds a
$\mathbb P^2$-bundle over $C$, hence a new central irreducible
component. To justify repetition, $N_{C/Q^\pm}$ is an extension
of $\mathcal O_B(-1)$ by $\mathcal O_B$ and splits, and
$H^1(N_{C/Q^\pm})=0$. The normal sequence in the total family
therefore splits off the trivial parameter direction. Its
section in the exceptional $\mathbb P^2$-bundle is disjoint
from the intersection with the strict central transform and
has the same normal sequence. We may repeat indefinitely.
Their differing component counts preclude
isomorphism as fillings. The local substitutions and the
section-doubling calculation are exact computations, recorded
with the other local checks.
\end{proof}

\subsection{A.4. Integral topology and nearby cycles}

\begin{proposition}[A.4.1; Proved]\label{push:A41}
With integral coefficients, the spaces constructed above have
\[
\begin{array}{c|rrrrrrr|r}
&b_0&b_1&b_2&b_3&b_4&b_5&b_6&\chi\\\hline
R&1&0&11&0&11&0&1&24\\
R^\pm&1&0&12&0&12&0&1&26\\
Q&1&1&8&7&10&0&1&12\\
Q^\pm&1&1&9&7&11&0&1&14
\end{array}
\]
All these homology groups are torsion free. The fundamental
groups of $R,R^\pm$ are trivial and those of $Q,Q^\pm$ are
$\mathbb Z$. The dual complex of either $Q^\pm$ is a circle.
\end{proposition}
\begin{proof}
We check the lattice step of R10 on which the calculation depends.
Let $\mathcal L=H_2(S;\mathbb Z)$, with its unimodular intersection
pairing, and put
\[
 F=[S_t],\qquad m=P-O,\qquad v=[6P]-[O]-3F.
\]
The intersections are
\[
 F^2=Fm=Fv=0,\quad m^2=-2,\quad v^2=-6,\quad mv=-1,
 \quad OF=Om=1,\quad Ov=0.
\]
Translation is the integral isometry
\[
 T_*a=a+(aF)v-(av-3aF)F.
\]
This is checked on $F,O,v$ and on the seven root classes of
the two reducible Kodaira fibres, which translation fixes.
It also gives $T^*m=m-F$. The span of $F,v,m$ is primitive:
split off the unimodular plane generated by $F,O$; the Gram
determinant of $m-F,v$ in its orthogonal complement is eleven,
so an index larger than one would have its square divide eleven.
Unimodularity therefore makes
$a\mapsto(aF,av,am)$ surjective onto $\mathbb Z^3$.

The normalization mapping cylinder gives
\[
 H_i(S)\xrightarrow{d_i}H_i(R)\longrightarrow H_i(Q)
 \longrightarrow H_{i-1}(S),\qquad
 d_i=s_{0*}-s_{\infty*}T_*.
\]
The maps are induced by actual embeddings, and the mapping
cylinder is homotopy equivalent to the pinching because the
two section embeddings are disjoint cofibrations. With $f$ the
projective-line fibre, the degree-two map is
\[
 d_2(a)=-(aF)v+(av-3aF)F+(am-aF)f.
\]
Its three coefficient columns have determinant $-1$, so the
image is primitive of rank three, the kernel has rank seven,
and the cokernel has rank eight. In degree four the image is
the primitive class Poincar\'e dual to $m$, since $mO=1$.
In degree zero the difference map is zero. The projective-bundle
formula and this exact sequence give the asserted groups of
$R,Q$. Van Kampen on the same mapping cylinder gives the
single stable loop and no other generator, hence
$\pi_1(Q)=\mathbb Z$.

Now let $h:R^+\to R$ be the modification constructed in A.2.2.
Its fibres are points off the nodal curve $C=S_\infty$, one
$\mathbb P^1$ over $C_{\mathrm{reg}}$, and a chain of two
$\mathbb P^1$'s over the node. The chart calculation identifies
the sheaf $R^2h_*\mathbb Z$ with
$\eta_*\mathbb Z_{\mathbb P^1}$, where
$\eta:\mathbb P^1\to C$ is the normalization. Here is the
integral check at the node. If $(a,b)$ are the cohomological
degrees on the two exceptional curves, generization to the
two branches is $a$ and $a+b$: on one branch the generic
exceptional fibre specializes to the first curve, on the
other it specializes to their sum. The matrix
$\left(\begin{smallmatrix}1&0\\1&1\end{smallmatrix}\right)$
is unimodular. These generizations follow by substituting
$X=za$, $v=aB$ in the blow-up charts, or by taking the two
limits of their exceptional fibres. Proper local neighbourhoods
retract onto these fibres by the positive toric contraction,
so the fibre calculation computes the higher direct images.
There is no other positive-degree direct image.

The Leray spectral sequence has only the row $H^*(R,\mathbb Z)$
and the row $H^{*-2}(\mathbb P^1,\mathbb Z)$; both are free
and in even degrees. It degenerates integrally for degree
reasons, giving the stated groups of $R^+$. The same local
neighbourhood retractions and Van Kampen show that $h$ induces
an isomorphism on fundamental groups, so $R^+$ is simply
connected. The case $R^-$ interchanges the branches.

For $Q^+$ use the normalization mapping cylinder once again.
The new degree-two difference map is the old $d_2$ with
one extra coordinate $aF$: the exceptional divisor restricts
to $F$ on the blown-up section and to zero on the other.
It has the same primitive rank-three image and the same
kernel; its cokernel rank increases by one. The degree-four
difference remains primitive, since its pushforward to $R$
already was primitive, and its cokernel likewise increases
by one. This proves the groups of $Q^+$. Interchanging the
sections changes the extra coordinate by a sign and does
not change any group. Smith normal forms are recorded in
\texttt{ab\_push\_local\_checks.out}. The dual complex has
one vertex and one self-edge, with no higher cells, and
Van Kampen again gives $\mathbb Z$.

As independent Euler checks, stratification gives
$\chi(Q)=\chi(R)-\chi(S)=24-12$. In the modification $h$,
$C_{\mathrm{reg}}\simeq\mathbb C^*$ contributes zero and
the exceptional chain above the node increases Euler number
by $3-1=2$. Thus $\chi(R^\pm)=26$ and
$\chi(Q^\pm)=26-12=14$.
\end{proof}

\begin{proposition}[A.4.2; Proved]\label{push:A42}
The untwisted smooth fibre has torsion-free integral homology
with Betti numbers
\[
 (b_0,b_1,b_2,b_3,b_4,b_5,b_6)=(1,1,8,14,8,1,1).
\]
The positive-$u$ monodromy is the identity outside degree three.
In degree three there is an integral exact sequence
\[
 0\longrightarrow \mathcal L/\langle F,v,m\rangle
 \longrightarrow H_3(X(u);\mathbb Z)
 \longrightarrow K\longrightarrow0,\qquad
 K=\{a\in\mathcal L:aF=av=am=0\},
\]
in which $T_u-I$ is zero on the left term and, on the quotient,
is the map
\[
 K\longrightarrow \mathcal L/\langle F,v,m\rangle,\qquad
 a\longmapsto a\bmod\langle F,v,m\rangle.
\]
The sign convention orients the neck circle by increasing
$\arg x$ in $xy=u$; the boundary convention for a transported
chain is chosen so that a positive $u$ loop adds this circle
times its boundary cycle. With these conventions
$(T_u-I)^2=0$, its rank is six, and its nonzero Smith factors
are $1,1,1,1,1,11$.
\end{proposition}
\begin{proof}
We use the actual semistable model A.2.2. Remove disjoint
open tubular neighbourhoods of the two conductor sections
from $R^+$, and denote the remaining compact manifold by
$Z$. Its two boundary circle bundles are identified by the
plumbing; let $E$ denote the first. The Euler class of
$E\to S$ is $n=m-F$. It is primitive. The Thom sequence
for $(R^+,Z)$ gives
\[
 H_5(Z)=\mathbb Z,\quad H_4(Z)=0,\quad
 H_3(Z)=\mathcal L/\langle F,m\rangle\simeq\mathbb Z^8,
 \quad H_2(Z)\simeq\mathbb Z^{10},\quad H_1(Z)=0.
\]
We spell out the rank and saturation assertions. If $H$ is
the projective-bundle class restricting to zero on $s_0$
and to $-m$ on $s_\infty$, and $A$ is the exceptional
divisor class, the two restriction maps on $H^2(R^+)$ are
\[
 (\alpha,h,a)\longmapsto\alpha+aF,\qquad
 (\alpha,h,a)\longmapsto\alpha-hm.
\]
Their combined image is primitive of rank twelve, and its
cokernel is $\mathcal L/\langle F,m\rangle$. In degree two,
intersection with the two sections is onto $\mathbb Z^2$:
a generic projective fibre meets both once, and an exceptional
projective fibre meets just the section containing the centre
once. In degree six the map is the diagonal
$\mathbb Z\to\mathbb Z^2$. These are precisely the Thom
maps giving the asserted groups.

The Gysin sequence of the primitive circle bundle gives
\[
 H_2(E)=n^\perp\subset\mathcal L,\qquad
 H_3(E)=\mathcal L/\mathbb Z n,\qquad
 H_1(E)=H_4(E)=0,\quad H_5(E)=\mathbb Z.
\]
The smooth fibre is obtained from $Z$ by the boundary
mapping cylinder. In its Mayer--Vietoris sequence the
degree-two difference map is the normalization difference
$d_2^+$ restricted to $n^\perp$. In the coefficient
coordinates $(s,t,h)=(aF,av,am)$ the constraint is $h=s$,
so this map has primitive rank two and kernel $K$.
The degree-three difference map is, modulo $F,m$,
$(T_*-I)a=(aF)v$, hence has primitive image $\mathbb Z v$.
These statements use the surjectivity of the three
intersection coordinates proved in A.4.1, so no integral
index is omitted. The degree-five difference is zero:
the plumbing reverses the normal circle orientation, and
the two boundary fundamental classes sum to zero in $Z$.
The degree-zero difference is zero as well. Exactness now
gives the claimed groups and the displayed degree-three
extension. Poincar\'e duality, or the same integral exact
sequence, confirms the absence of torsion in complementary
degrees.

We also compute monodromy directly from this decomposition.
In the neck $xy=u$, a positive circuit of $u$ rotates one
boundary normal coordinate once relative to the other.
A cycle in $Z$ is unchanged. A three-cycle whose cut-open
boundary represents $a\in K$ acquires the tube obtained
by sweeping $a$ once around that normal circle. The Gysin
description sends this tube to the class of $a$ in
$H_3(E)=\mathcal L/\mathbb Z n$, and then to its class in
$\mathcal L/\langle F,m,v\rangle$. This is exactly the
map in the statement. The same description in the other
degrees is zero, using the groups and difference maps just
computed. A tube has no cut-open boundary, which proves
$(T_u-I)^2=0$ without an appeal to a monodromy theorem.

Finally $K\cap\langle F,m,v\rangle=\mathbb ZF$, since
the Gram matrix of $m-F,v$ has determinant eleven.
Hence the rank is six. In the unimodular $E_8$ summand,
the primitive rank-two lattice spanned by $m-F,v$ has
discriminant eleven; its orthogonal complement maps into
the dual quotient with cokernel of order eleven. Thus the
nonzero Smith factors are five ones and eleven. The fully
integral check uses an explicit $E_8$ basis in
\texttt{ab\_push\_monodromy.py}; its output records the
kernel basis, quotient projection, and all entries of the
seven-by-seven monodromy block. This supplies a second,
exact lattice verification of the discriminant argument.
\end{proof}

\begin{proposition}[A.4.3; Proved; pointwise correction to Lead 1]\label{push:A43}
For $\mathcal X^\flat\to\Delta_u$ the nearby-fibre Euler
function on the double locus of $Q$ is zero off $S_\infty$
and one at every point of $S_\infty$, including its node.
Its integral is one. In particular the assertion that the
pointwise value on the smooth part of $S_\infty$ is zero
is false; its integrated contribution is zero.
For either semistable family $\mathcal X^\pm$ the nearby
Euler function is one off the double locus and zero on
it, giving $\chi(X(u))=26-2\cdot12=2$.
\end{proposition}
\begin{proof}
Off the zero divisor of $\phi$, the transverse nearby
fibre of $xy=u\phi$ is an annulus and has Euler number
zero. At a smooth zero, the local equation is $xy=uz$;
for fixed $u\ne0$ one solves for $z$, and the small
local nearby space is contractible. Its Euler number
is one. At the node it is the quadratic cone
$xy=uzv$, also contractible, with Euler number one.
The smooth part of the nodal elliptic curve is
$\mathbb C^*$ and has Euler number zero; the node
contributes one. The core $Q\setminus S$ is a
$\mathbb C^*$-bundle and contributes zero.
This proves the first count.

In the resolved family the local map on the double locus
is simply $ab=u$ times two transverse coordinates,
including above the old node. Its nearby annulus has
Euler number zero everywhere on that locus. Away from
it, the family is a submersion and the local value is
one. Integrating gives the second count. If the nodal
threefold in the first family is smoothed rather than
small-resolved, its ordinary-double-point Milnor fibre
is homotopy equivalent to $S^3$ and has Euler number
zero; the replacement changes $1$ to $0$, whereas
the small resolution changes $1$ to $2$. These agree
with the intrinsic virtual number in A.3.2.
\end{proof}

\begin{proposition}[A.4.4; Proved; limitation of Clemens--Schmid]\label{push:A44}
The classical K\"ahler or projective Clemens--Schmid
comparison cannot be applied to these untwisted models.
Their smooth fibres have $b_1=1$ and are not K\"ahler.
Moreover the exceptional small-resolution curve in
$R^\pm$ is null-homologous, so that normalization is
not K\"ahler either. The rational monodromy weight
filtration, defined algebraically from $N=T_u-I$ and
centred at degree three, has graded ranks $6,2,6$ in
weights $2,3,4$. This is a filtration of the computed
monodromy representation; we do not identify it with
a polarized limiting mixed Hodge structure.
\end{proposition}
\begin{proof}
The odd first Betti number excludes a K\"ahler metric.
In the generization matrix of A.4.1 a global section of
$R^2h_*\mathbb Z$ has equal branch values, hence degree
zero on the small exceptional curve. All classes pulled
back from $R$ also have degree zero on it. The Leray
calculation exhausts $H^2(R^\pm)$, so the curve has
zero homology class; it cannot have positive area for
a K\"ahler form. Finally $N^2=0$, rank $N=6$, and
$\dim H_3=14$, so the centred filtration is
$W_2=\operatorname{im}N$, $W_3=\ker N$, $W_4=H_3$.
Its graded dimensions are the asserted ones. These
topological computations replace, rather than assume,
the inapplicable K\"ahler comparison.
\end{proof}

\begin{corollary}[A.4.5; Proved; unrestricted toroidal topology]\label{push:A45}
For every output of A.2.4, A.4.7--A.4.9 determine the
Euler number, fundamental group, all integral homology
groups with torsion, nearby Euler function and marked
monodromy. The integral normalization computation is
given by the finite algorithm and the primitive
central pushforward matrices of A.4.9; the independent
ambient computation is A.4.8, and the torus-stratum
Euler count is A.4.7. The non-K\"ahler limitation on
classical Clemens--Schmid is the explicit one in
A.4.4; no polarized limiting Hodge structure is
asserted.
\end{corollary}

\begin{proposition}[A.4.6; Proved; an infinite explicit subclass]\label{push:A46}
Let $D_m$ be the central fibre after the $m$ maximal-cone
star subdivisions in A.2.1, starting with $Q^+$, where
$m\geq0$. Then
\[
 \begin{gathered}
 \pi_1(D_m)=\mathbb Z,\qquad
 H_*(D_m;\mathbb Z)\text{ is free},\\
 (b_0,\ldots,b_6)=(1,1,9+m,7,11+m,0,1+m).
 \end{gathered}
\]
and $\chi(D_m)=14+3m$. For $m\geq1$ its normalization is
the disjoint union of $\operatorname{Bl}_pR^+$, $m-1$
copies of $\operatorname{Bl}_{p_i}\mathbb P^3$, and a final
$\mathbb P^3$. The original two sections are still glued
by $T$. The new components form a chain attached to the
original component. Consecutive components meet in one
$\mathbb P^2$, with no triple intersections. The dual
complex is the original circle with a chain of $m$ edges
attached. The nearby Euler function is one off the double
locus and zero on it; its integral remains two. The nearby
integral monodromy remains exactly that of A.4.2.
\end{proposition}
\begin{proof}
The torus orbit of the cone $ABCE$ is a point on a single
local central branch, because precisely one of its four
rays has positive $u$-height. Hence it is a smooth point
of the central fibre, away from the conductor. A star at
the sum of the four basis rays is the ordinary blow-up
of that point in the smooth fourfold. The central divisor
is the union of its strict transform and the exceptional
$\mathbb P^3$, both with multiplicity one. Their
intersection is the exceptional $\mathbb P^2$ in the
blow-up of the central threefold, identified with a
hyperplane in $\mathbb P^3$. Their two normal lines there
are $\mathcal O(-1)$ and $\mathcal O(1)$.

The next distinguished cone $ABCv_1$ has only the new
central ray $v_1$. Its point lies on the new $\mathbb P^3$
away from its attaching hyperplane. Repeating therefore
blows up that component at a point away from the preceding
double locus and attaches the next $\mathbb P^3$. The
incoming hyperplane and outgoing exceptional plane are
disjoint in each middle $\operatorname{Bl}_{p_i}\mathbb P^3$.
This proves the complete normalization and gluing
description, the absence of triples and the dual complex.

For the integral calculation, blowing up a smooth point
of a compact complex threefold adds one free generator
in each of $H_2$ and $H_4$ and changes no other homology.
This also holds when the threefold has singularities away
from that point: remove a small ball and use the
Mayer--Vietoris sequence with boundary $S^5$; the
replacement retracts to $\mathbb P^2$. Now attach
$\mathbb P^3$ along that $\mathbb P^2$. In degrees
$0,2,4$ the inclusion into $\mathbb P^3$ is an integral
isomorphism, so the Mayer--Vietoris boundary map is
injective and primitive. It contributes no torsion or
odd homology. The new component contributes one free
$H_6$ generator, since $H_6(\mathbb P^2)=0$. Each step
therefore adds precisely one to $b_2,b_4,b_6$.
Van Kampen gives no change of fundamental group: the
blow-up ball replacement, the attaching plane and the
new projective component are simply connected.
Induction from A.4.1 proves the asserted groups.

Independently, the strict-transform blow-up adds two to
the central Euler number, the new $\mathbb P^3$ adds
four, and its intersection $\mathbb P^2$ subtracts
three. Thus each step adds three, agreeing with the
homology computation. The double locus is the disjoint
union of the original $S$ and $m$ planes, of Euler number
$12+3m$. Since the local equation is $u=x$ off this locus
and $u=xy$ on it, the nearby Euler integral is
$(14+3m)-(12+3m)=2$. All modifications are the identity
over $u\ne0$, so they do not alter the marked punctured
family or its homological monodromy. The exact saved
certificates check all determinants, boundary faces,
slice volumes and these integer formulas for $0\leq m\leq8$;
the displayed induction proves them for every $m$.
\end{proof}

\begin{theorem}[A.4.7; Proved; Euler number and fundamental group for every toroidal output]\label{push:A47}
Let $\Sigma$ be any output of A.2.4 for $u=r^n$, and let
$D_\Sigma$ denote the underlying reduced central space.
Let $k(\Sigma)$ be the number of selected pure top ray
orbits, and let $t_4(\Sigma)$ be the number of maximal
four-cones in one fundamental prism. Then
\[
 \chi(D_\Sigma)=11k(\Sigma)+t_4(\Sigma),
 \qquad \pi_1(D_\Sigma)=\mathbb Z,
 \qquad H_1(D_\Sigma;\mathbb Z)=\mathbb Z.
\]
These assertions include nonreduced central divisors.
At a point with precisely one local central branch of
multiplicity $\ell$, the nearby Euler function is $\ell$;
at a point with at least two local branches it is zero.
Its integral is always two. For every reduced output the
pure top count is $k(\Sigma)=n$, so
$\chi(D_\Sigma)=11n+t_4(\Sigma)$.
The higher integral homology of every arbitrary output is
not asserted by this theorem.
\end{theorem}
\begin{proof}
Over $B\setminus\{\infty\}$ the top fan gives a cycle
of $k(\Sigma)$ ruled bundles over
$S\setminus S_\infty$, glued along their zero and
infinity sections with the specified translations.
There are no triple intersections there. Its Euler
number is therefore
\[
 2k(\Sigma)\chi(S\setminus S_\infty)
       -k(\Sigma)\chi(S\setminus S_\infty)
 =11k(\Sigma),
\]
because $\chi(S)=12$ and $\chi(S_\infty)=1$.
This computation concerns the support of the divisor,
so the multiplicities do not enter.

At $q=r=0$ the strata are the torus orbits of cones whose
generators collectively have positive $q$-height and
positive $r$-height. Each such cone has trivial period
stabilizer. Indeed, an element $G_1^aG_2^b$ preserving
its finite set of primitive rays has a power fixing them
all. A ray of positive $q$-height then forces $a=0$;
any ray with positive $q$- or $r$-height forces $b=0$.
Consequently a stratum of cone dimension $d$ is an actual
$(\mathbb C^*)^{4-d}$, with no additional torus quotient.
The induced gluing units only identify different period
translates of the stratum. Its compactly supported Euler
number is zero unless $d=4$, in which case it is one.
Every maximal cone has both kinds of height, and each has
one representative in the interior of the fundamental
prism. Thus the cusp contributes exactly $t_4(\Sigma)$.
Additivity proves the displayed central Euler formula.

The local total space is smooth, and its parameter is
$r=z_1^{\ell_1}\cdots z_s^{\ell_s}$ times a holomorphic
unit, with $\ell_i>0$. Its Milnor fibre has the homotopy
type of a finite disjoint union of $(S^1)^{s-1}$.
For $s=1$ it consists of $\ell_1$ contractible pieces;
for $s\geq2$ its Euler number is zero. This proves the
pointwise assertion. Off the cusp the one-branch parts
are $\mathbb C^*$-bundles over the appropriate parts of
$S\setminus S_\infty$ and have Euler number zero; the
double parts have nearby Euler value zero. At the cusp,
positive-dimensional torus strata again contribute zero.
A maximal cone with exactly one positive-$r$ ray must
have three bottom rays. Its bottom face is consequently
one of the two prescribed triangles $ABC$ and $BCD$.
Each triangle belongs to exactly one such maximal cone.
Its determinant is the $r$-height $\ell$ of the fourth
ray, because the bottom triangle is unimodular.
Smoothness forces $\ell=1$. Exactly these two points
contribute one each. Thus the nearby Euler integral is
two for every output, even when some other components
have multiplicity greater than one.

For the fundamental group, first choose the reference
model obtained by base-changing $\mathcal X^+$ by
$u=r^n$ and resolving the surface equation $xy=r^n$
along its smooth conductor $S$. Its normalized central
components are $R^+$ and $n-1$ projective-line bundles
over $S$. All are simply connected, and their double
surfaces are copies of the simply connected $S$.
The gluing graph is an $n$-cycle, interpreted as the
self-gluing circle when $n=1$. Van Kampen therefore
gives fundamental group $\mathbb Z$ for this reference
centre.

Any two fans of A.2.4 have the explicit proper common
normal modification formed by intersecting their cones.
Resolve this analytic space to a smooth common model;
analytic resolution, including the noncompact version
near the compact central set, is supplied by
\cite[Theorem 1.6 and its following paragraph]{BM95AB}.
This use of a common modification is legitimate here
because the common fan has already constructed it.
It is not an assumed extension of an arbitrary punctured
isomorphism.

A proper bimeromorphic map $F:Y\to X$ between smooth
complex manifolds induces an isomorphism of fundamental
groups. To see this directly, choose an analytic
$A\subset X$ of complex codimension at least two such
that $F$ is an isomorphism outside $A$. General position
gives $\pi_1(X\setminus A)\simeq\pi_1(X)$, and gives
surjectivity of
$\pi_1(Y\setminus F^{-1}A)\to\pi_1(Y)$.
Under the common-open identification, composition with
$F_*$ is the displayed isomorphism to $\pi_1(X)$.
Thus the first map is also injective, and $F_*$ is an
isomorphism. Apply this argument to both maps from the
smooth common model, after restricting to a sufficiently
small parameter disc.

Finally, each smooth total space retracts to the support
of its central divisor, including in the nonreduced case.
The gradient proof in B.4.2 applies with the explicit
bound for $\rho=|r|^2$
\[
 \|\nabla\rho\|\geq c\rho^{1-1/(2L)},
 \qquad L=\ell_1+\cdots+\ell_s.
\]
On the complement of the divisor it follows from the
weighted arithmetic--geometric mean inequality applied
to $\sum_i\ell_i^2/|z_i|^2$. A holomorphic unit can be
absorbed into one coordinate locally. Compactness of
the central set gives a uniform exponent below one,
and the resulting finite remaining path-length bound
gives the continuous gradient retraction. Hence all
the centres and their smooth total neighbourhoods have
the same fundamental group as the reference centre.
Abelianization gives the stated $H_1$.

The program \path{code/ab_push_all_fan_topology.py}
checks the Euler and two-point nearby formulas for every
minimal bridge at $n=1,2,12$, eight point-star
modifications, and five nonreduced conductor-star
modifications. In the latter family the added pure top
rays are $(0,1,0,j)$ for $2\leq j\leq m+1$; the formula
gives $\chi(D_m)=14+12m$, independently agreeing with
successive blow-ups of the conductor surface of Euler
number twelve. The finite checks are recorded separately
from the proof for every fan.
\end{proof}

\begin{theorem}[A.4.8; Proved; integral homology of every toroidal output]\label{push:A48}
Let $D_\Sigma$ be the reduced central space of an output of
A.2.4 for $u=r^n$. Write $k$ for its number of pure top
ray orbits, $m$ for its number of mixed ray orbits, and
$t_4$ for its number of maximal four-cone orbits. Then
every integral homology group is free, and
\[
 \begin{split}
 (b_0,b_1,b_2,b_3,b_4,b_5,b_6)
   ={}&(1,1,8+k+m,7,\\
       &\hspace{20mm}9k+t_4-2m-1,0,k+m).
 \end{split}
\]
This includes nonreduced central divisors: the homology
is that of their underlying support. The marked nearby
fibre is $X(r^n)$, and its integral monodromy is
$I+nN$ in degree three, where $N$ is the map in A.4.2;
it is the identity in every other degree. In particular
the six nonzero Smith factors of this degree-three
monodromy minus the identity are
$n,n,n,n,n,11n$.
\end{theorem}
\begin{proof}
We first compute a reference centre by its normalization
mapping cylinder. Base-change $\mathcal X^+$ by $u=r^n$
and resolve $xy=r^n$ along its smooth conductor $S$.
Put $N_0=\mathcal O_S(m-F)$, using the lattice notation
of A.4.1. The normalization consists of $R^+$ and
$n-1$ copies of $\mathbb P_S(\mathcal O\oplus N_0)$,
in a cycle. The latter assertion follows directly on
the surface-resolution charts: if $x$ has transition
function $g$, the coordinates $x/r^i$ have the same
transition function, while their complementary
coordinates have transition function $g^{-1}$.
Only the closing identification is $T$.

In the integral mapping-cylinder sequence the degree-two
map has domain $\mathcal L^n$ and target of rank $11n+1$.
Project each intermediate ruled component to $S$.
The resulting $n-1$ difference blocks are integral
identity matrices. Eliminate them by integral row and
column operations. The remaining map from $\mathcal L$
is $d_2^+$ from A.4.1, together with $n-1$ coordinates
$a\mapsto a(m-F)$. Those coordinates are integral
combinations of the three primitive coordinates of
$d_2^+$. Thus the original degree-two map has primitive
image of rank $10(n-1)+3$, kernel $\mathbb Z^7$, and
cokernel $\mathbb Z^{n+8}$.

In degree four, the same elimination uses the $n-1$
identity blocks on the fundamental classes of the bases.
The remaining map from $\mathbb Z$ contains $d_4^+$,
which is primitive by A.4.1; the other entries are the
section differences $m-F$. Hence the degree-four map
is injective with primitive image of rank $n$, and
cokernel of rank $10n+1$. In degree zero it is the
incidence map of an $n$-cycle. In degree six the
normalization contributes one independent fundamental
class per component. Consequently the reference centre
has free integral homology and Betti tuple
\[
 (1,1,n+8,7,10n+1,0,n).
\]
This is also a direct integral calculation of all the
maps needed below; it uses no degeneration theorem.

We next justify the precise factorization used to compare
an arbitrary output with this reference. On each total
space take the toroidal boundary
\[
 E=\bigl(\{r=0\}\cup f^{-1}(\infty)\bigr)_{\rm red}.
\]
In the cusp charts it is the toric boundary, and outside
the cusp its nontrivial toroidal charts are the plumbing
charts along $S$. The unit transition functions preserve
monomial ideals. A.2.4 gives a finite cone complex modulo
the two periods, with the possible self-identifications
of its faces explicitly retained.

Make this complex strict by barycentric star subdivisions,
in decreasing cone dimension. At the moment a stratum
is blown up, all of its deeper self-intersections have
already been separated. Its closure is therefore smooth,
and different translates in the chosen orbit are disjoint.
The star at the sum of the primitive basis generators is
the ordinary blow-up of that closure. Repeating over the
finitely many cone orbits gives smooth strict toroidal
embeddings. This is the geometric barycentric subdivision
used in \cite[Section 5.6, pp.565--566]{AKMW02AB}; its
analytic realization by the same monomial ideals is
described in Section 5.5 of that paper. Pure $q$-strata
are allowed in this preparatory step. Thus we do not
assume that strictification preserves the punctured
$r$-family.

Intersect the two subdivided fans, and take a regular
subdivision, to obtain their explicit common strict
toroidal modification. Apply toroidal weak factorization
to each of its two subdivision maps. The needed theorem
is the finite polyhedral-complex factorization by
nonsingular star subdivisions and inverse subdivisions,
\cite[Theorems 1.6.1 and 2.7.1]{AKMW02AB}. The corrected
form for a morphism states that its centres are smooth
irreducible closed strata and that all intermediate
complexes map to the target
\cite[Section 2, p.2]{Matsuki00ErrAB}.

There is no invocation here of a general relative
analytic weak-factorization theorem. The complexes in
question are finite. The cited combinatorial
factorization applies to them; the analytic
subdivision--modification correspondence of
\cite[Section 5.5]{AKMW02AB} realizes each of its finitely
many stars by the indicated monomial blow-up.
The correspondence respects the face identifications
and the bundle units, so the local analytic blow-ups
glue. All intermediate morphisms are proper over the
target of the subdivision map, hence over a sufficiently
small $r$-disc. They remain smooth, and both parameters
remain monomials times units. This gives the required
finite chain of actual analytic blow-ups and blow-downs
between the total neighbourhoods.

We classify the connected centres in this chain. A
nontrivial centre has codimension at least two. For a
cone having a positive $q$-height ray, its period
stabilizer is trivial: a power of a stabilizing element
fixes each of its finitely many primitive rays; the
positive $q$-height first forces the exponent of $G_1$
to vanish and then the exponent of $G_2$ to vanish.
The same argument holds for pure bottom cones. Thus
there is no elliptic quotient hidden in those centres.
The normalization of the closure of their torus stratum
is the toric variety of its star. If the closed stratum
is smooth, this finite birational normalization map is
an isomorphism.

For a mixed cone the centre is contained in $q=r=0$;
its star is complete, and it is a compact smooth toric
point, curve, or surface. It has free integral
cohomology in even degrees only. For a pure top centre,
the cone must have dimension two. Its stabilizer is
precisely $\langle G_1\rangle$, and the chart quotient
and the plumbing identification show that its dense
open part is the corresponding open part of $S$.
The compact smooth centre is therefore a surface
birational to the rational surface $S$. Equivalently,
its proper map to the unmodified top stratum is a
birational map of smooth surfaces. It is a smooth
rational surface and again has free even cohomology.
It need not remain isomorphic to $S$ after earlier
mixed modifications.

For a pure bottom centre the cone has dimension two
or three. Its star gives a smooth toric surface or
curve proper over the $r$-disc. The induced character
$r$ is primitive on its quotient lattice, since it
was primitive before quotienting by a smooth cone
on which it vanishes. Its generic fibre is therefore
connected, and is respectively a complete toric
$\mathbb P^1$ or a point. In the surface case the
ordered two-dimensional fan gives a chain of rational
curves as the central support, possibly with positive
multiplicities. In the curve case the support is one
point. Monomial gradient retraction, as in A.4.7,
shows that the ordinary integral cohomology of the
centre is free and even, with $H^0=\mathbb Z$ and
$H^4=0$. These centres are not compact; this is why
they must be distinguished from the first two types.

For any closed smooth complex centre $C\subset Y$ of
complex codimension $c$, compact or not, the integral
blow-up splitting is
\[
 H^j(\operatorname{Bl}_C Y;\mathbb Z)
 \simeq H^j(Y;\mathbb Z)\oplus
       \bigoplus_{a=1}^{c-1}H^{j-2a}(C;\mathbb Z).
\]
For clarity, this statement requires no K\"ahler
hypothesis. Tubular excision identifies the relative
groups by the integral Thom class; the projective-bundle
theorem for the exceptional divisor supplies the powers
of its hyperplane class. Their exceptional Gysin maps,
together with pullback from $Y$, give the displayed
map. The relation that restriction of the exceptional
class is minus the hyperplane class makes the relative
matrix triangular with unit diagonal. Proper
blow-down pushforward, defined by Borel--Moore duality,
splits the pullback summand. This proves the splitting
also for the noncompact centres just listed.

It follows that all odd cohomology groups and all
torsion are unchanged along the factorization. Each
connected centre adds one to $b_2$ and one to the number
of prime boundary components. Each compact centre
adds one to $b_6$ and one to the number of central
prime components: in degree six the only new group
is its top cohomology. A pure bottom centre adds
neither, because its $H^4$ (and, for a curve, its
$H^2$) vanishes. The reverse operations subtract
the same contributions.

Consequently $b_2$ minus the number of all boundary
ray orbits, and $b_6$ minus the number of central ray
orbits, are invariants of this comparison. Both
endpoints have exactly one pure bottom ray orbit.
The reference has $n$ central components and
$b_2=n+8$, so the arbitrary output has
\[
 b_2=8+k+m,\qquad b_6=k+m,
 \qquad (b_1,b_3,b_5)=(1,7,0).
\]
All groups are free by the integral splitting in both
directions. The total neighbourhoods retract onto
their central supports by A.4.7, and those compact
analytic supports have finite triangulations. The
universal coefficient theorem therefore transfers
the same freeness and ranks to integral homology.
Finally the independently computed torus-stratum
Euler number $11k+t_4$ of A.4.7 gives
$b_4=9k+t_4-2m-1$.

Every original output, unlike temporary intermediate
factorizations, agrees with the marked family
$X(r^n)$ off $r=0$. A positive $r$-loop makes $n$
positive $u$-loops. Since $N^2=0$, its monodromy is
$(I+N)^n=I+nN$. Multiplication of the already computed
integral Smith matrix by $n$ gives the stated factors.

The exact program \path{code/ab_push_all_fan_homology.py}
checks the full reference normalization maps and their
Smith forms at $n=1,2,3,4,6,12$, the blow-up rank
increments for every centre type used here, and the
resulting formula on the certified fan families of
A.4.7. These finite checks supplement the integral
row-reduction and factorization argument for every
fan; they are not a bounded search substituted for it.
\end{proof}

\begin{proposition}[A.4.9; Proved; a second integral central-space computation]\label{push:A49}
The homology formula of A.4.8 can also be computed
on the central spaces, by their normalization
resolutions and the strata of each recorded star.
More precisely, let $Y'\to Y$ be one ordinary
blow-up in the algorithm of A.2.5, with closed
smooth centre $C$ of complex codimension $c$.
Let $D,D'$ be the reduced central supports and
$C_0=C\cap D$. In the derived category of integral
sheaves on $D$ there is an isomorphism
\[
 R(\pi_D)_*\mathbb Z_{D'}\simeq
 \mathbb Z_D\oplus
 \bigoplus_{a=1}^{c-1}(i_0)_*\mathbb Z_{C_0}[-2a],
 \tag{A.49a}
\]
where $\pi_D:D'\to D$ and $i_0:C_0\hookrightarrow D$
are the induced continuous maps. In particular
\[
 H^j(D';\mathbb Z)\simeq H^j(D;\mathbb Z)
 \oplus\bigoplus_{a=1}^{c-1}H^{j-2a}(C_0;\mathbb Z).
 \tag{A.49b}
\]
The maps in this splitting are integral, with unit
diagonal on the exceptional-fibre basis. Thus it
also determines the primitive boundary maps and
the absence of torsion in the normalization
mapping-cylinder computation along the entire
finite algorithm, not just the Euler characteristic.
\end{proposition}
\begin{proof}
We construct the sheaf maps rather than infer a
splitting from fibre ranks. Let $E=\mathbb P_C(N_{C/Y})$
be the exceptional divisor, let $p:E\to C$, and put
$\xi=c_1(\mathcal O_E(1))$. Pullback from $Y$ gives
the degree-zero summand. For $1\leq a\leq c-1$,
pullback by $p$, multiplication by $\xi^{a-1}$,
and the exceptional Gysin morphism give a derived
map
\[
 i_*\mathbb Z_C[-2a]\longrightarrow
 R\pi_*\mathbb Z_{Y'},\qquad i:C\hookrightarrow Y.
\]
The Gysin map exists integrally for a closed complex
divisor, by its oriented tubular Thom class; neither
the centre nor the ambient space must be compact.
At a point of $C$, proper base change identifies
the stalk of the target with
$H^*(\mathbb P^{c-1};\mathbb Z)$.
Restriction of the exceptional normal line is
$\mathcal O_E(-1)$, so the displayed maps send
the generators to
\[
 1,-\xi,-\xi^2,\ldots,-\xi^{c-1}.
\]
This is an integral basis. Off $C$ the map is
the identity. It is therefore a stalkwise
quasi-isomorphism
\[
 \mathbb Z_Y\oplus
 \bigoplus_{a=1}^{c-1}i_*\mathbb Z_C[-2a]
 \xrightarrow{\ \simeq\ }R\pi_*\mathbb Z_{Y'}.
\]

Now restrict this actual sheaf morphism to $D$.
The central support upstairs is the whole inverse
image $\pi^{-1}D$, including when the centre is
horizontal. Proper base change applies to this
Cartesian square. The inverse-image functor for
integral sheaves is exact, and the inverse image
of $i_*\mathbb Z_C$ is
$(i_0)_*\mathbb Z_{C_0}$. This proves (A.49a).
Taking derived global sections proves (A.49b).
This calculation concerns the stratified map of
central spaces. It does not first calculate the
cohomology of either ambient total neighbourhood.

Here is its explicit relation to the requested
normalization mapping cylinders. Resolve
$\mathbb Z_{D'}$ by the alternating ordered-local-branch
complex: its first term is the constant sheaf on
the normalization, its next term is on the normalized
double-branch space, and so on. Its augmentation is
an integral quasi-isomorphism, because at a point
with $s$ branches it is the augmented simplex
complex on those branches. No division by $s!$ is
used; one uses alternating integral cochains, or
equivalently the orientation local system of each
unordered branch set. Proper pushforward of this
resolution computes the left side of (A.49a).
The same resolution for $D$, together with the
centre complex, computes its right side.

The local matrices can be read directly from the
blow-up charts in A.2.5. If the centre lies on $k$
central branches among its $c$ defining coordinates,
the exceptional fibre is $\mathbb P^{c-1}$. The new
exceptional branch supplies that whole projective
space; the strict transform of old branch $i$
supplies its coordinate hyperplane $H_i$.
Intersections supply the corresponding coordinate
projective subspaces. In each even cohomological
degree the normalization differential is the
alternating restriction matrix with entries
$0,1,-1$. The cover consisting of the whole
projective space and its selected hyperplanes
has the elementary simplex contraction, leaving
one generator $\xi^a$ in each degree $2a$,
$0\leq a<c$. For a horizontal centre, every
central branch meeting its generic central
stratum supplies the whole exceptional fibre;
the same simplex contraction applies to these
copies. Higher boundary strata are treated by
the same coordinate restrictions. The global
Gysin maps just constructed glue the surviving
classes with unit coefficient, so this local
cancellation cannot hide an extension or an
integral index.

Starting from the explicitly computed reference
normalization matrices in A.4.8, (A.49a) therefore
computes the normalization complex along every
recorded blow-up or inverse blow-up. A compact
centre has $C_0=C$. A horizontal toric surface
centre retracts to its rational-chain $C_0$;
a horizontal curve centre has $C_0$ a point.
Their cohomology is free and even. Consequently
the central computation preserves odd cohomology
and torsion, adds one to $b_2$ for every new
boundary ray, and adds one to $b_6$ exactly for
a new central ray. It gives the same
$(b_1,b_3,b_5)=(1,7,0)$,
$b_2=8+k+m$ and $b_6=k+m$ as A.4.8.

Independently, stratification of $\pi_D$ into
$D\setminus C_0$ and $C_0$, with fibres a point
and $\mathbb P^{c-1}$, gives
\[
 \chi(D')-\chi(D)=(c-1)\chi(C_0).
\]
The absolute torus-stratum count of A.4.7 is
$\chi(D)=11k+t_4$. Together with the central
group calculation this gives again
$b_4=9k+t_4-2m-1$. Thus the two computations
agree integrally and in Euler number.

The program \path{code/ab_push_central_pushforward.py}
constructs every local alternating projective-subspace
matrix for $2\leq c\leq4$ and $1\leq k\leq c$,
as well as the possible horizontal cases. Its
Smith calculations check the unit cancellations
in every local dimension occurring in the
fourfolds. The sheaf construction proves their
global compatibility for arbitrary input fans.
\end{proof}


\section{B. Local limits and the global component obstruction}
\subsection{B.1. The finite root neighbourhoods}

\begin{proposition}[B.1.1; Refuted and corrected]\label{push:B11}
The product assertion in Lead 4 is false on the whole root
disc. There is, however, a proper flat holomorphic filling of
the marked root family over
$\Delta_{s_j}\times\Delta_r$, where $u=r^{12}$, with smooth
total space and reduced normal-crossings $r$-central divisor.
Its central surface over an ordinary value of $s_j$ is a cycle
of twelve ruled surfaces over the smooth elliptic curve
$E'_{j,s_j}$. Their normalizations are projective bundles
associated with the degree-zero line bundle represented by
$P_{s_j}-O_{s_j}$, with the appropriate pullbacks by translations.
They are not generically products with $\mathbb P^1$.
\end{proposition}
\begin{proof}
On the root disc the elliptic family $E'_j\to\Delta_{s_j}$ is
smooth. In the additive period coordinates $(\xi,\zeta_2)$,
put $W=\exp(-2\pi i\zeta_2)$. Quotienting by
$\ker\gamma\subset\Lambda$ gives a holomorphic
$\mathbb C^*$-bundle over this elliptic family. The period
$\widehat\gamma$ acts by
\[
 (\xi,W)\longmapsto
 (\xi+6\mu(s),a(s)uW),\qquad
 a(s)=\exp(-2\pi i\beta^{\rm part}(s)).
\]
This is the bundle and translation in R8, with inverse fibre
coordinate; the period identities check its transition
functions directly. The factor $a$ is a holomorphic unit on
the entire root disc. Apply the compact-core plumbing of
A.1.2 relatively over a closed smaller root disc. Here the
smoothing section never vanishes, so the only local equation
at the conductor is $xy=u$ after a change of frame. This
gives a smooth total space before the base change.

After $u=r^{12}$, the local equation is $xy=r^{12}$,
times the two smooth coordinates on $E'_j$. Resolve its
$A_{11}$ surface singularity by its minimal toric resolution.
In lattice coordinates this inserts each primitive height-one
ray between the two boundary rays, at unit intervals. The
construction is canonical for changes $x\mapsto a x$,
$y\mapsto a^{-1}y$ by a holomorphic unit, so these local
resolutions glue. Its affine charts have the equation
$r=XY$, or $r=X$ at a smooth point of the central divisor,
times the elliptic and root coordinates. Thus the total space
is smooth, the divisor is reduced with normal crossings, and
the joint map to $(s_j,r)$ is flat. The resolution is proper,
and the plumbing was proper, so the joint map is proper.
The comparison over $r\ne0$ is the original period quotient,
since the resolution is an isomorphism there. There are
twelve components in the cycle because one application of
$\widehat\gamma$ shifts the unfolded component index by
twelve.

The ruled normalization of a component is obtained by adding
the zero and infinity sections to the indicated
$\mathbb C^*$-bundle. For a degree-zero line bundle $L_0$
on an elliptic curve, $\mathbb P(\mathcal O\oplus L_0)$
is a product over that curve only if $L_0$ is trivial.
Indeed an isomorphism would give two independent nowhere
vanishing quotient sections of a common twist of
$\mathcal O\oplus L_0$; degree zero forces that twist and
$L_0$ to be trivial. Equivalently, the two distinguished
disjoint sections have normal bundles $L_0,L_0^{-1}$.
A finite isogeny can trivialize $L_0$ only when $L_0$ is
torsion, because the norm of its pullback is a positive
tensor power of $L_0$.

The section $P$ of the rational elliptic surface is not
torsion, as is checked by its positive Mordell--Weil height
in Engel's Lemma 2.1. If its restriction to the root disc
were identically killed by a positive integer, analytic
continuation on the smooth part of $S\to B$ would give the
same identity globally. Each nonzero torsion equation
therefore has a discrete zero set. Outside a countable union
of such sets, $P_s-O_s$ is non-torsion. This disproves the
product assertion, even after an arbitrary fixed finite
isogeny of the elliptic factor. The special torsion
isogenies at $s_j=0$ used in R11 remain valid.
\end{proof}

\begin{theorem}[B.1.2; Proved]\label{push:B12}
The twisted finite deck generators extend to the models of
B.1.1. On the twelve component labels their actions are
$i\mapsto i+4$ at $j=1$ and $i\mapsto i-3$ at $j=2$,
with labels in $\mathbb Z/12$. Every nonidentity group
element acts freely, including at every node of the central
cycle. The quotient is a smooth complex fourfold, proper
and flat over $(t,r)=(s_j^{m_j},r)$, with a reduced
normal-crossings $r$-central divisor. Its fibre divisor
over $t=p_j$ has multiplicity exactly $m_j$.

The reduction over $s_j=r=0$ is the central surface of R11:
a cycle of four ruled surfaces at $j=1$, and a cycle of
three at $j=2$. In the root presentation the normal
character to $s_j=0$ is $e^{-2\pi i/m_j}$ and the
character on $r$ is trivial.
\end{theorem}
\begin{proof}
We retain source and target arguments. Write primes for
evaluation at $g_js$, put $\zeta_j=e^{-2\pi i/m_j}$,
and set
\[
 R_1=\begin{pmatrix}-1/\tau&0\\(1-\mu)/\tau&1\end{pmatrix},
 \qquad
 R_2=\begin{pmatrix}1/\tau&0\\-\mu/\tau&1\end{pmatrix}.
\]
The identity $R_j\Pi_c=\Pi'_c A_j$ and the fixed-vector
identities $A_jv_j=v_j$ are checked symbolically in
\texttt{ab\_push\_finite\_decks.py}. In the coordinates
$(s,\xi,W,r)$ the two affine deck maps become
\begin{align*}
 s'&=\zeta_1s,\qquad
 \xi'=-\xi/\tau+(6\mu'+2\tau'-4)/3,\\
 W'&=r^4 U_1(s,\xi)W,\\
 U_1&=\exp\!\left[-2\pi i
  \left(\frac{1-\mu}{\tau}\xi+
        \frac{\beta^{{\rm part}\prime}+2\mu'}3\right)\right],
 \end{align*}
and
\begin{align*}
 s'&=\zeta_2s,\qquad
 \xi'=\xi/\tau+(-6\mu'-3\tau'+3)/4,\\
 W'&=r^{-3} U_2(s,\xi)W,\\
 U_2&=\exp\!\left[-2\pi i
  \left(-\frac{\mu}{\tau}\xi-
        \frac{\beta^{{\rm part}\prime}+3\mu'}4\right)\right].
 \end{align*}
Each $U_j$ is a nowhere vanishing holomorphic function on
the additive elliptic cover. Period covariance says
exactly that it defines the required bundle isomorphism
after quotient by $\ker\gamma$; it is not being treated
as a scalar function on the elliptic curve when it is not
one. Multiplication by a unit preserves every toric chart.
Multiplication by $r^a$ sends the chart of component index
$i$ to the chart of index $i+a$. Therefore the displayed
maps and their inverses extend holomorphically to the
resolved infinite chain and then to its twelve-component
quotient. The identities $g_1^3=g_2^4=1$ on the quotient
hold on the dense punctured family, and hence everywhere.

A nonidentity power has component shift $4,8$ in the first
case and $-3,-6,-9$ in the second. None is zero modulo
twelve. The same shifts act on the edge labels of the
cycle, so no node can be fixed either. At $s\ne0$ the
rotation of $s$ already excludes fixed points. The finite
quotient map is consequently a holomorphic local
isomorphism. Its source is smooth, so its target is smooth.
Properness follows from the proper root model and the
finite base map. Flatness can be checked after the
surjective local-isomorphism cover: there $s^{m_j}$ and
$r$ are a regular sequence, and their fibres have pure
dimension two. The local central equation remains
$r=XY$ or $r=X$. Finally $t=s^{m_j}$ shows that the
Cartier divisor over $p_j$ is exactly $m_j$ times its
reduction, including where this divisor meets a central
node. This also proves the asserted normal characters.

At $s=0$ use the explicitly marked isogenies
\[
 L_1(\xi,\eta)=(\eta,\xi-\eta/3),\qquad
 L_2(\xi,\eta)=(\eta,\xi-\eta/2).
\]
Their kernel generators are $(2/3,\rho)$ and $(1/2,i)$,
and the finite deck maps lift to
\[
 (\xi,\eta)\mapsto(\xi+\kappa_1/3,(\rho-1)\eta),
 \qquad
 (\xi,\eta)\mapsto(\xi-\kappa_2/4,-i\eta).
\]
In the Tate coordinate $v=e^{-2\pi i\xi}$ these are the
shifts $4,-3$ just computed. The kernel generators act
by nonzero translations on the smooth elliptic factor
$F_j$, hence freely even at the two fixed points of each
Tate component. Every other element with nonzero deck
power has a nonzero component shift. The full product
groups are therefore free on the product semistable
models. Their component orbits have cardinalities three
and four, giving precisely four and three ruled
components. The original lattice isogenies intertwine
the maps on the dense tori and their unit-interval fans,
so the quotient identifications extend across all nodes.
This verifies the particular comparison with R11 that
we use.

For the punctured family obtained by descending the
$c$-family, the holomorphic correction $\eta_j(s)$ in
R2 causes no additional local obstruction here. At the
shift $c\mapsto c+12$, the reference translations
$-4\Pi_c(\widehat u+\widehat w)$ and
$3\Pi_c(\widehat u+\widehat w)$ are periods.
The remaining correction is equivariant and independent
of $c$. Translation by $-c\eta_j(s)/12$ on the additive
cover conjugates that descent to the period descent.
It commutes with the finite deck action by equivariance.
Thus the constructed finite models fill the actual local
pullback by $w=r^6$ after a specified holomorphic
identification over $r\ne0$. This local gauge is not
asserted to extend across $r=0$ on overlaps with other
base neighbourhoods.
\end{proof}

\subsection{B.2. The base-changed cusp fan}

\begin{proposition}[B.2.1; Proved]\label{push:B21}
For the plain base change $u=r^n$ of the untwisted cusp
model, the smooth period-invariant refinements which
retain the prescribed bottom triangulation, are jointly
flat over $(q,r)$, and have reduced $r$-central divisor
are indexed by $k=0,\ldots,n$. In a fundamental prism
write its four bottom rays as $A,B,C,D$ as in A.2.1,
and put $\rho_i=(0,i,0,1)$ for $0\leq i\leq n$.
The maximal cones for the index $k$ are
\[
 ABC\rho_k,\quad BCD\rho_k,\quad
 AB\rho_i\rho_{i+1}\ (i<k),\quad
 CD\rho_i\rho_{i+1}\ (i\geq k).
\]
There are thirteen such models for $n=12$ and three
for $n=2$. Each has smooth total space and a reduced
normal-crossings central divisor with $n$ components
and no triple locus. Its component labels over the
punctured $q$-disc are globally $\mathbb Z/n$; in
particular their monodromy around $q=0$ is the identity.
\end{proposition}
\begin{proof}
The map of cocharacter lattices for the base change is
\[
 (y_1,y_2,h_q,h_r)\longmapsto
 (y_1,y_2,h_q,nh_r).
\]
The old top ray indexed by $l$ pulls back to
$(0,nl,0,1)$. Smoothness on the two-dimensional top
face forces insertion of every primitive ray
$\rho_{nl+i}$, $0<i<n$: a skipped interval of length
$d>1$ would retain lattice index $d$ on that face.
Joint flatness excludes any new ray with both base
heights positive, and the fixed bottom excludes new
rays of $h_r=0$. Reducedness permits only top rays
of $h_r=1$. These are consequently all the rays.

On the left face $AC\rho_0\rho_n$, a triangulation
with these collinear top vertices has a unique triangle
containing the bottom edge $AC$. Write its third
vertex as $\rho_a$. The triangles on either side
of that bridge are then forced. The analogous right
face has bridge $BD\rho_b$. The first period takes
$A,C$ to $B,D$ and fixes every top ray, so periodicity
requires $a=b=k$. The fixed bottom triangles $ABC$
and $BCD$ must then cone to $\rho_k$. The remaining
successive top intervals force exactly the displayed
cones, leaving no choice. Conversely those cones
have determinant $\pm1$ and common-face
intersections. Their normalized volumes sum to
$n+2$, the volume of the prism, so they cover it.
The determinant, intersection and volume assertions
for both requested exponents are certified in
\texttt{ab\_push\_basechange\_fan.py} using rational
arithmetic.

The period actions are
\[
 G_1:(y_1,y_2)\mapsto(y_1+h_q,y_2),\qquad
 G_2:(y_1,y_2)\mapsto(y_1,y_2+h_q+nh_r).
\]
Thus the choices propagate consistently from the
fundamental prism. The proper toric subdivisions
descend by the covering construction of A.1.2.
Every ray has pure $q$-height or pure $r$-height, so
in a regular affine chart $q$ and $r$ are monomials
in disjoint sets of coordinates. This proves joint
flatness. All nonzero $r$-heights equal one, and
there are at most two top rays in any maximal cone;
this proves reduced normal crossings and absence of
triple points. Properness follows by descent from
the proper toric modification of the base-changed
proper cusp model.

The central components and all their intersection
maps are specified by the stars of the $\rho_i$.
Modulo the periods their indices are $i\bmod n$;
adjacent indices meet in the conductor elliptic
surface. Away from the cusp their normalizations
are the ruled bundles of the expanded plumbing.
At the unique bridge class the star has both bottom
triangles, and gives the nodal-curve blow-up and
small resolution described in A.2.2, on the
corresponding ruled bundle. This star description
also fixes the gluing when the bridge is on an
inserted component. The endpoint choices $k=0,n$
are the canonical resolutions of the base changes
of the two original models. Intermediate values
are related to them by the successive side flops.

Finally, $G_1$ fixes each index and $G_2$ changes it
by $n$. The units in the cusp periods are
single-valued holomorphic functions of $q$ and do
not change a ray label. The toric charts themselves
are defined over $q$, not a root of $q$. Hence
these labels remain global on the punctured
$q$-disc, proving the monodromy assertion. Further
smooth periodic subdivisions with mixed rays give
infinitely many models flat over $r$ alone, as in
A.2.1; they are excluded by the stated joint-flatness
condition.
\end{proof}

\begin{corollary}[B.2.2; Proved; general plain-cusp normalization]\label{push:B22}
For $n=2$ and $n=12$, A.2.4--A.2.5 give all
finite-orbit smooth invariant plain-cusp refinements,
their proper flat families, normalization charts,
finite blow-up and blow-down presentations,
multiplicities and all multiple loci. A.4.7--A.4.9
give their complete integral topology and nearby
Euler function. These statements concern the plain
untwisted cusp. They assert no compatibility with
the finite logarithmic-transform pieces of B.1.
\end{corollary}

\begin{corollary}[B.2.3; Proved; unrestricted toroidal parametrization]\label{push:B23}
For $u=r^{12}$ and $u=w^2$, all finite-orbit smooth invariant
subdivisions of the plain cusp fan are obtained by A.2.4
with $n=12$ and $n=2$, respectively. The construction there
gives their proper flat smooth-total-space families,
central multiplicities, normalization charts and all
multiple loci. Reducedness is exactly the restriction
$\ell=1$. This is a classification of the plain untwisted
cusp models; it does not supply the missing transition to
the finite logarithmic-transform models in B.3.
\end{corollary}

\subsection{B.3. The prescribed pieces do not glue}

\begin{theorem}[B.3.1; Refuted; Lead 5]\label{push:B31}
The finite models of B.1 and any of the plain
$u=r^{12}$ cusp models of B.2 cannot be glued to a
proper complex space over $B\times\Delta_r$ with
the prescribed central fibres on all three base
neighbourhoods. In particular Lead 5's asserted
simultaneous gluing is false. This conclusion does
not rule out a different cusp model of the same
punctured $S^6$ family.
\end{theorem}
\begin{proof}
If a further plain cusp refinement changes the top fan,
its central fibre over an ordinary $q\ne0$ has additional
components or different multiplicities and already cannot
agree with the twelve-cycle model of B.1 on the overlap.
We may therefore assume its top fan is the standard reduced
twelve-cycle. Additional mixed-height rays lie over $q=0$
and do not change the following argument over the punctured
base.

Suppose such a gluing exists and put
$B^\circ=B\setminus\{0,1,\infty\}$. Normalize
its reduced central fibre over $B^\circ$. In
each of the prescribed plumbing charts this is
a proper smooth family whose fibres are the
disjoint union of twelve connected ruled
surfaces. The charts and their transition
isomorphisms therefore give a locally constant
twelve-element set of connected components.
Equivalently, the Stein factor of this
normalization is an unramified degree-twelve
cover of $B^\circ$: the algebra of fibrewise
holomorphic functions is locally $\mathcal O^{12}$,
since each proper connected ruled surface has
only constant holomorphic functions. This
description constructs the cover and requires
no conjecture about components at the missing
three points.

The double curves give the adjacency graph of
a twelve-cycle. Thus its transition
permutations lie in the dihedral group
$D_{12}$, the automorphism group of that graph.
We record a peripheral permutation by its
deck-identification map. Holonomy defined by
lifting a clockwise path and returning its
endpoint to the initial root slice uses the
inverse permutations throughout; the relation
and the following contradiction are unchanged.
At $0$ the monodromy is a rotation by four
steps; at $1$ it is a rotation by minus three
steps, by the explicit root deck maps in
B.1.2. Relabelling a cycle can reverse either
orientation, but changes these only to
$\pm4$ and $\pm3$. At the cusp the
monodromy is the identity by B.2.1.

The peripheral relation in
$\pi_1(B^\circ)$ is $g_1g_2g_0=1$ with
the consistently chosen clockwise loops.
In any common labelling, the product of the
two finite rotations would therefore have
to be the identity. Its possible rotation
numbers are
\[
 4+3,\quad4-3,\quad-4+3,\quad-4-3
 \equiv7,1,11,5\pmod {12}.
\]
None is zero. This contradicts the peripheral
relation. The four residues and the finite
component cycles are checked exactly by
\texttt{ab\_push\_finite\_decks.py}.
\end{proof}

\begin{proposition}[B.3.2; Proved; necessary condition only]\label{push:B32}
If a corrected global construction retains the
finite models and a twelve-cycle over
$B^\circ$, then, in the common orientation
with finite shifts $4,-3$, its cusp
component monodromy must be the twelve-cycle
of shift $-1$. The associated component
cover of $B^\circ$ is connected. Its compact
normalization over $B$ has degree twelve,
ramification indices $3,4,12$, and genus
three. In this deck convention it is the
normalization of
\[
 y^{12}=\frac{t^4}{(t-1)^3}.
\]
These assertions do not construct
the required central complex space.
\end{proposition}
\begin{proof}
The peripheral relation gives the remaining
rotation as $-4+3=-1$. The subgroup generated
by $4,-3$ is all of $\mathbb Z/12$, proving
connectedness. The local cycle decompositions
have four cycles of length three, three
cycles of length four, and one of length
twelve. Filling a punctured-disc cover by
its normal branched cover gives ramification
contribution
$4(3-1)+3(4-1)+(12-1)=28$.
The Riemann--Hurwitz formula gives
$2-2g=24-28=-4$, hence $g=3$.
The displayed Kummer equation has orders
$4,-3,-1$ at the three points and hence
the same marked cyclic monodromy. The
classification of connected covering
spaces identifies its restriction to
$B^\circ$ with the required component
cover; normalization gives the compact
comparison. Reversing the transport
convention replaces $y$ by $y^{-1}$.
Only this component cover has been
constructed; no topology is assigned here
to a hypothetical threefold above it.
\end{proof}

\begin{proposition}[B.3.3; Proved; obstruction to a specified repair]\label{push:B33}
A cyclic cusp-root construction using a
fractional translation by the two cusp
periods cannot repair B.3.1 while retaining
the reduced Mumford fibre for $r\ne0$.
More precisely, let $q=s^d$, $u=r^n$,
and suppose a cyclic deck generator has
translation $d^{-1}(aG_1+bG_2)$ in the
additive notation for the two periods.
Assume its $r$-component shift is integral.
Then its $r$-component shift is
$nb/d\bmod n$, whereas the multiplicity
of every generic quotient $q$-component
is a positive multiple of
\[
 \frac{d}{\gcd(d,a,b)}.
\]
A nonzero $r$-component shift therefore
forces a nonreduced $q$-fibre. In
particular a shift $-1\bmod12$ in this
class of constructions forces multiplicity
at least twelve.
\end{proposition}
\begin{proof}
In the coordinates of A.1.2 the two
periods have $q$-valuation vectors
$(1,0)$ and $(0,1)$; their $r$-valuation
vectors are $(0,0)$ and $(0,n)$.
After $q=s^d$, the period lattice on
the $q$-central component indices is
$d\mathbb Z^2$. A regular toroidal
resolution has indices in
$(\mathbb Z/d)^2$. The proposed deck
generator translates them by $(a,b)$.
The order of that translation is
$d/\gcd(d,a,b)$, so the subgroup
preserving a component has order
$h=\gcd(d,a,b)$. Its subgroup fixing
a generic point has some order
$h_0$ dividing $h$; unit factors may
make $h_0$ smaller. That isotropy
acts on the normal coordinate $s$
by a primitive $h_0$th root of
unity. The quotient normal
coordinate is $s^{h_0}$, and
$q=(s^{h_0})^{d/h_0}$ has
multiplicity $d/h_0$, a positive
multiple of $d/h$. Unit factors
cannot decrease this lower bound.

Reducedness would require $h_0=d$,
and therefore $h=d$,
so both $a,b$ are divisible by $d$.
Then $nb/d$ is zero modulo $n$.
If $nb/d$ is relatively prime to
$n$, its order in $\mathbb Z/n$
is $n$, which divides the order
$d/h$ of the full component
translation; this gives the final
lower bound. The residue calculations
and examples for the requested
exponents are independently recorded
in \texttt{ab\_push\_cusp\_root\_repair.py}.
This rules out the specified repair,
not a cusp model with a different
geometry or a different comparison
to the punctured family.
\end{proof}

\begin{theorem}[B.3.4; Proved; split-normalization trace obstruction]\label{push:B34}
Let $C$ be the genus-three component cover
in B.3.2, and let $\mathscr E\to C$ be
the relatively minimal elliptic surface
obtained by pulling back $S\to B$.
Write $\sigma$ for the order-twelve deck
transformation, with its natural lift to
$\mathscr E$. Then
\[
 \operatorname{Tr}_{\sigma}
   \mathscr E(\mathbb C(C))
 :=\left\{\sum_{i=0}^{11}\sigma^i Q:
                 Q\in\mathscr E(\mathbb C(C))\right\}
 =12\mathbb Z P.
\]
In particular there is no section $Q$
for which $(t_Q\circ\sigma)^{12}=t_{6P}$.
Thus the proposed repair that takes the
normalization over $C$ to be ruled over
this split elliptic surface, with adjacent
conductor map $t_Q\circ\sigma$, is
impossible.

This conclusion is restricted to the
stated split normalization. A ruled
normalization over a nontrivial genus-one
torsor need not admit a section or such
an expression for its affine deck lift;
it is not excluded by this theorem.
\end{theorem}
\begin{proof}
There is a convenient model of the cover
and its elliptic surface:
\[
 C:\ b^3=a^4+1,\qquad
 t=b^3,\quad y=b/a,\qquad
 \mathscr E:\ Y^2=X^3-3bX+2a^2.
\]
Indeed $y^{12}=t^4/(t-1)^3$.
The substitution from the original
Weierstrass coordinates is
$x=a^2b^4X$, $y_{\rm old}=a^3b^6Y$.
The affine discriminant of the new
equation is $1728$; it has no singular
fibre on the affine curve. The curve
$C$ has one point at infinity, with
pole orders $3,4$ for $a,b$.
Put $z=a/b$ and $c=z^3a=z^4b$ there.
Then $z$ is a uniformizer and
\[
 c^3(1-c)=z^{12},\qquad c(0)=1.
\]
In the integral Weierstrass coordinates
$X_0=z^2X$, $Y_0=z^3Y$, the equation is
\[
 Y_0^2=X_0^3-3cX_0+2c^2.
\]
Its discriminant has order twelve and
its $c_4$ is a unit. The unique singular
fibre of the minimal regular surface
is therefore $I_{12}$. The fundamental
line bundle is $\mathcal O_C(\infty)$,
so $\chi(\mathcal O_{\mathscr E})=1$.
This is a projective elliptic surface.

Let $\omega=e^{2\pi i/3}$ and
$\lambda=e^{2\pi i/12}$. The deck lift is
\[
 \sigma(a,b,X,Y)=(ia,\omega b,\lambda^2X,\lambda^3Y).
\]
It fixes the section $P=(0,-\sqrt2a)$.
Two further sections are
\[
 R=\left(\frac{a^2}{b-1},
                 \frac{\sqrt3a}{b-1}\right),
 \qquad
 R'=\sigma R=
   \left(\frac{a^2}{b-\omega},
                 \frac{\sqrt3\omega a}{b-\omega}\right).
\]
Direct substitution verifies their
equations. The three orbit points
$R,\sigma R,\sigma^2R$ lie on the line
\[
 Y=\frac{\sqrt3b}{a}X-\sqrt3a.
\]
Thus $R+\sigma R+\sigma^2R=O$ by the
elliptic chord law, and $\sigma^3R=R$.
The ansatz that produced these sections,
its other solutions, and this line
identity are checked symbolically in
\texttt{ab\_push\_trace\_sections.py}.

We give the intersection computation
needed for the height lattice; this
also fixes the component indices used
below. The only pole of $R$ on the affine
curve is at $(a,b)=(0,1)$. There $a$ is
a uniformizer and $b-1=a^4/3+O(a^8)$,
so $X=3a^{-2}+O(a^2)$ and
$Y=3\sqrt3a^{-3}+O(a)$.
Consequently $R\cdot O=1$.
The same computation at $(0,\omega)$
gives $R'\cdot O=1$, while $P\cdot O=0$.
At the other two points with $a=0$,
$P$ and $R$ meet simply: the fibre is
smooth at $(X,Y)=(0,0)$, and the
difference of their $Y$ coordinates
has a nonzero linear term in $a$.
Thus $P\cdot R=2$, and likewise
$P\cdot R'=2$. The sections $R,R'$
meet simply at $(a,b)=(0,\omega^2)$.
They cannot meet elsewhere on the
affine curve, since their $X$
coordinates are distinct for $a\ne0$.

To check infinity, for the section
$\sigma^jR$ put $k=\omega^j$. Its
integral coordinates are
\[
 X_0=\frac{c^2}{c-kz^4},\qquad
 Y_0=\frac{\sqrt3kc z^4}{c-kz^4}.
\]
At the node use analytic coordinates
\[
 U=Y_0-(X_0-1)\sqrt{X_0+2},\qquad
 V=Y_0+(X_0-1)\sqrt{X_0+2}.
\]
They satisfy
\[
 UV=(1-c)\{3X_0-2(c+1)\},
\]
whose right side is a unit times $z^{12}$.
The section expansions are
\[
 U=-\frac{\sqrt3k^2}{6}z^8+O(z^{12}),
 \qquad V=2\sqrt3kz^4+O(z^8).
\]
Hence $R,R'$ both meet component $4$
of the resolved $I_{12}$ cycle, with
distinct coordinates on that component.
Reversing the cycle numbering replaces
both indices by $8$ and has no effect
on the argument. The section $P$ meets
the identity component, at the smooth
point $(X_0,Y_0)=(0,-\sqrt2)$.
Thus no further intersections occur at
infinity, and $R\cdot R'=1$.

The height pairing formula and its
positivity are those of
\cite[Theorems 8.4 and 8.6, table (8.16)]{Shioda90}.
For an $I_{12}$ fibre its correction is
$\operatorname{contr}(i,j)=\min(i,j)-ij/12$.
The preceding intersections give, in
the ordered basis $P,R,R'$,
\[
 G=\begin{pmatrix}
 2&0&0\\
 0&4/3&-2/3\\
 0&-2/3&4/3
 \end{pmatrix},\qquad \det G=8/3.
\]
In particular these are three independent
sections. The Leray sequence and
$R^1\pi_*\mathcal O_{\mathscr E}
=\mathcal O_C(-\infty)$ show that its
irregularity equals $g(C)=3$; the
right-hand sheaf has no global section.
This also agrees with the Picard-variety
statement in \cite[Theorem 4.1]{Shioda90}.
The equalities $\chi=1$ and
$e(\mathscr E)=12$ give $p_g=3$,
$b_2=22$ and $h^{1,1}=16$.
The zero section, a fibre and the eleven
nonidentity $I_{12}$ components generate
a rank-thirteen trivial lattice.
The identification of sections with
the N\'eron--Severi quotient
\cite[Theorem 1.3]{Shioda90}, together
with $\rho\le h^{1,1}$, now shows that
the Mordell--Weil rank is exactly three
and $\rho=16$.

There is no nonzero torsion section.
Indeed its height would be zero, so
for some $0\le j<12$ the height formula
would give
\[
 0=2+2(Q\cdot O)-\frac{j(12-j)}{12}.
\]
The correction is at most three and
$Q\cdot O$ is a nonnegative integer.
This forces $Q\cdot O=0$ and
$j(12-j)=24$, which has no integral
solution.

Let $L=\mathbb ZP+\mathbb ZR+\mathbb ZR'$.
The integral N\'eron--Severi sublattice
generated by the trivial lattice and
these three sections has determinant
$-32$. This follows either by eliminating
the trivial lattice with the displayed
height matrix, or from the integral
intersection matrix recorded in
\texttt{ab\_push\_trace\_obstruction.py}.
All its entries have just been computed:
sections have square $-1$, their mutual
intersections are $2,2,1$, and $R,R'$
meet precisely component $4$.
If $d=[\mathscr E(\mathbb C(C)):L]$,
the section--N\'eron--Severi quotient
identification gives the same index
for this sublattice in
$\operatorname{NS}(\mathscr E)$.
Its integral determinant implies
$d^2\mid32$, hence $d\in\{1,2,4\}$.
In particular $4Q\in L$ for every section.

In the rational basis $P,R,R'$ the
deck action and its trace are
\[
 M=\begin{pmatrix}1&0&0\\0&0&-1\\0&1&-1\end{pmatrix},
 \qquad \sum_{i=0}^{11}M^i=
       \operatorname{diag}(12,0,0).
\]
Suppose first that the trace of $Q$ is
$6P$. By the index bound we may write
\[
 Q=\frac12P+\frac14(uR+vR'),\qquad u,v\in\mathbb Z.
\]
Let $j\in\mathbb Z/12$ be its component
index at infinity. Since both $R,R'$
meet component four, the height formula
modulo integers gives
\[
 \frac{2u-v}{6}\equiv\frac j3\pmod{\mathbb Z},
 \qquad
 \frac{-u+2v}{6}\equiv\frac j3\pmod{\mathbb Z}.
\]
Equivalently $2u-v-2j$ and $-u+2v-2j$
are divisible by six. Reducing modulo
two shows that both $u,v$ are even.
Write $u=2u'$, $v=2v'$. The component
map is a group homomorphism, so
\[
 2j\equiv4(u'+v')\pmod{12}.
\]
Therefore $j$ is even. The geometric
height formula then lies in
$\tfrac13\mathbb Z$. In contrast the
height from $G$ is
\[
 \langle Q,Q\rangle
 =\frac12+\frac{u'^2-u'v'+v'^2}{3},
\]
which lies in the disjoint coset
$\tfrac12+\tfrac13\mathbb Z$.
This is a contradiction.

The same argument determines the full
trace image. If $Q=\alpha P+\beta R+\gamma R'$
in the rational span, pairing with the
narrow section $P$ gives
$2\alpha=\langle Q,P\rangle\in\mathbb Z$.
Every odd half-integral value of $\alpha$
leads to the identical congruences and
height contradiction, with $1/2$
replaced by an odd square divided by
two. Thus $\alpha\in\mathbb Z$ and
the trace is in $12\mathbb ZP$.
The section $P$ realizes the generator.
The exact matrices, local series,
torsion test, and all relevant residue
classes are verified in
\texttt{ab\_push\_trace\_obstruction.py}.
Finally the identity
$(t_Q\circ\sigma)^{12}=t_{\operatorname{Tr}_\sigma Q}$
proves the claimed obstruction, with
the scope stated in the theorem.
\end{proof}

\begin{corollary}[B.3.5; Proved; exact Mordell--Weil lattice]\label{push:B35}
In the notation of B.3.4,
\[
 \mathscr E(\mathbb C(C))
   =\mathbb ZP\oplus\mathbb ZR\oplus\mathbb ZR',
 \qquad \det\operatorname{NS}(\mathscr E)=-32.
\]
\end{corollary}
\begin{proof}
The affine coordinate ring of $C$ is
$\mathbb C[a,b]/(b^3-a^4-1)$, with
unique representatives of degree at most
two in $b$. The pole orders $3i+4j$ of
$a^ib^j$, $0\le j\le2$, are distinct,
so cancellation cannot create functions
with pole order one or two at infinity.
Consequently
\[
 H^0(C,\mathcal O_C(2\infty))=\mathbb C,
 \qquad
 H^0(C,\mathcal O_C(3\infty))
       =\mathbb C\oplus\mathbb Ca.
\]
A section disjoint from $O$ therefore
has $X=x_0$, $Y=y_0+y_1a$ for constants
$x_0,y_0,y_1$. Substitution in the
Weierstrass equation forces $x_0=y_0=0$
and $y_1^2=2$. Thus the only nonzero
sections with intersection zero with
$O$ are $P,-P$.
Every other nonzero section has
intersection at least one with $O$,
and hence height at least $4-3=1$.

B.3.4 proves that every section has
integral coefficient along $P$ in the
rational height decomposition. The
remaining lattice can have index only
one, two or four over $\mathbb ZR+\mathbb ZR'$.
Any proper such overlattice contains a
nonzero element of order two modulo
this sublattice. After subtracting an
integral combination of $R,R'$ it is
one of $R/2$, $R'/2$, $(R+R')/2$.
Each has height $1/3$, contradicting
the preceding lower bound. Thus the
index is one. The determinant already
computed in B.3.4 is the full
N\'eron--Severi determinant.
The pole-basis calculation, coefficient
equations and three half-class heights
are included in the exact certificate
\texttt{ab\_push\_trace\_obstruction.py}.
\end{proof}

\begin{proposition}[B.3.6; Proved; an unramified torsor alternative]\label{push:B36}
Put $K=\mathbb C(t)$, $L=\mathbb C(C)$,
$G=\operatorname{Gal}(L/K)=\langle\sigma\rangle$,
and let $E/K$ be the generic elliptic
curve of $S\to B$. There is a nonzero
class
\[
 \xi\in\operatorname{Sha}(E_L/L)^G
 \quad\hbox{whose transgression in}\quad
 H^2(G,E(L))=\mathbb ZP/12\mathbb ZP
 \quad\hbox{is }6P.
\]
Here $\operatorname{Sha}$ denotes the
subgroup of genus-one torsors trivial
over the henselian local fields of all
points of $C$.
In particular the trace obstruction of
B.3.4 does not obstruct all genus-one
torsors. Such a class has a relatively
minimal projective genus-one surface
$V\to C$, locally isomorphic to
$\mathscr E\to C$, and an affine lift
$A:V\to V$ of $\sigma$ with
\[
 A^{12}=t_{6P}.
\]
This is an existence statement for a
possible normalization base, not a
construction of a threefold central
fibre or of its smoothing.
\end{proposition}
\begin{proof}
We use Galois cohomology with its
continuous cochains. The five-term
Hochschild--Serre sequence gives
\[
 H^1(K,E)\longrightarrow H^1(L,E)^G
 \xrightarrow{\ d_2\ }H^2(G,E(L))
 \longrightarrow H^2(K,E).
\]
The transgression sequence is
\cite[Chapter III, Section 4, Theorem 2
and the following remark, pp.129--130]{HS53};
its continuous form follows by passing
to finite Galois quotients for the
finitely many values of each continuous
cochain. For this complex function field,
$H^2(K,E)=0$, as in
\cite[Lemma 4.7.8]{CDL24}.
One can also see this from
$\operatorname{cd}K=1$: the torsion
submodule of $E(\overline K)$ has no
second cohomology, and its quotient is
a rational vector space, whose positive
continuous cohomology vanishes by
averaging on finite quotients.

Cyclic group cohomology gives
$H^2(G,E(L))=E(K)/\operatorname{Tr}_\sigma E(L)$.
By B.3.4 this is $\mathbb Z/12$, with
generator the class of $P$. Hence there
is a $G$-invariant torsor class
$\xi_0\in H^1(L,E)$ with $d_2\xi_0=6P$.
We next remove all its local ramification
without changing this transgression.

At a place $v$ of $B$ choose a point
$w$ of $C$ above it and its decomposition
subgroup $G_w$. The only nontrivial
decomposition groups have orders
$3,4,12$, at $0,1,\infty$ respectively.
The component groups of the original
fibres are $\mathbb Z/3$, $\mathbb Z/2$
and zero. Thus $6P$ belongs to the
identity component of the local N\'eron
group at every ramified place.
Multiplication by the ramification
index is surjective on that identity
group over the henselian local field:
on the residue group it is surjective
on $\mathbb G_a(\mathbb C)$ or
$\mathbb G_m(\mathbb C)$, and its
derivative is invertible; Hensel's
lemma lifts a division point.
Choose such a division point over
$K_v^h$. Viewed over $L_w^h$, it is
invariant, so its local norm is $6P$.
The restriction of the transgression
class to $H^2(G_w,E(L_w^h))$ is
therefore zero. For unramified places
the decomposition group is trivial,
so the same assertion is automatic.

Naturality and exactness of the local
five-term sequences now show that
$(\xi_0)_w$ is the restriction of some
$\eta_v\in H^1(K_v^h,E)$.
Only finitely many of these classes
need be nonzero. The original elliptic
surface is nontrivial, since
$j=1728t$. Ogg--Shafarevich localization
is therefore surjective:
\[
 H^1(K,E)\longrightarrow
 \bigoplus_{v\in B}H^1(K_v^h,E)
 \longrightarrow0.
\]
This is the combination of exact
sequence (4.7.4) and Theorem 4.7.9
of \cite{CDL24}; the same surjectivity
and its geometric fibre-twist
interpretation are given in
\cite[Section 4, equation (4)]{Xiao91}.
Take a global $\eta$ with these local
values and put
$\xi=\xi_0-\operatorname{res}_{L/K}\eta$.
It is locally trivial at every point
of $C$, remains $G$-invariant, and has
$d_2\xi=6P$. It is nonzero because its
transgression is nonzero.

The standard interpretation of the
Weil--Ch\^atelet group represents
$\xi$ by a smooth projective genus-one
curve over $L$; taking its relatively
minimal regular model over $C$ gives
$V$. Local triviality identifies each
henselian model with that of $E_L$.
Over $\mathbb C$ the corresponding
\'etale identifications yield
holomorphic identifications over small
base discs. Equivalently, $V$ is obtained
by gluing copies of $\mathscr E$ over
these discs by translations by local
sections. These translations extend
over the regular $I_{12}$ model.
Thus $V$ is locally isomorphic to
$\mathscr E$, has smooth total surface,
no multiple fibres, and the same single
$I_{12}$ fibre. Properness and flatness
hold locally over the base and hence
globally. This use of torsors and local
twists is the one in
\cite[Sections 3--4]{Xiao91}; it does
not assert an explicit equation for
the selected torsor.

The $G$-invariance of $\xi$ gives an
affine isomorphism of its generic fibre
covering $\sigma$. It extends to the
relatively minimal regular surface:
in this particular situation the
extension can be checked on a common
resolution. If an exceptional
$(-1)$-curve for one projection is not
contracted by the other, its image is
a vertical rational curve and has
self-intersection at least $-1$,
since blowing down only increases
self-intersection. On either locally
trivial model all vertical rational
curves are components of the $I_{12}$
fibre, of square $-2$. Thus every such
$(-1)$-curve is contracted by both
projections and can be removed from
the common resolution. Repeating this
proves that the map is an isomorphism.
The same argument gives the local
comparison with the Jacobian model
used above. Transporting the base by
$\sigma$ changes none of these facts.
Write $A_0$ for this lift. Its twelfth
power is translation by a section
$N\in E(L)^G=\mathbb ZP$.
Choosing the powers $1,A_0,\ldots,A_0^{11}$
as lifts of $G$ gives the cyclic
factor set
$N\lfloor(i+j)/12\rfloor$.
Its cohomology class is precisely the
descent obstruction $d_2\xi$, so
$N=(6+12m)P$ for an integer $m$.
Replacing $A_0$ by $A=t_{-mP}\circ A_0$
gives $A^{12}=t_{6P}$.
The finite carry-cocycle identities,
its exact order two modulo twelve,
and the three local component checks
are recorded in
\texttt{ab\_push\_torsor\_cocycles.py}.

Neither this cohomological construction
nor local triviality specifies the
normal line bundles, the threefold
cusp charts, or a smoothing with the
given punctured-family identification.
Those remain requirements of B.4.1.
\end{proof}

\begin{proposition}[B.3.7; Proved; an auxiliary family with the wrong cusp]\label{push:B37}
For the torsor and affine lift $(V,A)$
of B.3.6, there is a line bundle $M$ on
$V$ representing the narrow Jacobian
section $P$, with
\[
 M^2=-2,\qquad
 M\otimes(A^*M)^{-1}
       \simeq\pi^*\mathcal O_C(\infty).
\]
These data give a proper flat nodal
plumbing $\mathcal Z^\flat\to\Delta_h$
and two uniform smooth-total-space
semistable modifications
$\mathcal Z^\pm\to\Delta_h$, together
with holomorphic maps to $B$.
For $h\ne0$, the nodal fibre has
twelve ordinary double points and
Euler number twelve; each small-resolved
fibre has Euler number twenty-four.
The cusp fibre of the latter map to
$B$ has multiplicity twelve.
Consequently this auxiliary construction
does not fill the prescribed $S^6$
family. No proper birational modification
over $B$ can turn its multiple cusp
into a reduced cusp.
\end{proposition}
\begin{proof}
The section $P$ determines a relative
Picard class of multidegree zero on
every fibre component of $V\to C$.
It lifts to a line bundle $M$ because
$\operatorname{Br}(C)=0$; this is the
relative Picard consequence of
\cite[Proposition 0.9.2 and the opening
of Section 4.2, p.370]{CDL24}.
In local trivializations of the torsor,
translation acts trivially on the
relative degree-zero Picard class.
Thus the local Poincar\'e lines glue
up to base lines, with precisely that
Brauer obstruction. In particular we
may and do retain multidegree zero
on all twelve fibre components.

Choose a finite multisection cover of
degree $d$ which trivializes the torsor
generically, and resolve its base change.
There the pulled-back $M$ becomes the
narrow Poincar\'e line for the pulled-back
$P$, up to a base line. Its square is
$-2d$, by height base change
\cite[Proposition 8.12, equation (8.25)]{Shioda90}.
Base lines do not change the square
of a fibre-degree-zero line. The
intersection pullback formula therefore
gives $M^2=-2$.

Since $\sigma P=P$, the line
$J=M\otimes(A^*M)^{-1}$ is trivial
on every fibre, including the
$I_{12}$ fibre. Grauert base change
and the evaluation map give
$J=\pi^*J_C$ for a line on $C$.
The line $M\otimes(t_{6P}^*M)^{-1}$
has base degree twelve. To check the
sign and degree, pass to the same
trivializing cover. On a Jacobian,
the degree of the restriction of
$\mathcal O(P-O)$ to a section $Q$
minus its degree on $O$ is
$-\langle P,Q\rangle$, by the height
formula and the narrowness of $P$.
For $Q=6P$ this is $-12d$.
Descending the degree factor $d$
proves the assertion. Iterating $A^*$
twelve times now gives
$12\deg J_C=12$, so $\deg J_C=1$.

The forms $da/b^2$, $a\,da/b^2$ and
$da/b$ form a basis of $H^0(C,K_C)$.
Their $\sigma$-eigenvalues are
$\lambda^7,\lambda^{10},\lambda^{11}$.
Thus the characteristic polynomial
on $H^1(C,\mathbb Z)$ is
\[
 p(x)=\Phi_{12}(x)\Phi_6(x)
     =x^6-x^5+x^3-x+1,
 \qquad
 p(x)+(1-x)(x^5+x^2+x)=1.
\]
It follows that $1-\sigma^*$ is an
automorphism of $\operatorname{Pic}^0(C)$,
with inverse $(\sigma^*)^5+(\sigma^*)^2+\sigma^*$.
Twisting $M$ by a base line changes
$J_C$ by $N\otimes(\sigma^*N)^{-1}$.
The displayed inverse therefore supplies
a base twist making
$J_C=\mathcal O_C(\infty)$.
Let $\phi$ be its section with divisor
the reduced $I_{12}$ fibre $F_\infty$.

We spell out the comparison with A.1.2
to identify the family actually obtained.
Set $R_V=\mathbb P_V(\mathcal O\oplus M)$,
and identify its zero and infinity
sections by $s_0(x)\sim s_\infty(Ax)$.
On the node neighbourhood over $V$ use
\[
 a\in M_x,\quad b\in M_{Ax}^{-1},
 \qquad a\otimes b=h\phi_x.
\]
The $a$ collar is glued to the annular
core at $(x,a)$; the $b$ collar is glued
at $(Ax,b^{-1})$. Choose the Hermitian
metric and collar radii as in A.1.2.
Compactness of $V$ and the bounded norm
of $\phi$ give disjoint collars,
Hausdorffness and compact representatives
by exactly the same inequalities.
These are holomorphic gluing maps of
fourfolds over $\Delta_h$.
Their local hypersurface rings are flat
over $h$, and the compact representatives
prove properness. The map to $B$ is
well defined because $t\circ\sigma=t$.
No identification with the requested
punctured $S^6$ family is being made.

At a smooth point of $F_\infty$ the
equation is $ab=hz_C$; at one of its
twelve nodes, $z_C=uv$ and it is
$ab=huv$. For $h\ne0$ the latter
has one ordinary double point.
These twelve centres are distinct
even if $A$ rotates the fibre
components. A centre has $a=b=0$
and is in neither annular collar;
no overlap identifies it with the
centre indexed by $Ax$.
The double locus is $V$, not $V/A$.
The two uniform prism subdivisions
of A.2.2 extend over these charts:
the affine action preserves the
orientation of the $I_{12}$ cycle,
and all remaining changes are units.
They give proper small modifications,
smooth total spaces and reduced
normal-crossings $h$-central fibres.
Equivalently their normalizations are
the blow-up of $R_V$ along $F_\infty$
in one section, followed by the twelve
chosen small resolutions. The normal
product becomes
$\pi^*\mathcal O_C(\infty)(-F_\infty)=\mathcal O_V$.

The nearby function of the nodal
plumbing is one on $F_\infty$ and
zero on its complement in the conductor.
Since $e(F_\infty)=12$, its integral
is twelve. Each small resolution
replaces a point by $\mathbb P^1$
and adds one, giving twenty-four.
Independently $e(R_V)=24$; above each
of the twelve nodes the modified
normalization replaces a point by
a chain of two projective lines,
of Euler number three. The open
parts of the fibre are copies of
$\mathbb C^*$ and contribute zero.
The modified normalization therefore
has Euler number $48$, its pinching
has Euler number $36$, and the
semistable nearby count is
$48-2e(V)=24$.

Finally $z=a/b$ is a coordinate on
$C$ near infinity and
$q_B=1/t=z^{12}/c^3$, with $c$ a unit.
At a generic point of the auxiliary
cusp fibre this gives multiplicity
twelve. In fact the map $q_B$ is
locally a twelfth power throughout
the smooth nearby threefold; a
holomorphic root of the unit can
be chosen locally. On any common
resolution of a proper birational
comparison over $B$, every divisor
over that fibre consequently has
$q_B$-valuation divisible by twelve.
A reduced fibre on another smooth
model would have a strict transform
of valuation one, a contradiction.
The finite Euler, polynomial and
multiplicity checks are in
\texttt{ab\_push\_auxiliary\_plumbing.py}.
\end{proof}

\begin{proposition}[B.3.8; Proved; analytic unramified torsors do not evade the Hodge obstruction]\label{push:B38}
Let $V\to C$ be any compact smooth analytic genus-one surface
locally isomorphic, as a torsor over every base disc, to the
Jacobian surface $\mathscr E\to C$ of B.3.4. Algebraicity
or finite order of this analytic torsor class is not assumed.
Then $b_1(V)=6$, and $V$ is K\"ahler. For every holomorphic
line bundle $M$ on $V$, the projective bundle
$\mathbb P_V(\mathcal O\oplus M)$ is K\"ahler. Every smooth
compact threefold bimeromorphic to such a bundle satisfies
the $\partial\bar\partial$-lemma.

Consequently the necessary non-$\partial\bar\partial$
normalization component in B.4.5 cannot be obtained simply
by replacing the projective torsor of B.3.6 by a nonalgebraic
locally trivial analytic torsor, followed by a ruled bundle
and bimeromorphic modifications. This excludes that specified
repair; it does not assert that every possible cusp
normalization has this form.
\end{proposition}
\begin{proof}
Write $p:\mathscr E\to C$ for the projective Jacobian.
The calculation in B.3.4 gives
$R^1p_*\mathcal O_{\mathscr E}=\mathcal O_C(-\infty)$,
and $g(C)=3$. The coherent Leray sequence therefore gives
$h^1(\mathscr E,\mathcal O_{\mathscr E})=3$. Projectivity
and Hodge decomposition imply $b_1(\mathscr E)=6$.
The topological Leray sequence with rational coefficients is
\[
 0\longrightarrow H^1(C;\mathbb Q)
 \longrightarrow H^1(\mathscr E;\mathbb Q)
 \longrightarrow H^0(C,R^1p_*\mathbb Q)
 \longrightarrow H^2(C;\mathbb Q)
 \longrightarrow H^2(\mathscr E;\mathbb Q).
\]
The last arrow is injective because $p$ has a section.
The first arrow is an isomorphism of six-dimensional
vector spaces. Exactness thus implies
$H^0(C,R^1p_*\mathbb Q)=0$.

The local torsor identifications differ by translations.
On a smooth elliptic fibre translations act trivially on
$H^1$. On the $I_{12}$ fibre they preserve the orientation
of the component polygon and therefore act trivially on
its one-dimensional $H^1$ as well. These identifications
commute with the local specialization maps. Hence, for
$\pi:V\to C$, they glue to
$R^1\pi_*\mathbb Q\simeq R^1p_*\mathbb Q$ on the whole
curve, including infinity. The same Leray edge sequence
for $V$ now gives
$H^1(V;\mathbb Q)\simeq H^1(C;\mathbb Q)$.
Thus $b_1(V)=6$ is even. Buchdahl's theorem
\cite[Theorem 11, p.297]{Buchdahl1999} proves that $V$
is K\"ahler without an algebraicity hypothesis.

Choose a Hermitian metric on $\mathcal O\oplus M$.
The induced metric on the tautological relatively ample
line of its projectivization has a closed real $(1,1)$
curvature form $\theta$ positive in the vertical direction.
If $\omega$ is a K\"ahler form on $V$, compactness implies
that $\theta+A\pi^*\omega$ is positive in every direction
for a sufficiently large constant $A$. It is closed, so
it is a K\"ahler form on the projective bundle. A smooth
compact bimeromorphic model is therefore in Fujiki class
$\mathcal C$ and satisfies the $\partial\bar\partial$-lemma;
the domination statement is recalled, with its precise
reference to the $\partial\bar\partial$ theorem, in
\cite[Section 2.1, p.8]{Chen2026HodgeRoot}. Apply B.4.5
for the final obstruction.
\end{proof}


\subsection{B.4--B.5. The remaining filling problem}

\begin{problem}[B.4.1; Open]\label{push:B41}
Construct, or obstruct, a proper flat
holomorphic filling of the actual $S^6$
family after a finite base change which
repairs the cusp incompatibility of B.3.1.
If it exists, determine an equivariant
model for descent to $w$, its central
multiplicities and a resolution, and the
least semistable base-change exponent.
The local constructions in B.1 and B.2
do not settle this problem.
\end{problem}

The obstruction cannot be removed by citing
$H^1(B,\mathcal V)=0$ from R2. That
vanishing solves an additive translation
cocycle for a fixed nonzero parameter and
its integer shifts. On a finite overlap
the translation contains
\[
 \frac{\log s_j}{2\pi i}\,\Pi_c v_j,
 \qquad c=-\frac{12}{2\pi i}\log r.
\]
Its vertical exponential contains a product
of the two logarithms. No assertion in
the cited vanishing controls its growth
at $r=0$ or makes it a unit in the central
toric charts. B.3.1 proves that at least
one of those extension requirements must
fail for the prescribed pieces. A
construction with a corrected cusp remains
necessary before descent or characteristic
numbers can be computed.

\begin{theorem}[B.4.2; Proved; fundamental-group obstruction]\label{push:B42}
Let $f:\mathcal Y\to\Delta$ be a
proper holomorphic map from a smooth
complex manifold, with smooth connected
fibres over $\Delta^*$ and reduced
normal-crossings central fibre $D$.
There is a continuous specialization
map from a nearby fibre $Y_t$ to $D$
which induces a surjection
\[
 \pi_1(Y_t)\longrightarrow\pi_1(D).
\]
No K\"ahler hypothesis is needed.
Hence the central fibre of any such
semistable filling of $S^6$ is simply
connected. In particular it cannot
be a simple pinching of a connected
normalization along one pair of
disjoint nonempty conductor sections,
whether its elliptic normalization
base is split or is a genus-one torsor.
\end{theorem}
\begin{proof}
Equip $\mathcal Y$ and $\Delta$ with
their divisorial logarithmic structures
from $D$ and $0$. Locally, after
absorbing a unit, the map is
$t=z_1\cdots z_k$. Its monoid chart
is the diagonal
$\mathbb N\to\mathbb N^k$; it is
fine, saturated, exact and log smooth.
Thus \cite[Theorem 5.1, p.2226]{NO10}
applies in the complex analytic
category. The induced map of real
oriented logarithmic blow-ups is a
proper topological fibre bundle over
$[0,\epsilon)\times S^1$.
Fix $\theta=\arg t$. Trivializing
this bundle along the radial interval
identifies $Y_t$ homeomorphically with
the fibre $D^{\log}_\theta$ over
$(0,\theta)$. Its natural projection
\[
 p_\theta:D^{\log}_\theta\longrightarrow D
\]
is proper and surjective. Above a
point lying on $k$ local branches,
its fibre consists of the phases
$(\theta_1,\ldots,\theta_k)$ whose
product is the fixed phase $\theta$,
including the harmless phase of the
local unit. It is a torus
$(S^1)^{k-1}$, in particular connected.

Here is a covering-space proof that
$p_\theta$ is surjective on fundamental
groups; path lifting for the collapse
is not assumed. The normal-crossings
space $D$ is locally contractible
and has a connected universal cover
$\widetilde D$. Pull it back along
$p_\theta$. The map from the pullback
to $\widetilde D$ is still proper,
surjective and has connected fibres.
Its total space is connected: a
nontrivial separation into open and
closed sets would, by properness and
connected fibres, give such a
separation of $\widetilde D$.
On the other hand its connected
components are indexed by the cosets
of the image of
$\pi_1(D^{\log}_\theta)$ in $\pi_1(D)$.
There is therefore only one coset.
Composing $p_\theta$ with the chosen
radial homeomorphism proves the claim.

We give an independent elementary
surjectivity argument. Put
$W=f^{-1}(\overline\Delta_\epsilon)$
and $W^*=W\setminus D$. The negative
gradient of $|f|^2$, for a Hermitian
metric, retracts $W$ onto $D$.
For completeness, in a product chart
$f=z_1\cdots z_k$ the arithmetic--
geometric mean inequality gives
\[
 \|\nabla |f|^2\|
   \ge c\,|f|^{2-1/k}.
\]
A finite cover of compact $D$ gives
a common exponent $\alpha<1$ in
$\|\nabla\rho\|\ge c\rho^\alpha$,
where $\rho=|f|^2$; for a fourfold
one may take $\alpha=7/8$.
The remaining trajectory length is
bounded by
$c^{-1}\rho^{1-\alpha}/(1-\alpha)$.
Hence trajectories converge uniformly
to $D$, and parametrization by the
level of $\rho$ gives a continuous
strong deformation retraction fixing
$D$. General position moves every
loop of $W$ off the real-codimension-two
normal-crossings strata, so
$\pi_1(W^*)\to\pi_1(W)$ is onto.
The smooth proper fibration over the
punctured disc gives
\[
 \pi_1(Y_t)\longrightarrow\pi_1(W^*)
       \longrightarrow\mathbb Z\longrightarrow0.
\]
A small transverse meridian at a
smooth point of a reduced component
maps to the generator of $\mathbb Z$
and bounds a disc in $W$.
Thus the image of $\pi_1(Y_t)$ alone
generates $\pi_1(W)=\pi_1(D)$.

For the last assertion, let $R$ be
the connected normalization and
$D_0,D_\infty$ its two disjoint
conductor sections. Choose a
continuous function $h:R\to[0,1]$
equal to zero on $D_0$ and one on
$D_\infty$. The function
$\exp(2\pi i h)$ descends to the
pinching. A path from a point of
$D_0$ to its paired point in
$D_\infty$ becomes a loop whose
image has winding number one.
Consequently the pinching has a
surjection of its fundamental group
onto $\mathbb Z$. This contradicts
the preceding necessary condition
when $Y_t\simeq S^6$.
\end{proof}

\subsection{B.4 supplement. A Hodge obstruction to the central strata}

\begin{theorem}[B.4.3; Known consequences with explicit derivation]
Let $g:\mathcal Z\to\Delta_s$ be a proper holomorphic map from a
smooth complex fourfold, smooth over $\Delta_s^*$, with reduced
normal-crossings central fibre $D$. Suppose its nearby fibres are
biholomorphic to the six-sphere members of the specified family,
after any finite base change $w=s^N$. Then
\[
 H^1(D;\C)=0,\qquad h^1(D,\OO_D)\geq1.
\]
Consequently the natural map
$H^1(D;\C)\to H^1(D,\OO_D)$ is not surjective. In particular,
$D$ cannot be the analytification of a proper complex algebraic
variety with normal-crossings singularities.

If $D$ has \emph{simple} normal crossings, at least one of its
nonempty closed strata fails the $\partial\bar\partial$-lemma.
In particular, no such simple normal-crossings filling can have
all irreducible components K\"ahler or all irreducible components
Moishezon. These obstructions hold for every $N\geq1$.
\end{theorem}
\begin{proof}
The fundamental-group surjection in B.4.2 gives
$\pi_1(D)=1$, and hence $H^1(D;\C)=0$. We check the other
input directly. The audited root-curve direct images in E.2 and
E.3, for $\mathfrak B=\PP(3,4)$, are
\[
 \bar f_*\OO_X=\OO_{\mathfrak B},\qquad
 R^1\bar f_*\OO_X=\OO_{\mathfrak B}\oplus
                          \OO_{\mathfrak B}(-1).
\]
The coarse curve is $\PP^1$, and taking invariants by the finite
stabilizers is exact in characteristic zero. Thus
$H^1(\mathfrak B,\OO)=H^2(\mathfrak B,\OO)=0$.
Moreover $H^0(\mathfrak B,\OO)=\C$ and
$H^0(\mathfrak B,\OO(-1))=0$: in the weighted homogeneous
coordinate ring the latter would require a monomial of weighted
degree $-1$. The degree-one Leray edge sequence therefore gives
\[
 H^1(X,\OO_X)\simeq
 H^0(\mathfrak B,R^1\bar f_*\OO_X)=\C.
\]
This confirms the precise R4 input, rather than inferring
$h^1(\OO)$ from tangent-bundle cohomology. The elementary weighted
degree counts are recorded in \path{code/b_hodge_obstruction.py}.
Upper semicontinuity for the proper flat map $g$ gives
$h^1(D,\OO_D)\geq1$; flatness also follows directly from the
local equation $s=z_1\cdots z_k$.

There is a family-independent check of the required lower bound.
For any compact complex threefold $X$ diffeomorphic to $S^6$,
$c_1(X)=c_2(X)=0$ in integral cohomology, so analytic
Hirzebruch--Riemann--Roch gives
$\chi(\OO_X)=\langle c_1c_2,[X]\rangle/24=0$.
The analytic index theorem used here applies without a K\"ahler
assumption; see \citet[Theorem 4.3]{AtiyahSinger1968IIIE}.
A holomorphic three-form is closed and, since $H^3(X;\C)=0$,
exact. Stokes' theorem applied to its product with its complex
conjugate, whose suitable constant multiple is pointwise
nonnegative, shows that it vanishes. Serre duality then gives
$h^3(\OO_X)=0$. It follows that
$h^1(\OO_X)=1+h^2(\OO_X)\geq1$.

For the algebraic assertion we use the precise properness
hypothesis in \citet[Section 1.D, p.120]{GrafKovacs2014Root}.
For a proper complex algebraic variety with Du Bois singularities,
the canonical map from singular cohomology to coherent cohomology
is surjective in every degree. Normal-crossings singularities are
Du Bois, including self-intersections of a global component.
Applied in degree one, this contradicts the two groups just
computed. No algebraicity of the total space is required for
this contradiction; algebraicity of the central fibre alone
would suffice.

For the simple normal-crossings assertion, suppose instead that
every nonempty intersection of irreducible components of $D$
is a $\partial\bar\partial$-manifold. This is exactly condition
(K1) in \citet[Definition 2.3]{Chen2026HodgeRoot}.
Theorem 2.4(4) of that source, which states the relevant
Steenbrink degeneration theorem under (K1), gives degeneration
at $E_1$ of the nearby Hodge--de Rham spectral sequence.
All its hypotheses have been specified: $g$ is proper,
$\mathcal Z$ is smooth, $D$ is reduced and simple normal
crossings, and its closed strata satisfy (K1). Neither a
K\"ahler metric on $\mathcal Z$ nor a local invariant-cycle
theorem is being assumed. Degeneration in total degree one
would imply
\[
 0=b_1(X_s)=h^0(X_s,\Omega^1_{X_s})+h^1(X_s,\OO_{X_s})
          \geq1,
\]
which is impossible.

Closed smooth submanifolds of a compact K\"ahler manifold are
K\"ahler. Closed strata of a Moishezon component are Moishezon:
a projective modification, followed by taking a component of
the inverse image that dominates the stratum and a general
relative linear section if needed, gives a generically finite
projective cover. Its meromorphic functions show that the
stratum has maximal algebraic dimension. Smooth Moishezon
manifolds satisfy the $\partial\bar\partial$-lemma, as recalled
in \citet[Section 2.1, p.8]{Chen2026HodgeRoot}. Thus either
of the two stated conditions on components would imply (K1).
Finally, finite base change does not change the biholomorphism
type or either cohomological value of a nonzero fibre, so none
of the contradictions depends on $N$.
\end{proof}

This is an application of existing Hodge-theoretic results, not
a claim of a new degeneration theorem. Searches on 15 November
2026 for six-sphere semistable degenerations, normal-crossings
limits, Moishezon central fibres, and Conjecture 8.2 with Hodge
theory did not locate this specific application. We opened
Chen's version 2 of 16 October 2026 and checked Definition 2.3
and Theorem 2.4(4), including the condition inserted before
parts (3)--(4); omitting that condition would make the
statement false. We also opened the original Graf--Kov\'acs
article and checked the proper Du Bois assertion in Section
1.D. The separate source audit records the searches.

The theorem prevents a repair made solely from compact
Moishezon components in a simple normal-crossings model.
It does not prove that an arbitrary analytic filling must have
such components. In particular, no argument here makes an
unknown three-dimensional component over the cusp Moishezon.
The existence and least base-change exponent in B.4.1 therefore
remain open in the unrestricted analytic category.
The Hopf-component example of
\citet[Theorem 1.4 and Example 2.1]{Nishiguchi1988IIRoot}
shows why an analytic degeneration cannot simply be assumed to
have class-$\mathcal C$ strata. The toric-Kato construction of
\citet[Theorem 7.1]{IstratiEtAl2022Root} is another relevant
non-K\"ahler analogue. Neither source supplies the required
cusp comparison for the six-sphere family.

\begin{proposition}[B.4.4 (Proved): extension of the marked fibration]
Suppose $g:\mathcal Z\to\Delta_s$ is a proper holomorphic map
from a smooth complex fourfold, with reduced normal-crossings
central fibre, and its restriction over $\Delta_s^*$ is
isomorphic to a finite base change of the specified six-sphere
family. The given marked holomorphic map
\[
 f^*:\mathcal Z|_{\Delta_s^*}\longrightarrow\PP^1_t
\]
extends to a meromorphic map $\mathcal Z\dashrightarrow\PP^1_t$.
Its graph gives a proper projective modification
$\mathcal Z'\to\mathcal Z$, an isomorphism over
$\Delta_s^*$, on which the marked map is holomorphic.
No assertion that $\mathcal Z'$ is smooth or semistable is
included in this statement.
\end{proposition}
\begin{proof}
We shrink the disc whenever necessary. Put
$L=K_{\mathcal Z/\Delta_s}^{-2}$. The relative canonical
bundle here is the line bundle
$K_{\mathcal Z}\otimes g^*K_{\Delta_s}^{-1}$; its
restriction to a nonzero fibre is that fibre's canonical bundle.
The precise R1 input is
$K_{X_s}^{-2}\simeq f_s^*\mathcal O_{\PP^1}(1)$,
with connected fibres and no further sections. Therefore
$h^0(X_s,L|_{X_s})=2$, and its complete linear system is
the given map $f_s$, up to a projective change of coordinate.
The three critical values are intrinsically distinguished by
their fibres: multiplicity three, multiplicity four, and reduced
normal crossings, respectively.

Grauert's proper direct-image theorem
\citep[Section 6, Hauptsatz I]{Grauert60push} makes
$\mathcal F=g_*L$ coherent. It is torsion-free: multiplication
by a nonzero base germ is injective on $L$, because the local
semistable equation makes $g$ flat, and hence is injective on
its direct image. A coherent torsion-free sheaf on the smooth
disc is locally free. Its rank is two. On the punctured disc,
constant $h^0=2$ and the analytic base-change theorem checked
in E.3 identify its fibres with the complete spaces of sections.
Choose a frame $(\sigma_0,\sigma_1)$ of $\mathcal F$.
Evaluation gives two global holomorphic sections of $L$ on
$\mathcal Z$, and thus a meromorphic map
$F_0=[\sigma_0:\sigma_1]$ to $\PP^1$.

For completeness its graph is analytic and proper without an
appeal to an unproved extension of a punctured graph. In a local
trivialization of $L$, write the sections as $a,b$. Remove their
common divisorial factor in the regular local ring of the smooth
space $\mathcal Z$. The equation
$v a-u b=0$ in $\mathcal Z\times\PP^1_{[u:v]}$
then has a unique component dominating $\mathcal Z$; it is
the closure of the graph where $(a,b)\ne(0,0)$. These local
components glue. Denote the resulting reduced graph by
$\mathcal G$. The projection $\mathcal G\to\mathcal Z$
is proper and projective, and is an isomorphism over
$\Delta_s^*$. The map
\[
 H:\mathcal G\longrightarrow\PP^1\times\Delta_s,
 \qquad H=(F_0,g),
\]
is proper. Locally $\mathcal G$ is a hypersurface in the smooth
space $\mathcal Z\times\PP^1$. The vanishing of the maximal
Jacobian minors for its equation and the two target coordinates
defines the closed analytic locus $C$ where $H$ is not a
submersion; singular points of $\mathcal G$ are included.

The image $H(C)$ is a closed analytic subset of the target.
One may deduce this precise proper-image assertion from the same
direct-image theorem: the coherent sheaf
$H_*\mathcal O_C$ has support exactly $H(C)$, since the
constant section $1$ is nonzero at every nonempty proper fibre.
Over $\Delta_s^*$ this image consists exactly of three distinct
critical-value sections. It cannot have a two-dimensional
component, since such an analytic component would contain an
open set of the connected two-dimensional target and hence meet
the punctured region in dimension two. Remove any component
contained in $\PP^1\times\{0\}$. The remaining reduced
analytic curve is proper with zero-dimensional fibres over the
disc, and is therefore finite over it.

Each of the three labelled sections on the punctured disc belongs
to a degree-one branch of that finite curve. Indeed their labels
are intrinsic and cannot be permuted by continuation. The
normalization of a degree-$d$ branch over a punctured disc has
cyclic monodromy of length $d$; a single-valued labelled section
therefore has $d=1$. The normalization extends it holomorphically
to a map from the full disc to $\PP^1$. Let these three maps be
$c_1(s),c_2(s),c_0(s)$, in the multiplicity-three,
multiplicity-four and cusp order. They may collide at $s=0$.

For $s\ne0$ the unique projective transformation sending
$(c_1,c_2,c_0)$ to $(0,1,\infty)$ has meromorphic coefficients
at $0$. In a common affine chart its value at $z$ is
\[
 A_s(z)=
 \frac{(z-c_1(s))(c_2(s)-c_0(s))}
      {(z-c_0(s))(c_2(s)-c_1(s))};
\]
homogeneous brackets give the same formula if a value is infinite.
The denominators are not identically zero, since the three values
are distinct on the punctured disc. Thus $A_s\circ F_0$ is
meromorphic on $\mathcal Z$. On every nonzero fibre it agrees
with the actual marked $f_s$, because both are the complete
bi-anticanonical map and have the same three labelled values.
This equality is an equality of maps over the punctured disc,
not just a fibrewise biholomorphism statement.

Finally the graph of $A_s\circ F_0$ is constructed by the same
local two-section argument after clearing powers of the base
coordinate in its meromorphic coefficients. Its projection to
$\mathcal Z$ is the asserted proper projective modification,
and its second projection is the required holomorphic marked map.
\end{proof}

The proposition removes a possible essential projective-coordinate
ambiguity. It does not extend a punctured isomorphism to any of
the auxiliary models of B.3, and does not make a component over
$t=\infty,s=0$ Moishezon. Those would be additional statements.

\begin{theorem}[B.4.5 (Proved): the obstruction on normalization components]
Let $D$ be a compact reduced complex space with normal-crossings
singularities, allowing self-intersections of global components.
Suppose that every connected component of its smooth normalization
satisfies the $\partial\bar\partial$-lemma. Then the natural map
\[
 H^1(D;\C)\longrightarrow H^1(D,\OO_D)
\]
is surjective.

Consequently every smooth-total-space reduced normal-crossings
filling of the specified six-sphere family, after any finite base
change, has a normalization component which fails the
$\partial\bar\partial$-lemma. In particular its normalization
cannot consist entirely of compact K\"ahler or Moishezon manifolds.
This conclusion includes normal crossings which are not simple.
\end{theorem}
\begin{proof}
We give a degree-one argument directly on $D$; no nearby
Hodge-de Rham degeneration theorem is needed. Let
$D^{[p]}$ be the smooth compact branch space whose points are
ordered $(p+1)$-tuples of distinct local branches through one
point of $D$. Equivalently, we normalize the loci where these
branches meet, retaining their ordering. These are distinct
\emph{local branches}, even when they belong to one global
irreducible component. The finite holomorphic maps
$\pi_p:D^{[p]}\to D$ and their face maps are specified by
forgetting one branch. The symmetric group acts by reordering;
we take the alternating summand, with the sign character, in
each degree. Averaging over this finite group is exact over
$\C$. Write $\delta$ for the alternating sum of the face
pullbacks.

Both augmented complexes
\[
 \C_D\longrightarrow
 \bigl((\pi_p)_*\C_{D^{[p]}}\bigr)_{p\geq0}^{\mathrm{alt}},
 \qquad
 \OO_D\longrightarrow
 \bigl((\pi_p)_*\OO_{D^{[p]}}\bigr)_{p\geq0}^{\mathrm{alt}}
\]
are resolutions. This can be checked locally without a global
ordering. In a chart $z_1\cdots z_k=0$, the constant complex is
the augmented cochain complex of the simplex on the $k$ branches.
For the holomorphic complex, fix a monomial in the convergent
power series ring. It survives on exactly the intersections
indexed by subsets of those $i$ for which its $z_i$-exponent is
zero. If there are no such indices it already vanishes on $D$;
otherwise its coefficient complex is again an augmented simplex
complex. The usual contraction is a finite combination of
coordinate restrictions, so preserves convergence. This proves
exactness for analytic functions, not just for formal series.
The descriptions agree under permutations of the local branches.
The finite simplex and permutation identities for the at most
four branches in our threefold application are checked exactly
in \path{code/b_nc_degree_one.py}.

Resolve the holomorphic terms by smooth Dolbeault forms and
the constant terms by smooth de Rham forms on the branch spaces.
These compute the respective hypercohomologies. Indeed finite
pushforward is exact here, and partitions of unity from local
ambient smooth functions pull back to make the resulting sheaves
fine. The ordinary de Rham lemma and Dolbeault lemma apply on
each smooth branch space. Projection to the $(0,\bullet)$
part defines the natural comparison between the two total
complexes and induces $H^1(D;\C)\to H^1(D,\OO_D)$.

Represent a coherent degree-one class by a total cocycle
$(\alpha,f)$, where $\alpha$ is a smooth $(0,1)$-form on
$D^{[0]}$ and $f$ is an alternating smooth function on
$D^{[1]}$. Choose the total-complex signs so that its equations
are
\[
 \bar\partial\alpha=0,\qquad
 \delta\alpha=\bar\partial f,\qquad \delta f=0.
\]
On a normalization component, $\partial\alpha$ is of type
$(1,1)$, is $\partial$-exact and is $\bar\partial$-closed.
The $\partial\bar\partial$-lemma gives a smooth function $h$
with $\partial\alpha=\partial\bar\partial h$.
Do this on every component, and put
\[
 \alpha'=\alpha-\bar\partial h,\qquad
 f'=f-\delta h.
\]
The pair represents the same coherent class and still satisfies
the displayed cocycle equations; now $d\alpha'=0$.
Holomorphic face maps commute with both differential operators,
so
\[
 \partial\bar\partial f'
       =\partial\delta\alpha'=\delta\partial\alpha'=0.
\]
A smooth complex-valued pluriharmonic function on a compact
connected complex manifold is constant: apply the maximum
principle to its real and imaginary parts in local complex
coordinates. Therefore $f'$ is constant on each connected
component of $D^{[1]}$. It follows that
$df'=0$ and $\delta\alpha'=0$.

Thus the same pair $(\alpha',f')$, with $\alpha'$ viewed as
a complex-valued smooth one-form, is a cocycle of the total
de Rham branch complex. Its image is the original coherent
class. This proves surjectivity. Notice that no
$\partial\bar\partial$ hypothesis on $D^{[1]}$ or higher
branch spaces was used; compactness suffices for the only
function argument there.

For the asserted application, B.4.2 and B.4.3 give
$H^1(D;\C)=0$ and $h^1(D,\OO_D)\geq1$.
Surjectivity is impossible, so some normalization component
fails the lemma. Compact K\"ahler and smooth Moishezon
manifolds satisfy it, as in the primary comparison checked
in B.4.3. The same argument applies after every finite base
change, since its smooth fibres are still the specified
six-spheres.
\end{proof}

This strengthens the simple-normal-crossings obstruction in
B.4.3 by using only degree one. It also avoids an unverified
extension of an algebraic Du Bois theorem to arbitrary analytic
spaces. Searches for normal-crossings cohomology with K\"ahler
normalization, the normalization Dolbeault complex, and
degree-one surjectivity located the broader Hodge-theoretic
setting discussed by \citet{Chen2026HodgeRoot} and the proper
algebraic statement of \citet{GrafKovacs2014Root}; we do not
claim priority for this elementary consequence. The proof above
states and verifies the particular analytic comparison required
here. It still does not prove that all possible analytic
normalization components satisfy its hypothesis.

\begin{proposition}[B.4.6; Proved; a restricted non-K\"ahler cap obstruction]\label{push:B46}
Let $D$ be a compact reduced normal-crossings complex
space. Suppose its irreducible components have been
partitioned into closed unions $D_0,\ldots,D_s$, with
their reduced induced structures. Assume that
\[
 H^1(D_i,\mathbb C)\longrightarrow H^1(D_i,\mathcal O)
 \quad\hbox{is surjective for every }i,
 \tag{B.46a}
\]
and that the analogous map is injective for every
pairwise intersection $D_i\cap D_j$. Then
\[
 H^1(D,\mathbb C)\longrightarrow H^1(D,\mathcal O_D)
 \quad\hbox{is surjective}.
 \tag{B.46b}
\]

In particular these hypotheses hold for the following
specified cap assembly. The horizontal block $D_0$ is
a two-section pinching of a ruled bundle over an
analytic locally trivial torsor $V\to C$ of B.3.8,
with conductor map $A$ covering $\sigma$ and
$A^{12}=t_{6P}$. The normalization may undergo the
modifications of B.3.7, or any smooth bimeromorphic
modifications for which the two conductor restriction
maps on $H^1$ are the same pullbacks as before.
The other blocks are smooth hyperbolic toric Kato
threefolds or compact smooth spaces satisfying
(B.46a), for example class-$\mathcal C$ spaces.
Assume every pairwise overlap is either a reduced
union of invariant divisors of a hyperbolic toric
Kato threefold, or a compact normal-crossings surface
whose normalized components and normalized double
curves are rational. No such assembly can be a
semistable central fibre of the $S^6$ family.

This excludes the stated cap construction even when
an invariant Kato surface in an overlap kills a
fundamental-group generator of a Kato component.
It is not a classification of arbitrary analytic
cusp components or of their possible overlaps.
\end{proposition}
\begin{proof}
Write $D^{[p]}$ for the disjoint union of intersections
of $p+1$ distinct blocks. The augmented closed-cover
complexes resolve both $\mathbb C_D$ and $\mathcal O_D$.
For the coherent assertion this can be checked in a
normal-crossings coordinate chart. The local equation
of $D$ is a product of distinct coordinates, and each
block groups a disjoint subset of those factors.
The corresponding ideals are generated by products
in disjoint sets of variables. The ordinary alternating
intersection complex is exact, as one sees on each
monomial, where it is the augmented complex of a
simplex. The same contraction acts coefficientwise
on convergent power series. For constant sheaves,
small intersections are contractible unions of
coordinate subspaces, and the same stalkwise simplex
complex is exact. This proves both resolutions also
when the global blocks are self-crossing.

Their hypercohomology spectral sequences have
\[
 E_1^{p,q}(\mathcal F)=H^q(D^{[p]},\mathcal F),
 \qquad \mathcal F=\mathbb C\ \hbox{or}\ \mathcal O.
\]
Every connected compact reduced branch intersection
has only constant holomorphic functions. Thus the
comparison maps in the entire $q=0$ row are
isomorphisms. They induce isomorphisms on
$E_2^{1,0}$ and $E_2^{2,0}$. Surjectivity in (B.46a)
and injectivity on pairwise intersections imply
surjectivity on $E_2^{0,1}$: lift a coherent
degree-one class on the blocks, then its topological
restriction difference vanishes because its coherent
image vanishes and the pairwise comparison is
injective.

Now use the natural five-term exact sequences
\[
 0\longrightarrow E_2^{1,0}
 \longrightarrow H^1(D,\mathcal F)
 \longrightarrow E_2^{0,1}
 \xrightarrow{d_2} E_2^{2,0}.
\]
The isomorphism on the last term makes a lift of a
coherent $d_2$-closed element topologically
$d_2$-closed. The isomorphism on the first term then
gives (B.46b). If the block comparisons are
isomorphisms, the same argument also proves
injectivity, provided the map on $E_2^{0,1}$ is
injective.

We verify the horizontal-block assertion rather than
apply the pairwise criterion directly to its two
conductor branches. By B.3.8,
$H^1(V,\mathbb C)=H^1(C,\mathbb C)$ and
$H^1(V,\mathcal O)=H^1(C,\mathcal O_C)$.
The first equality is the Leray calculation there;
the second also follows from
$R^1\pi_*\mathcal O_V=\mathcal O_C(-\infty)$,
which has no global section. The action of $A$ on
these groups is $\sigma^*$, since it covers $\sigma$
and the displayed maps are pullbacks from the base.
The characteristic polynomial
\[
 p(x)=\Phi_{12}(x)\Phi_6(x),\qquad p(1)=1,
\]
shows that $1-A^*$ is invertible on both groups.
The normalization sequence of the two-section
pinching, for constants and for holomorphic
functions, therefore gives
\[
 H^1(D_0,\mathbb C)=\mathbb C,
 \qquad H^1(D_0,\mathcal O)=\mathbb C.
\]
Both classes come from the same difference map on
constant functions of the conductor. Their natural
comparison is an isomorphism. Projective-bundle
pullback identifies the normalization groups with
those of $V$. The stipulated modifications preserve
these groups and both restriction maps; in particular
the explicit modifications of B.3.7 have this property.
This grouped computation is essential: on an individual
conductor $V$, the comparison has dimensions six and
three and is not injective.

For a hyperbolic toric Kato manifold $Z$, its
fundamental group is $\mathbb Z$, as in
\cite[Section 2, p.1352]{IstratiEtAl2022Root};
\cite[Proposition 8.4 and Theorem 8.5]{IstratiEtAl2022Root}
give $H^0(Z,\Omega_Z^1)=0$ and
$H^1(Z,\mathcal O_Z)=\mathbb C$.
The exact sequence
$0\to\mathbb C\to\mathcal O\xrightarrow{d}d\mathcal O\to0$
shows that $H^1(Z,\mathbb C)\to H^1(Z,\mathcal O)$
is injective, because a global section of
$d\mathcal O$ is a closed holomorphic one-form.
The two dimensions are one, so it is an isomorphism.
The same reasoning applies to hyperbolic toric Kato
surfaces. Compact smooth rational varieties have
both groups zero. Class-$\mathcal C$ blocks satisfy
(B.46a) by their $\partial\bar\partial$ property.

Finally let $F$ be a reduced union of invariant
divisors in a hyperbolic toric Kato threefold.
Its normalized components are compact toric surfaces
or hyperbolic toric Kato surfaces
\cite[proof of Theorem 8.8, p.1380]{IstratiEtAl2022Root}.
Its normalized double curves are rational, since
every invariant curve is rational in the hyperbolic
case \cite[Theorem 5.2]{IstratiEtAl2022Root}.
Use the ordered-local-branch resolution from B.4.5.
The comparisons on component $H^1$ are isomorphisms,
and the groups on normalized double curves vanish.
The same five-term argument shows that
$H^1(F,\mathbb C)\to H^1(F,\mathcal O_F)$ is an
isomorphism. For an overlap with rational normalized
components and double curves, the argument is
identical with zero component $H^1$. Thus every
pairwise overlap in the statement has the required
injective comparison, proving (B.46b) for the
specified assembly.

For an actual semistable $S^6$ centre, B.4.2 gives
$H^1(D,\mathbb C)=0$, while upper semicontinuity
and the nearby value $h^1(\mathcal O)=1$ in B.4.3
give $h^1(D,\mathcal O_D)\geq1$. This contradicts
(B.46b). The obstruction uses the cohomology maps
on the overlaps; it does not assume that every
Kato component's loop survives in the union.
\end{proof}


\paragraph{The remaining analytic comparison.}
We also examined whether general results on degenerations of tori
would force an arbitrary filling into the excluded toroidal class.
The distinction between a construction theorem and such a comparison
theorem is essential. \citet[Theorem 5.1.3 and Sections 5.2--5.4]
{KKN2008Root} construct proper analytic models of log complex tori
without assuming a polarization. Their input is a specified log
complex torus and an admissible pairing; the theorem does not assert
that every proper analytic filling is one of these models. We have
not constructed that global log input with the required finite-place
and cusp markings for the six-sphere family.

The classification in \citet[Theorem 3.1]{KamenovaToriRoot} assumes
a K\"ahler total space and trivial canonical bundle. Those hypotheses
have not been proved for a slice of an arbitrary filling here.
\citet[Section 3, Proposition 3.1]{SchroeerSiebert2002Root} explicitly
study normal-crossings degenerations with Hopf normalization, but
their smooth fibres are primary Kodaira surfaces. This is a source
of alternative analytic models to test, not a comparison with our
punctured family.

Even extending the marked map in B.4.4 does not extend a chosen
punctured isomorphism to a second model. Properness alone cannot
justify that inference: the automorphism
$(s,[z_0:z_1])\mapsto(s,[e^{1/s}z_0:z_1])$ of
$\Delta^*\times\PP^1$ has no meromorphic extension across zero,
although both product families are proper over the disc. Its
restriction to $[1:1]$ has an essential singularity. Thus an
unrestricted proof still needs a specific extension or boundedness
argument for the comparison, or an actual non-toroidal cusp
construction with the prescribed punctured restriction.


\begin{proposition}[B.4.7 (Proved): a parabolic analytic model and its comparison failure]
There is a proper semistable analytic family of compact complex
surfaces $p:\mathcal H\to\Delta$ with smooth total space such that:
\begin{enumerate}
\item its punctured restriction is a holomorphic product with the
scalar primary Hopf surface
\[
 H_-=(\C^2\setminus\{0\})/\langle -e^{-\pi}\id\rangle;
\]
\item its central fibre is $Y=H_1\cup_E H_2$, where both $H_i$
are isomorphic to
\[
 H_+=(\C^2\setminus\{0\})/\langle e^{-2\pi}\id\rangle,
 \qquad E=\C/(\Z+i\Z);
\]
\item $N_{E/H_1}\simeq N_{E/H_2}\simeq\mathcal O_E$, and
\[
 \pi_1(Y)=0,\qquad h^1(Y,\mathcal O_Y)=1,
 \qquad H_*(Y;\Z)=(\Z,0,0,\Z^3,\Z^2);
\]
\item no punctured product isomorphism has a meromorphic extension
whose graph is proper over the two total spaces near the central
fibres.
\end{enumerate}
Consequently $\mathcal H\times\PP^1\to\Delta$ is an actual
semistable family of compact threefolds with simply connected centre
and $h^1(\mathcal O)=1$, whose normalized components fail the
$\partial\bar\partial$ lemma.  Its smooth fibres have Euler number
zero and fundamental group $\Z$.  It is not a filling of the
six-sphere family.
\end{proposition}
\begin{proof}
We construct the family by an explicit instance of Fujimoto's
generalized logarithmic transformation
\citep[Definition 1.1, pp.~85--86]{Fujimoto1990AB}.
Let
\[
 b:\mathcal B=\operatorname{Bl}_{(p,0)}(\PP^1\times\Delta)
       \longrightarrow\PP^1\times\Delta.
\]
The central curve is $F_1\cup F_2$, where $F_1$ is the strict
transform of $\PP^1\times\{0\}$ and $F_2$ is the exceptional
projective line.  They meet transversely in one point.  The two
line bundles
\[
 L=b^*\operatorname{pr}_1^*\mathcal O_{\PP^1}(1),
 \qquad A=\mathcal O_{\mathcal B}(F_2)
\]
have degree-restriction matrix
\[
 \begin{array}{c|rr}
       &L&A\\\hline
 F_1&1&1\\
 F_2&0&-1
 \end{array}.
\]
In particular this matrix is unimodular.  Take nonzero holomorphic
transition functions $l_{ab},a_{ab}$ for these two line bundles.
On small intersections choose logarithms and put
\[
 t_{ab}=\left[
       \frac{(1+i)\log l_{ab}-i\log a_{ab}}{2\pi i}
                         \right]\in E.
\]
This is independent of the logarithm choices, since $1+i,-i$
belong to the lattice of $E$.  On triple intersections the sum
is again a lattice element.  The $t_{ab}$ therefore form a
holomorphic cocycle with values in the translation group $E$.
Glue $U_a\times E$ by these translations, obtaining a
holomorphic principal $E$-bundle $\mathcal H\to\mathcal B$.
It is a proper submersion, since it is locally a product with a
compact elliptic curve.  Its total space is smooth and Hausdorff.
Compose with $\mathcal B\to\Delta$.

The Chern vector of the bundle, in the lattice $\Z+i\Z$, is
\[
 \kappa=(1+i)c_1(L)-i c_1(A).
\]
Its restrictions to $F_1,F_2$ are respectively $1$ and $i$.
On a nonzero fibre its restriction is $1+i$.  These statements
also follow directly from the displayed transition functions.
Locally at the crossing of $F_1$ and $F_2$ the base map is
$xy=u$, and the principal elliptic bundle is locally trivial.
The composed family is thus locally $(x,y,z)\mapsto xy$,
with $z$ a coordinate on $E$.  It is flat, has smooth total
space and reduced simple normal crossings centre.  It is proper
because both maps in the composition are proper.

We identify each surface explicitly.  For a primitive vector
$c\in\Z+i\Z$, choose $d$ so that $(c,d)$ is an oriented
lattice basis.  The homomorphism
\[
 \C^*\longrightarrow E,\qquad
 z\longmapsto\left[\frac{c\log z}{2\pi i}\right]
\]
has kernel generated by
$\alpha_c=\exp(2\pi i d/c)$, with $|\alpha_c|<1$.
The principal $E$-bundle on $\PP^1$ associated to
$\mathcal O(1)^\times$ by this homomorphism is the quotient
of $\mathcal O(1)^\times$ by that cyclic kernel.  Fibrewise
inversion identifies $\mathcal O(1)^\times$ with
$\mathcal O(-1)^\times\simeq\C^2\setminus\{0\}$;
inverting the cyclic generator does not change its group.
The resulting surface is the scalar Hopf surface $H_{\alpha_c}$.
For the three vectors in the construction we may take
\[
 (c,d)=(1,i),\quad(i,-1),\quad(1+i,i),
\]
giving respectively
\[
 \alpha_1=\alpha_i=e^{-2\pi},\qquad
 \alpha_{1+i}=\exp\bigl(2\pi i(1+i)/2\bigr)=-e^{-\pi}.
\]
On $\mathcal B|_{\Delta^*}=\PP^1\times\Delta^*$,
$A$ has the canonical nowhere-zero section of its exceptional
divisor.  Its cocycle is therefore a coboundary there, while
$L$ is the constant $\mathcal O(1)$ bundle.  The construction
gives an actual holomorphic product $H_-\times\Delta^*$,
not merely an isotrivial set of fibres.

Each central $H_i\to F_i$ is a principal elliptic bundle.
Its intersection with the other component is a smooth fibre,
so a local trivialization over a base disc gives
$N_{E/H_i}=\mathcal O_E$ exactly.  Thus the product of the
two normal lines is trivial, with the smoothing section supplied
by the coordinate $u=xy$.

Here are the integral topology maps.  A scalar Hopf surface is
diffeomorphic to $S^1\times S^3$: divide the radial coordinate
by multiplication by $|\alpha|$ and untwist its constant phase
rotation along the radial interval.  Its fibre inclusion has
fundamental-group map
\[
 \Lambda=\Z^2\longrightarrow\Lambda/\Z c_i.
\]
For $c_1=(1,0)$ and $c_2=(0,1)$, these are the two coordinate
quotients.  The van Kampen pushout kills both generators of
$\Lambda$, and hence $\pi_1(Y)=0$.  In the normalization
Mayer--Vietoris sequence the degree-one map, with the second
inclusion negated, is
\[
 H_1(E;\Z)\longrightarrow H_1(H_1;\Z)\oplus H_1(H_2;\Z),
 \qquad \begin{pmatrix}0&1\\-1&0\end{pmatrix}.
\]
It is unimodular.  The degree-zero map is the primitive column
$(1,-1)^t$, the degree-two target is zero, and $H_3(E)=0$.
The exact sequence now gives
\[
 H_1Y=H_2Y=0,\quad
 0\longrightarrow\Z^2\longrightarrow H_3Y
       \longrightarrow\Z\longrightarrow0,
 \quad H_4Y=\Z^2.
\]
All groups displayed in the proposition are therefore free.

For coherent cohomology, translations act trivially on
$H^1(E,\mathcal O_E)$, so for either central principal bundle
the sheaf $R^1\pi_*\mathcal O$ on $\PP^1$ is the constant
line bundle.  Its Leray sequence shows that restriction
\[
 H^1(H_i,\mathcal O_{H_i})\longrightarrow H^1(E,\mathcal O_E)
\]
is an isomorphism.  The normalization exact sequence therefore
has degree-one map $(s,t)\mapsto s-t$ from $\C^2$ to $\C$.
The preceding degree-zero map is surjective.  Its kernel has
dimension one, proving $h^1(Y,\mathcal O_Y)=1$.
This also displays exactly why the overlap-injectivity hypothesis
in B.4.6 fails: $H^1(E;\C)\to H^1(E,\mathcal O_E)$ has a
one-dimensional kernel.

For completeness, the two scalar Hopf surfaces used above are
not biholomorphic.  An isomorphism lifts to an automorphism of
$\C^2\setminus\{0\}$ that intertwines the two cyclic generators
or their inverses.  Hartogs extension applies to the lift and
its inverse; both fix zero and have invertible derivative there.
Differentiating the intertwining identity gives
$e^{-2\pi}=(-e^{-\pi})^{\pm1}$, which is false.  They have
no rational curves, since a map $\PP^1\to H_\alpha$ lifts to
$\C^2\setminus\{0\}$ and is constant.  The usual resolution
of a bimeromorphic map of smooth compact surfaces therefore
also shows that these two Hopf surfaces are not bimeromorphic.
This agrees with the explicit scalar case of
\citet[Section 2.2]{AbashevaDeev2024AB}.

There is an even shorter argument excluding a meromorphic
extension of the punctured product comparison.  A sufficiently
small closed semistable neighbourhood in $\mathcal H$ retracts
onto $Y$, by the monomial-gradient argument of B.4.2.  Its
fundamental group is zero.  The product neighbourhood
$H_-\times\Delta$ has fundamental group $\Z$.  If the
punctured comparison had a proper meromorphic graph, resolving
that graph would give a common smooth proper modification of
the two neighbourhoods.  A proper bimeromorphic map between
smooth complex manifolds preserves $\pi_1$, by the
codimension-two avoidance argument in A.4.7.  This is a
contradiction.  Both models already have the indicated
holomorphic map to the same $\PP^1$; the example therefore
retains an extended base map while the total-space comparison
fails.

Finally take the product with $\PP^1$.  Its central integral
Betti numbers are
\[
 (1,0,1,3,2,3,2),
\]
and $h^1(\mathcal O)$ remains one.  The nearby fibre is
$H_-\times\PP^1$, of fundamental group $\Z$ and Euler
number zero.  These explicit invariants prevent this example
from being mistaken for a filling of $S^6$.
\end{proof}

\begin{remark}[Relation to the primary literature and to the remaining cusp]
The Cartier-divisor cocycle above is established machinery, not
a new logarithmic-transformation construction.  Fujimoto's
Definition~1.1 permits an arbitrary Cartier divisor; his
Remark~1.4 attributes non-K\"ahler torus degenerations to Ueno.
His Proposition~1.14 extends principal elliptic bundles across
subsets of codimension at least three.  Example~1.15 concerns
failure of extension as a principal elliptic bundle in
codimension two and explicitly leaves general relative
compactification open there; it supplies no unrestricted
nonexistence theorem for the present cusp
\citep[pp.~85--86, 93--94]{Fujimoto1990AB}.

Hopf components in analytic semistable degenerations are already
present in the original work of
\citet[Theorem 1.4 and Sections 2--3]{Nishiguchi1988IIRoot}.
The explicit diagonal Hopf quotients, elliptic conductor gluing
and normal-line calculation in
\citet[Section 3, especially Propositions 3.1--3.4]{SS00}
give a closely related analytic model.  Those are comparisons,
not smoothing theorems for the present threefold.  The construction
above supplies its own complete proper family and does not invoke
such a theorem by analogy.

There are two concrete obstacles to directly inserting the product
cap into B.3.  The divisor $E\times\PP^1$ in
$H_i\times\PP^1$ has normal degree zero along the $\PP^1$
ruling.  The exceptional divisor obtained by blowing up an
elliptic curve in a horizontal normalization has normal degree
$-1$ on its exceptional $\PP^1$ fibres.  For a two-component
attachment with no triple divisor on this interface, these cannot
be dual normal lines.  In the general normal-crossings formula
a triple divisor must also be included: a section
$E\times\{p\}$ has degree $+1$ on the ruling and could cancel
this particular mismatch.  No obstruction to such a higher-intersection
attachment is asserted.  Moreover the actual generic two-tori
do not possess the elliptic fibration of a principal elliptic
bundle, as the following exact calculation shows.

Fix an ordinary base point for which
$\mu\notin\Q\tau+\Q$.  Such points exist: otherwise the
ordinary base would be the countable union of the analytic loci
where one rational relation $\mu=r\tau+s$ holds; Baire's
theorem would force one such relation identically, contradicting
that $P$ is a nontorsion section.  For integral vectors
$v=(a,b,c,d)$ and $v'=(A,B,C,D)$, write
\[
 \Pi v=(6a\mu+b\tau+c,\ a\beta+b\mu+d).
\]
The coefficient of $\beta$ in
$\det(\Pi v,\Pi v')$ is
\[
 (Ab-aB)\tau+(Ac-aC).
\]
If it vanishes and one of $a,A$ is nonzero, the first three
coordinates of $v,v'$ are rationally proportional.  The
remaining determinant is
$(6a\mu+b\tau+c)(D-(A/a)d)$ after choosing $a\ne0$;
the first factor is nonzero, so a zero determinant would make
the two vectors dependent.  The case $a=0,A\ne0$ is the
same statement with the vectors interchanged.  If $a=A=0$,
the determinant is
\[
 (bD-Bd)\tau+(cB-Cb)\mu+(cD-Cd),
\]
and the rational independence of $1,\tau,\mu$ again forces
all the minors to vanish.  Thus every independent integral
pair spans a complex line for at most one value of $\beta$.
There are countably many pairs.  For all other admissible
parameters the two-torus is simple and contains no elliptic
subtorus.  In particular the distinguished line $\C\delta$
meets the period lattice in $\Z\delta$, so its complexified
orbit is $\C^*$, rather than an elliptic curve.

A principal elliptic-bundle construction may therefore describe
a central cap, but cannot by itself equal the entire prescribed
punctured cusp family.  It would need an additional affine
identification that destroys its generic elliptic fibration,
together with the actual nodal $q$-interface and the finite-place
component rotations.  These comparisons remain open.
\end{remark}

\begin{remark}[Limits of automatic analytic comparison]
The abstract identity $\Aut(X)=\C^*$ does not supply a
meromorphic compactification of the action. On an ordinary
two-torus the distinguished orbit is a nonconstant holomorphic
map from $\C^*$. If it extended meromorphically across both
ends to $\PP^1$, it would extend holomorphically: a
meromorphic map from a smooth curve to a compact complex
manifold has a removable indeterminacy. Every holomorphic
map $\PP^1\to T^2$ is constant, since it lifts to $\C^2$.
This is a contradiction. Thus an algebraic structure on the
abstract automorphism group is insufficient here.

The compactification theorem of
\citet[Theorem 1.2]{BiliottiGhigi2021Root} assumes Fujiki
class $\mathcal C$; that hypothesis is unavailable.
The local isomorphism theorem of
\citet[Theorem 3]{Meersseman2011Root} assumes pointwise
isomorphism on the whole base, including its central point,
and constant vector-field dimension. It cannot identify the
unknown central fibre. The Hopf compactification example of
\citet[Example 1]{KatzarkovEtAl2025Root} likewise illustrates
that an analytic open part need not determine a bimeromorphic
compactification. Their definability-based comparison is not
used: no definability of our punctured comparison has been
proved. B.4.7 supplies a proper family exhibiting the relevant
failure directly. None of these comparisons closes B.4.

Non-K\"ahler modifications of elliptic degenerations are
established. The product case in
\citet[Section 3, Example 3.4(1)]{Fujimoto1988Root}, using
the preceding construction, already gives such degenerations
with elliptic abelian nearby surfaces. Its elliptic-bundle
structure prevents a direct identification with our generic
simple two-tori. For the related codimension-two extension
question, \citet[Theorem 7.1.3]{Nakayama2002Root} proves an
extension criterion for locally bimeromorphically projective
elliptic fibrations over a punctured polydisc: the fibration
must be bimeromorphically projective over that punctured base.
It rules out the sectionless elliptic example discussed there.
Our regular fibres have dimension two and no asserted elliptic
quotient, so that theorem does not establish a nonexistence
result for the present filling.
\end{remark}

\begin{proposition}[B.4.8 (Proved): a normal-line obstruction on every semistable good-reduction interface]
Let $a,b$ be holomorphic units on a disc $\Delta_q$, with
$0<|a(q)|<1$, and let $m\ne0$ be an integer.  Put
\[
 \mathcal E_a=(\C^*\times\Delta_q)/
       \bigl((z,q)\sim(a(q)z,q)\bigr).
\]
This is a smooth proper elliptic family over the full disc.
Over $\Delta_q^*$ define a line bundle $\mathcal L_m$ by
\[
 (z,v,q)\sim\bigl(a(q)z,\ q^m b(q)v,\ q\bigr).
\]
Its restriction to each elliptic fibre has degree zero.
After any finite base change $q=s^N\varepsilon(s)$, where
$N>0$ and $\varepsilon$ is a holomorphic unit, the following
extension is impossible: a proper semistable family
$p:\mathcal C\to\Delta_s$ of connected curves, with smooth
total surface, an isomorphism
$\mathcal C|_{\Delta_s^*}\simeq\mathcal E_a|_{\Delta_s^*}$,
and a line bundle on $\mathcal C$ restricting to
$\mathcal L_m$ under that isomorphism.

The assertion allows any finite collection of rational tails
in the central curve and arbitrary translations in the
punctured isomorphism.  It remains true after tensoring
$\mathcal L_m$ by the restriction of a line bundle extending
on the proposed interface.  It does not assert an obstruction
for arbitrary singular, non-flat, or higher-dimensional
interface networks.

In particular the axis $z_1=0$ in the diagonal Hopf surfaces
\[
 H_{q,a(q)}=(\C^2\setminus\{0\})/
       \langle(z_1,z_2)\mapsto(qz_1,a(q)z_2)\rangle
\]
cannot acquire such a semistable completion together with
its normal line, after any finite base change.
\end{proposition}
\begin{proof}
Choose a logarithm of $a$ on the disc and set
$\tau_a=\log a/(2\pi i)$.  Then $\Im\tau_a>0$, and
\[
 \mathcal E_a\simeq
 (\C\times\Delta_q)/(\Z+\tau_a(q)\Z).
\]
The real coordinates $z=\exp(2\pi i(x+\tau_a(q)y))$
give a smooth trivialization by the marked torus
$(x,y)\in\R^2/\Z^2$.  In particular its monodromy on
$H_1$ is the identity.  The line factors in this marking
are $1$ along the $x$ cycle and $q^m b(q)$ along the
$y$ cycle.  On a positive parameter circle, the latter
factor has winding $m$.  Equivalently, the clutching
calculation gives
\[
 \langle c_1(\mathcal L_m),[T_x]\rangle=0,
 \qquad
 \langle c_1(\mathcal L_m),[T_y]\rangle=m,
 \qquad
 \langle c_1(\mathcal L_m),[\mathcal E_{a,q}]\rangle=0.
\]
Here $T_x,T_y$ sweep the indicated cycles around the
parameter circle; the orientation of $T_y$ is chosen to
agree with the clutching convention in the displayed
line factor.  Reversing that convention changes both
sides' sign and has no effect on nonvanishing.  One can
also compute the integer on the square
$[0,1]_y\times S^1_q$: the two $y$ edges are glued by
$q^m b(q)$, whose degree is $m$.

Under $q=s^N\varepsilon(s)$ the degree becomes $mN$.
The unit $\varepsilon(s)^m b(q(s))$ has zero winding.
Thus the mixed Chern class is the nonzero covector
\[
  (0,mN)\in
 H^1(\mathcal E_{a,q};\Z)\otimes H^1(S_s^1;\Z).
\]
This is an integral calculation; no limiting logarithm
or asymptotic estimate is involved.

Suppose an extension $\mathcal C$ existed.  Its central
fibre is connected.  Flatness and properness give arithmetic
genus one, and the normalization exact sequence gives
\[
 1=\sum_i g(C_i)+b_1(\Gamma),
\]
where $C_i$ are the smooth normalization components and
$\Gamma$ is the dual graph, including a loop for a
self-node.  Consequently either there is a unique
elliptic component and $\Gamma$ is a tree, or all
components are rational and $\Gamma$ has one cycle.

The second case has nontrivial monodromy.  This can be
seen directly in the local plumbing $uv=s$.  Transport
around $s=0$ is a positive Dehn twist on its neck annulus.
An edge outside the unique graph cycle gives a separating
neck and acts trivially on $H_1$.  Each of the $k\ge1$
edges in the graph cycle gives a nonseparating neck.
These disjoint necks represent the same primitive class
up to sign on the smoothed genus-one surface.  The
transvection is unchanged by reversing that class.
Their product therefore has, in a symplectic basis, the
matrix
\[
  \begin{pmatrix}1&k\\0&1\end{pmatrix}
\]
or its inverse under the opposite monodromy convention.
It is not the identity.  The given punctured elliptic
family has identity monodromy, so this case is impossible.
This elementary plumbing argument is all that is needed
here; no algebraic stable-reduction or comparison theorem
for arbitrary analytic total spaces is being assumed.

There is therefore an elliptic core $C_0$ with rational
trees attached at finitely many points.  Choose two
oriented loops on $C_0$ avoiding those points and with
intersection number one.  The map $p$ is a submersion
near their union.  Compactness of the loops and a smooth
horizontal distribution provide, after shrinking the
disc, families of loops over the full disc.  Each swept
torus on its boundary bounds the resulting solid torus
in $\mathcal C$.  Hence the first Chern class of every
line bundle on $\mathcal C$ evaluates to zero on both
such swept tori.

The two transported loops form a basis of the nearby
elliptic fibre's $H_1$: cutting off the rational trees
adds only genus-zero pieces along separating necks,
and their intersection number remains one.  Under the
given punctured isomorphism, their swept tori therefore
project to a basis of the mixed group
$H_1(\mathcal E_{a,q};\Z)\otimes H_1(S_s^1;\Z)$.
They may differ from the marked swept tori by integral
fibre classes.  This accounts in particular for an
arbitrary loop of fibre translations in the isomorphism.
The fibre degree of $\mathcal L_m$ is zero, so these
additional fibre classes do not change either pairing.
The nonzero covector $(0,mN)$ cannot vanish on both
basis vectors.  This contradicts the bounding-solid-torus
calculation and proves the assertion.

Finally an extending tensor factor has zero pairing on
both bounding tori by the same argument.  It cannot
cancel the obstruction.  Vertical divisors and any
Cartier triple divisor on such an extended interface
are included among these harmless tensor factors.

For the stated Hopf consequence, the axis is precisely
$\mathcal E_{a,q}$.
The derivative of the deck transformation in the $z_1$
direction multiplies a normal vector by $q$; its action
on the axis coordinate is multiplication by $a(q)$.
These are exactly the two factors defining
$\mathcal L_1$.  This is also the explicit diagonal-Hopf
normal-line convention of
\citet[Section 2.2]{AbashevaDeev2024AB}.
The assertion just proved applies with $m=1$ and $b=1$.
It also applies to the first multiplier $q^m b(q)$ for
any positive $m$.  At the actual cusp the nonzero unit
$e^{2\pi i\mu(q)}$ has absolute value different from one
at $q=0$.  Indeed the nodal cubic is
$Y^2=(X-1)^2(X+2)$; its normalization
$X=h^2-2$, $Y=h(h^2-3)$ and the multiplicative coordinate
$(h-\sqrt3)/(h+\sqrt3)$ show that the marked section at
$(X,Y)=(0,\pm\sqrt2)$ has multiplier one of
$-(5\pm2\sqrt6)$.
Choose that multiplier or its inverse to obtain
such a contracting $a(q)$.
This explains why the candidate axis has good reduction.
No claim is made that an arbitrary filling is obliged
to contain this particular axis or normal factor.
\end{proof}

\begin{remark}[What a corrected cap must change]
The obstruction concerns an individual normal line,
not just the tensor product of two normal lines in
the $T^1$ condition.  Opposite windings can cancel in
that tensor product while neither line extends on its
own proposed smooth branch.  Thus cancelling the
scalar $q$ by $q^{-1}$ is insufficient for this
good-reduction interface.  A successful correction
would have to change the interface or its punctured
normal bundle outside the class specified above.
The constant-elliptic example B.4.7 has no such
winding and remains consistent with B.4.8.
\end{remark}

\begin{proposition}[B.4.9 (Proved): ordinary logarithmic torus cocycles and the remaining descent classes]
Let $\mathcal T\to U$, $U=\Delta_q^*\times\Delta_r^*$,
be the Jacobian two-torus family of the marked cusp,
after $w=r^6$.  A holomorphic homomorphism of group
families
\[
 \C^*\times U\longrightarrow\mathcal T
\]
is determined by a flat integral cycle.  For this family
all such cycles are
\[
 c=a\widehat w+b\widehat\delta,\qquad a,b\in\Z.
\]
In the exponential cusp coordinates its homomorphism
is $z\mapsto(z^a,z^b)$.  Thus the ordinary Cartier
logarithm construction along a globally specified
cocharacter has transition maps
\[
 (Z,W)\longmapsto(g_{ij}^{\,a}Z,g_{ij}^{\,b}W).
\]
If these units extend over the prescribed fan charts,
the maps extend there and preserve each toric
boundary component.  Changing only these transition
units cannot repair the component permutations
$+4,-3,0$ refuted in B.3.1.

This is not a classification of all torus-torsor
descent.  In fact, for the sheaf $\mathscr T$ of
holomorphic sections of $\mathcal T$,
\[
 H^1(U,\mathscr T)\simeq H^2(U,\Lambda)\simeq\Z^2.
\]
The last group is represented by the coinvariant
classes of $\widehat\gamma,\widehat u$.  The ordinary
fixed-cocharacter Cartier constructions have zero
image in this group.  The actual cusp family with
the coherent closing translations of C.3.5 also
has zero torsor class, by the explicit holomorphic
section below.  Thus nonzero classes in this
$\Z^2$ do not have its prescribed marked punctured
cusp germ.  No compactification of an arbitrary
torsor is asserted.
\end{proposition}
\begin{proof}
First consider any smooth holomorphic family of
complex tori with integral period local system
$\Lambda$ and relative Lie bundle $\mathcal V$.
Lift a group homomorphism from $\C^*$ along the
universal covering $z=\exp(2\pi i t)$.
The lifted homomorphism from the additive group
$\C$ to the universal covering vector space of
each torus is $t\mapsto t v$ for a holomorphic
section $v$ of $\mathcal V$.  Periodicity under
$t\mapsto t+1$ says precisely that $v$ belongs
to the period lattice in every fibre.  In any
local marking its integral coefficients are
continuous, and hence locally constant.  It
therefore defines a flat global section of
$\Lambda$.  Conversely any such section gives
the homomorphism by this formula.  The image
need not be a closed elliptic subgroup; this
argument applies without a polarization and
also to simple two-tori.

In the ordered basis
$(\widehat\gamma,\widehat u,\widehat w,\widehat\delta)$,
the actual commuting linear monodromies have
nilpotent parts
\[
 N_q=A_0-I=
 \begin{pmatrix}
 0&0&0&0\\0&0&0&0\\0&1&0&0\\-1&0&0&0
 \end{pmatrix},\qquad
 N_r=-12\widehat\delta\otimes\gamma.
\]
The sign of $N_r$ is the positive-loop closing
convention of C.3.5; reversing it leaves all
groups below unchanged.  Their common integral
kernel is exactly $\Z\widehat w\oplus\Z\widehat\delta$.
The two corresponding columns of $\Pi_c$ are
$(1,0)^t,(0,1)^t$.  Exponentiating the homomorphism
therefore gives $(z^a,z^b)$, as asserted.

Fujimoto's logarithmic cocycle formula
\citep[Definition 1.1]{Fujimoto1990AB} works in
the same way for a complex torus of any dimension:
apply the displayed homomorphism to the transition
functions $g_{ij}$ of a line bundle.  Equivalently
use $c\log(g_{ij})/(2\pi i)$ modulo the period
lattice.  A different logarithm adds an integral
cycle, and the triple-overlap defect is an
integral cycle as well.  With the present common
invariant cycles this is exactly multiplication
by $(g_{ij}^a,g_{ij}^b)$.  Torus multiplication
extends to every toric variety of a fixed fan
and preserves every torus orbit and boundary
ray.  For an extending unit it is an automorphism
on the full chart.  Consequently it contributes
no permutation of the central component labels.
The product-of-permutations obstruction of B.3.1
is unchanged.  This conclusion keeps the fan,
period action, and local component models fixed.

To exhibit the scope left outside that calculation,
retract $U$ to its real two-torus.  The cellular
cochain complex with coefficients in the local
system is the two-generator Koszul complex.
Its last differential has image
$N_q\Lambda+N_r\Lambda$, and hence
\[
 H^2(U,\Lambda)
   =\Lambda/(N_q\Lambda+N_r\Lambda)
   =\Lambda/(\Z\widehat w+\Z\widehat\delta)
   \simeq\Z\widehat\gamma\oplus\Z\widehat u.
\]
The relation lattice is primitive: $N_q$ already
has image $\Z\widehat w+\Z\widehat\delta$.
The exact matrix and Smith calculation is saved
in \path{code/ab_push_hopf_winding.py}.

The analytic exponential sequence is
\[
 0\longrightarrow\Lambda\longrightarrow
 \mathcal O(\mathcal V)\longrightarrow
 \mathscr T\longrightarrow0.
\]
The product of two punctured discs is Stein,
and $\mathcal O(\mathcal V)$ is coherent.
Cartan's theorem B therefore kills its positive
cohomology
\citep[Section 5]{Cartan1952B_AB}.
The long exact sequence gives the stated
isomorphism $H^1(U,\mathscr T)\simeq H^2(U,\Lambda)$.
These classes do define holomorphic torsors by
translation gluing, but this calculation contains
no compactness assertion at either missing axis.

For the last assertion, naturality of the two
exponential sequences sends the first Chern class
of a Cartier line bundle to its multiple of the
flat cycle $c$ in $H^2(U,\Lambda)$.  Both common
invariant cycles lie in $N_q\Lambda$, so this
multiple is zero.  A holomorphic trivializing
gauge over $U$ need not extend at the missing
axes.  The proof thus neither identifies
arbitrary descent data with monomial cocycles
nor asserts a meromorphic comparison of
arbitrary analytic fillings.

We finally compute the actual affine class rather
than infer it from these linear monodromies.  Retain
the cusp functions of C.3.5,
\[
 q=e^{2\pi i s},\qquad
 C_0(q)=\beta^{\rm part}+s,\qquad
 \ell=-s/12+\kappa_0(q),
 \qquad D(q)=C_0(q)-12\kappa_0(q).
\]
The functions $\mu,C_0,\kappa_0$ are single-valued
and holomorphic at $q=0$.  The actual cusp unit
representative of the even-shift translation is
\[
 \bar b_{h,c}(q)=
 \left(-\frac h2\mu(q),\,
       -\frac h{12}(c+D(q))-\frac{h^2}{24}\right).
\]
On the logarithmic parameter cover define
\[
 \sigma(q,c)=
 \left(-\frac c2\mu(q),\,
       -\frac{c^2}{24}-\frac{cD(q)}{12}\right).
\]
Direct subtraction gives the exact identity
\[
 \sigma(q,c+h)-\sigma(q,c)=\bar b_{h,c}(q).
\]
In particular it is equivariant for the deck
transformation $c\mapsto c-12$ corresponding to
a positive $r$ loop, and it is single-valued in
$q$.  The original cusp has no extra affine
$q$ translation ($l_0=0$).  Thus $\sigma$ descends
to a holomorphic section of the actual affine
torus family.  Equivalently its $q$ and $r$
affine lifts commute exactly: the $r$ translation
$\bar b_{-12,c}$ is independent of the logarithm
branch $s$.  The resulting integral commutator
representing the Chern class is zero already in
$\Lambda$, before passage to coinvariants.

Multiplication by the exponential of $\sigma$
is a torus-unit automorphism of every cusp fan
chart for a fixed nonzero parameter.  Since its
coefficients are holomorphic in the full $q$
disc, this comparison includes the prescribed
Mumford cusp fibre for $r\ne0$.  The quadratic
dependence on $c=-(12/(2\pi i))\log r$ is retained:
it supplies no extension across $r=0$ and no
global comparison with the finite-place models.
Changing coherent closings by a scalar depending
only on the parameter has zero $q$ commutator
and leaves this Chern-class conclusion unchanged.
The explicit polynomial identity is checked in
\path{code/ab_push_hopf_winding.py}.
\end{proof}

\begin{proposition}[B.4.10 (Proved): an explicit parabolic family and its nonextending exceptional correction]
Fix an integer $n\ge1$ and holomorphic units $a(s),\kappa(s)$ on a
disc, with $0<|a(s)|<1$. After shrinking the disc there is a proper
holomorphic submersion $\mathcal K\to\Delta_s$, with an embedded smooth
proper elliptic family
\[
 \mathcal E=(\C^*\times\Delta_s)/
       ((v,s)\sim(a(s)v,s)).
\]
Its central surface is a parabolic Inoue surface $S_n$, with
$\pi_1(S_n)=\Z$, $b_2(S_n)=n$, and elliptic axis of square $-n$.
For $s\ne0$, the surface $\mathcal K_s$ is obtained from the diagonal
Hopf surface $H_{\kappa(s)s^n,a(s)}$ by $n$ successive blow-ups at one
axis point and its indicated infinitely near strict-axis points.
In a fixed annulus marking of $\mathcal E$, its normal line has factor
\[
 g_s(v)=\kappa(s)(a(s)v+s)^n.
\]
This line extends across $s=0$ and is isomorphic to the holomorphic
family with factor $\kappa(s)a(s)^n v^n$.
For constant $a,\kappa$ it is a constant family of degree $-n$ lines.

The relation with the normal line of the unmodified Hopf axis is
\[
 N_{\mathcal E_s/\mathcal K_s}
   =L_{\kappa(s)s^n}\otimes\OO_{\mathcal E_s}(-n[-s]),
       \qquad s\ne0.
\]
Here $n[-s]$ is the weighted total exceptional correction, not the
intersection with the reduced final exceptional locus. It has no
Cartier extension on $\mathcal E$, even after a finite base change or
the addition of a vertical component. Thus this construction does
not contradict B.4.8: the marked divisor used to change the normal
line fails to extend.

For $n=8$, the degree $-8$ axis normal can be chosen to be the inverse
of $\OO_E(2)$ on any chosen smooth elliptic complete intersection
$E=(2,2)\subset\PP^3$. This supplies the normal-line matching with
$\operatorname{Bl}_E\PP^3$ described below, including one positive
triple section. It is an auxiliary construction and first-order
matching result. No attachment to the prescribed punctured
$S^6$ family, compact third component, or global smoothing is asserted.
\end{proposition}
\begin{proof}
We give the gluing domains as well as the local formulas. Let
$B_R\subset\C^2$ be a ball, and let
$\pi_n:\widehat B_R\to B_R$ be the following modification.
Blow up the origin, then successively the intersection of the newest
exceptional curve with the strict transform of the axis $x=0$.
In the last chart,
\[
 \pi_n(z,w)=(zw^n,w).
\]
Define the parameter-dependent embedding in that chart by
\[
 \sigma_s(z,w)=
       \bigl(\kappa(s)z,\ a(s)w+(1-a(s))s\bigr).
\]
Choose constants $A<\rho<1$ and $K>0$ with
$|a(s)|\le A$ and $|\kappa(s)|\le K$ on a closed smaller disc.
Shrink it further so that
$|a(s)w+(1-a(s))s|\le\rho R$ on $\overline B_R$.
Then
\[
 |\pi_n\sigma_s(z,w)|^2
 \le K^2R^2(\rho R)^{2n}+\rho^2R^2<R^2
\]
when $R$ is small enough. The last-chart coordinates of these images
are bounded, so this is compact containment of
$\sigma_s(\overline B_R)$ in $\widehat B_R$, uniformly in $s$.
Notice that no special smallness condition on the Picard parameter
$\kappa$ is required.

Remove the embedded ball and identify the outer and inner collars by
$\sigma_s\pi_n$. The outer collar avoids the exceptional locus, and
$\pi_n$ is a biholomorphism there; thus the collar identification is
a biholomorphism. Carry out the same construction over the parameter
disc. The embeddings and the identifications are holomorphic and
preserve $s$. Uniformly disjoint collars give a Hausdorff smooth
threefold, and the projection is a holomorphic submersion in every
chart. Uniform compact fundamental regions over compact parameter
sets prove properness. This is the usual Kato construction, with
the explicit moving embedding above; see
\citet[Section 2, pp.1351--1353]{IstratiEtAl2022Root}.
The general occurrence of deformations from Kato surfaces to
blown-up Hopf surfaces is classical; see
\citet[Introduction, pp.1750--1751]{Teleman2010AB}.
We use explicit formulas here to retain the elliptic marking and
its exceptional correction.

The contracted map $\pi_n\sigma_s$ fixes $(0,s)$.
In the coordinate $v=w-s$ it becomes
\[
 f_s(x,v)=
    \bigl(\kappa(s)x(a(s)v+s)^n,\ a(s)v\bigr).
\]
At $s=0$ this is the parabolic monomial germ with exponent matrix
$\left(\begin{smallmatrix}1&n\\0&1\end{smallmatrix}\right)$.
The chosen modification is exactly the $n$-fold modification in
\citet[Example 9.10, p.1386]{IstratiEtAl2022Root}.
Consequently the central surface is the stated parabolic Inoue
surface; its elliptic axis is disjoint from the rational cycle.
The construction gives $\pi_1=\Z$ and $b_2=n$, or equivalently one
can use the cited description of this surface. These are invariants
of the auxiliary surface, not of a proposed $S^6$ central fibre.

The strict axis is preserved in the collar gluing. Its coordinate
$v$ is multiplied by $a(s)$. To see its full quotient, the axis disc
in the blowdown coordinates is
\[
 D_s=\{v:|v+s|<R\}.
\]
For small $s$, it contains $0$ and satisfies $a(s)D_s\Subset D_s$.
Its dilates have union $\C$ and shrinking intersection $\{0\}$.
Thus $D_s\setminus a(s)D_s$, with the two boundary collars glued,
is a fundamental annulus for $\C^*/a(s)^\Z$.
This proves the exact elliptic marking throughout the parameter disc.

There are two normal frames, whose distinction is essential.
In the last blow-up chart the collar map $\sigma_s\pi_n$ has factor
\[
 g_z(v)=\kappa(s)(v+s)^n.
\]
On a fixed annulus avoiding $v=0$, shrink $s$ so that both
$v+s$ and $a(s)v+s$ are nonzero. The blowdown normal coordinate is
$x=z(v+s)^n$, and in that frame the factor is
\[
 g_x(v)=\kappa(s)(a(s)v+s)^n.
\]
The change between the frames is a nonzero gauge on these annuli.
The central factor is $\kappa(0)a(0)^n v^n$ and has degree $-n$,
with the factor-of-automorphy convention of
\citet[Proposition 5.6, pp.1369--1370]{IstratiEtAl2022Root}.

An explicit regular gauge verifies the whole family. On a fixed
annulus and its overlap with its $a(s)$-translate put
\[
 H_s(v)=\prod_{j\ge0}\left(1+\frac{a(s)^j s}{v}\right)^n.
\]
The product converges locally uniformly in $(v,s)$ and is nonzero
after shrinking the parameter disc. It satisfies
\[
 \frac{H_s(a(s)v)}{H_s(v)}
      =\left(1+\frac{s}{a(s)v}\right)^n.
\]
Replacing $x$ by $x/H_s(v)$ changes $g_x$ to
$\kappa(s)a(s)^n v^n$.
The collar maps are linear in the normal coordinate, so the same
calculation identifies an actual tubular line neighbourhood, not
only its first infinitesimal normal bundle.

For completeness, the generic Hopf blowdown also has an explicit
developing map. Fix $s\ne0$ and set
\[
 G_s(v)=\prod_{j\ge1}\left(1+\frac{a(s)^j v}{s}\right)^n,
       \qquad
 \Psi_s(x,v)=(xG_s(v),v).
\]
This product is entire in $v$, converges locally uniformly, and obeys
\[
 G_s(a(s)v)=\frac{G_s(v)}{(1+a(s)v/s)^n},\qquad
 \Psi_s f_s=\operatorname{diag}(\kappa(s)s^n,a(s))\Psi_s.
\]
For sufficiently small nonzero $s$, both diagonal multipliers have
absolute value less than one. Compose $\Psi_s$ with $\pi_n$ on a
Kato fundamental region. The displayed equivariance glues it to a
holomorphic map to $H_{\kappa(s)s^n,a(s)}$.
It never reaches the excluded origin: that would require
$x=v=0$, which is the fixed point $(0,s)$ removed with the inner ball.
All zero-factor curves of $G_s$ have
$v=-s/a(s)^j$, so they lie in the one axis orbit $[-s]$.
The original modification also contracts over $v=-s$.
At that representative,
\[
 G_s(-s)=\prod_{j\ge1}(1-a(s)^j)^n\ne0.
\]
Hence locally the map is precisely $\pi_n$ followed by a coordinate
unit. Away from that orbit it is a local biholomorphism of degree
one. Compactness makes it a proper modification, and the displayed
local model identifies it as the stated $n$ successive blow-ups.
This is a punctured-parameter construction: no meromorphic extension
of $G_s$ through $s=0$ is being claimed.

Each blow-up of a point on the strict elliptic curve subtracts that
point from its normal line. Following all $n$ stages gives
\[
 N_s=L_{\kappa(s)s^n}\otimes\OO_{\mathcal E_s}(-n[-s]).
\]
The multiplicity $n$ records total exceptional transforms.
The reduced final exceptional locus meets the strict axis only once.
One may check the formula directly: the standard theta factor
$\theta_a(av)=-v^{-1}\theta_a(v)$ shows that $\OO([p])$
has factor $-p/v$. At $p=-s$ the tensor product above has factor
\[
 \kappa(s)s^n(v/s)^n=\kappa(s)v^n,
\]
which is equivalent to $\kappa(s)a(s)^n v^n$ by a nonzero monomial
gauge. This checks the signs and degrees independently of the
successive-blow-up formula.

The divisor $n[-s]$ cannot extend. For constant $a$, choose two
distinct classes $[v_1],[v_2]\in\C^*/a^\Z$.
The sequences $s_j=-a^jv_i$ tend to zero, whereas the corresponding
divisors are constantly $n[v_i]$ and have different limits in
$\operatorname{Sym}^n(E_a)$.
For varying $a(s)$, choose the holomorphic logarithm
$\tau_a(s)=\log a(s)/(2\pi i)$ on the disc.
The Abel--Jacobi coordinate of $n[-s]$ is
$n\log(-s)/(2\pi i)$ modulo $\Z+\tau_a(s)\Z$.
On a positive parameter loop it winds $n$ times around a primitive
cycle in the smooth elliptic family. An extending section over the
disc would have zero winding. A finite base change of order $N$
multiplies this nonzero integer by $N$.
Adding a vertical elliptic component does not help: remove that
Cartier component, and the remaining proper horizontal divisor of
degree $n$ would be finite over the disc and define a continuous
section of the relative symmetric power, which has just been ruled
out. Nonproper exceptional divisors in deformations of Kato surfaces
are a classical phenomenon, studied by
\citet[Section 1.2]{DlousskyTeleman2012AB}; the calculation here
identifies the precise weighted divisor in this marked family.

This also explains the scope of ordinary elementary modifications.
If a degree-zero line $L_m$ from B.4.8 is changed by point blow-ups,
then its new normal is $L_m\otimes\OO(-D_{\rm exc})$.
If both that normal and the weighted Cartier correction extended on
a semistable good-reduction interface, their tensor product would
extend $L_m$, contradicting B.4.8. Drifting translations cannot
remove both requirements. The present Kato family escapes exactly
because $D_{\rm exc}=n[-s]$ has no such extension.

Finally let $E=(2,2)\subset\PP^3$ be smooth and set
$L=\OO_E(2)$, of degree eight. Its normal in $\PP^3$ is
$L\oplus L$. For $A=\operatorname{Bl}_E\PP^3$, the exceptional
divisor is $D=E\times\PP^1$ and, in the lines convention,
\[
 N_{D/A}=L\boxtimes\OO_{\PP^1}(-1).
\]
A bare Hopf-product cap has normal $M\boxtimes\OO$ with
$\deg M=0$. The required triple-divisor line would then be
$(L\otimes M)^{-1}\boxtimes\OO(1)$, which has no nonzero section
because its elliptic degree is $-8$. Thus no effective triple divisor
repairs that particular bare interface.

For the parabolic surface with $n=8$, changing $\kappa$ varies the
flat factor through all of $\operatorname{Pic}^0(E)$:
constant factors are parametrized by $\C^*/a^\Z$.
Choose its axis normal to equal $L^{-1}$.
In $B=S_8\times\PP^1$ the divisor $D=E\times\PP^1$ then has normal
$L^{-1}\boxtimes\OO$. With triple section $T=E\times\{p\}$,
\[
 N_{D/A}\otimes N_{D/B}\otimes\OO_D(T)
   =(L\boxtimes\OO(-1))\otimes
     (L^{-1}\boxtimes\OO)\otimes(\OO\boxtimes\OO(1))
   \simeq\OO_D.
\]
This proves precisely the first-order matching claimed.
\end{proof}

\begin{remark}[The third component and the actual cusp remain open]
Choose a smooth quadric $Q\subset\PP^3$ containing $E$.
Its strict transform in $A$ is still $Q\simeq\PP^1\times\PP^1$:
blowing up its Cartier divisor $E$ does not change $Q$.
Its normal is
$\OO_Q(2,2)(-E)\simeq\OO_Q$, and it meets $D$ in a section.
The matching surface $S_8\times\{p\}$ in $B$ also has trivial
normal. A third threefold would therefore need divisors $Q$ and
$S_8$ meeting along $E$, with normals
$\OO_Q(-E)$ and $\OO_{S_8}(-E)$ respectively.
These are exactly the normal conditions for the central surface
$Q\cup_E S_8$ in a semistable surface smoothing.
If such a smoothing exists, its topological Euler number is
$4+8-2\chi(E)=12$ and its holomorphic Euler number is one.
These necessary data do not supply a compact third threefold or
prove the existence of that smoothing.

A full assembly must also preserve the required coherent cohomology
from B.4.5; a lone parabolic component attached to rational components
can lose its $H^1(\OO)$ through the elliptic overlap.
Additional conductor surfaces meeting $E$ would contribute further
triple divisors and must be included in the normal equation.
Nor is a local change of the cusp parameter automatically a permitted
global base change of the prescribed family.
Thus B.4.1 and B.5.1 remain open after this auxiliary construction.
\end{remark}

% Owner: boundary_d.  Independent proof; reserved by root as B.4.11.
% This is an auxiliary contracted surface, not the requested S^6 filling.
\begin{proposition}[B.4.11 (Proved): globalization at the contracted parabolic elliptic axis]
Let $S=S_8(a,\kappa)$ be the central parabolic surface in B.4.10,
let $E\subset S$ be its elliptic axis, and put
$N=N_{E/S}$, so $\deg N=-8$.
The contraction $\pi:S\to Y$ of $E$ is a compact normal surface
with a single isolated singular point $p$.
Its germ at $p$ is the vertex of the affine cone over
$(E,N^{-1})$.  In particular,
\[
 H^2(Y,\Theta_Y)=0.
\]
The restriction morphism from a semiuniversal deformation of $Y$
to a semiuniversal deformation of $(Y,p)$ is regular.
Consequently every prescribed analytic one-parameter deformation
of this isolated singularity is induced, after shrinking its
parameter disc, by a proper deformation of the compact surface $Y$.
If the prescribed local deformation smooths the singularity, the
global deformation has smooth fibres away from the central value.

This concerns globalization of a specified local deformation of
an auxiliary compact surface.  It does not assert an attachment to
the punctured $S^6$ family or a compactification of the resulting
threefold over the open parameter disc.
\end{proposition}
\begin{proof}
We first retain the analytic neighbourhood, not merely its normal
line.  At the central parameter the Kato gluing of B.4.10 is
\[
 (x,v)\longmapsto(\kappa a^8xv^8,av).
\]
It is exactly linear in $x$.  Thus a neighbourhood of $E$ is a
neighbourhood of the zero section in the negative line bundle
$N$, with the indicated factor of automorphy.
Its contraction is the cone associated with the section ring
\[
 R(E,N^{-1})=\bigoplus_{m\ge0}H^0(E,N^{-m}).
\]
A line bundle of degree eight on an elliptic curve is very ample
and normally generated.  The cone is normal, its only singularity
is its vertex, and the blow-up of the vertex maximal ideal is
$\operatorname{Tot}(N)$ near its zero section.
One can see the last assertion directly: the degree-one sections
generate the section ring and have no common zero on $E$;
their ratios on the blow-up recover the projective embedding of
$E$, and their common radial coordinate is the coordinate in $N$.
In particular, with $I=\mathcal O_S(-E)$ and $\mathfrak m_p$ the
vertex maximal ideal,
\[
 \mathfrak m_p\mathcal O_S=I.
\]
The local contraction glues to the identity on $S\setminus E$,
giving the stated proper contraction of the compact surface.

Write $C$ for the disjoint rational cycle of the parabolic surface.
The meromorphic two-form
$d\log x\wedge d\log v$ is invariant under the displayed
monomial gluing.  In the toric blow-up charts it has simple poles
along precisely the axis and the rational cycle, and no zeros.
Consequently
\[
 K_S=\mathcal O_S(-E-C).
\]
Also $b_1(S)=1$, by B.4.10.
Put $\mathcal F=\Theta_S(-\log E)$, the locally free sheaf of
vector fields tangent to $E$.
Serre duality and the displayed canonical divisor give
\[
 H^2(S,\mathcal F)^*
 =H^0\bigl(S,\Omega^1_S(\log E)\otimes K_S\bigr)=0.
\]
Indeed there are sheaf injections
\[
 \Omega^1_S(\log E)(-E-C)
       \lhook\joinrel\longrightarrow\Omega^1_S(-C)
       \lhook\joinrel\longrightarrow\Omega^1_S.
\]
Locally the first sends the logarithmic generators
$dx/x,dy$, multiplied by $x$, to the ordinary forms $dx,xdy$.
Finally $H^0(S,\Omega^1_S)=0$ follows already from $b_1(S)=1$.
For completeness, a holomorphic one-form $\alpha$ on a compact
complex surface is closed: Stokes' theorem applied to
$\alpha\wedge\overline{d\alpha}$ makes the integral of the
positive volume form associated with
$d\alpha\wedge\overline{d\alpha}$ vanish.
If $\alpha\ne0$, its real and imaginary parts give independent
real de Rham classes.  An exact real linear combination would
be the differential of a global real pluriharmonic function,
which is constant on the compact surface.  This would give
$b_1(S)\ge2$, a contradiction.

We compute the one possible local higher direct image explicitly.
The restriction of $\mathcal F$ to $E$ is the Atiyah extension
of $N$:
\[
 0\longrightarrow\mathcal O_E
 \longrightarrow\mathcal A_N:=\mathcal F|_E
 \longrightarrow\Theta_E\longrightarrow0.
\]
In normal coordinates related by $x_j=g_{ij}(v)x_i$, the
vertical Euler field is $x_i\partial_{x_i}$, and the changes of
horizontal lifts have cocycle $d\log g_{ij}$.
Thus this extension has class $c_1(N)$ in
$H^1(E,\Omega^1_E)$, which is nonzero because $\deg N=-8$.
Since $\Theta_E\simeq\mathcal O_E$, the connecting map
$H^0(E,\Theta_E)\to H^1(E,\mathcal O_E)$ is an isomorphism.
The same cohomology sequence therefore gives the natural
isomorphism
\[
 H^1(E,\mathcal A_N)\xrightarrow{\ \sim\ }
 H^1(E,\Theta_E)\simeq\C.
\]

The successive quotients for the formal neighbourhood of $E$ are
\[
 I^m\mathcal F/I^{m+1}\mathcal F
     =\mathcal A_N\otimes N^{-m}\qquad(m\ge0).
\]
For $m\ge1$ both line bundles in the induced extension have
degree $8m>0$ on the elliptic curve, so their first cohomology
vanishes.  It follows that all restriction maps
\[
 H^1\bigl(E,\mathcal F/I^{m+1}\mathcal F\bigr)
 \longrightarrow H^1(E,\mathcal F/I^m\mathcal F)
 \quad(m\ge1)
\]
are isomorphisms.  Their common value is
$H^1(E,\mathcal A_N)\simeq\C$.
The analytic theorem on formal functions
\citep[Section~6, Hauptsatz~IIa, p.~59]{Grauert60push}
identifies this inverse limit with the completed stalk of
$R^1\pi_*\mathcal F$ at $p$.
This coherent sheaf is supported at the single point $p$,
so its stalk is a finite-length module and equals its completion.
We have proved
\[
 R^1\pi_*\mathcal F\simeq\C_p,
 \qquad
 H^0(Y,R^1\pi_*\mathcal F)
    \xrightarrow{\ \sim\ }H^1(E,\Theta_E).
\]
The second map is the natural restriction followed by the
Atiyah-extension map just computed, not a dimension count alone.

There is also an equality of sheaves
\[
 \pi_*\mathcal F=\Theta_Y.
\]
To check it at $p$, a holomorphic derivation of the cone
integrates locally to a holomorphic flow: embed the cone in a
polydisc, lift its coefficient functions, and use the ordinary
holomorphic vector-field flow, which preserves the defining ideal.
The flow preserves the unique singular point, hence its maximal
ideal.  Functoriality of the blow-up of that ideal lifts the flow
to the resolution, preserving the exceptional curve.
Its derivative is therefore a section of $\mathcal F$.
Conversely a section of $\mathcal F$ differentiates pullbacks
of holomorphic functions on $Y$ and descends to a derivation,
because $\pi_*\mathcal O_S=\mathcal O_Y$.
Both operations agree off $E$ and are inverse there; uniqueness
extends the equality over $E$.

The global-to-local cohomology map is surjective for a concrete
reason.  Write $a=e^{2\pi i\tau}$ with $\operatorname{Im}\tau>0$,
and vary $\tau$ to $\tau+u$ in the explicit Kato construction,
keeping its translation parameter zero.
The same uniform collar argument as in B.4.10 gives a proper
holomorphic deformation of the pair $(S,E)$.
Its Kodaira--Spencer class belongs to $H^1(S,\mathcal F)$.
Under the natural composite
\[
 H^1(S,\mathcal F)\longrightarrow
 H^0(Y,R^1\pi_*\mathcal F)
 \xrightarrow{\ \sim\ }H^1(E,\Theta_E)
\]
it is the Kodaira--Spencer class of
$\C/(\Z+(\tau+u)\Z)$, which is nonzero.
For example, differentiating its second deck transformation
$z\mapsto z+\tau+u$ gives the usual nonzero period cocycle.
This argument works also at the elliptic curves with extra
automorphisms; it does not differentiate the coarse $j$-invariant.
The target is one dimensional, so the first arrow is onto.

The degree-two part of the Leray sequence is
\[
 H^1(S,\mathcal F)\longrightarrow
 H^0(Y,R^1\pi_*\mathcal F)\longrightarrow
 H^2(Y,\pi_*\mathcal F)\longrightarrow H^2(S,\mathcal F).
\]
The first map is surjective and the last group is zero.
Using $\pi_*\mathcal F=\Theta_Y$ proves
$H^2(Y,\Theta_Y)=0$.

We spell out the local-to-global criterion used in the final
step, since the rational-cycle contraction theorem does not
apply to this elliptic contraction.
For a compact complex space with isolated singularities,
$H^2(\Theta)=0$ makes the restriction from global to local
deformations formally smooth.
Here is the obstruction argument for the present single
singularity.  Take a finite Stein cover with one small
neighbourhood of $p$ and all other members smooth.
For a small extension of Artinian bases, prescribe the lifted
deformation on the singular member and take smooth local
lifts on the other members.  Their pairwise identifications
lift on the smooth Stein intersections.  The failure of the
cocycle identity on triple intersections is a Cech two-cocycle
in $\Theta_Y$ tensored with the square-zero kernel.
Its class vanishes by the proved $H^2$-vanishing; correcting
the identifications by a one-cochain glues the desired lift.
This proves the relative infinitesimal lifting property.

The analytic semiuniversal deformation theory turns this
property into regularity of the induced morphism of analytic
base germs.  The general criterion is
\citet[Lemma~1, p.~93]{Manetti1991D}, stated for reduced compact
complex spaces with isolated singularities and proved by the
local-to-global spectral sequence of the cotangent complex.
It is also used in
\citet[Section~3.2, Proposition~3.3 and the discussion preceding
Proposition~3.4]{DlousskyTeleman2020D}.
There is no hypothesis about the kind of exceptional curve
in this criterion.  A regular morphism of analytic
germs is locally a projection from a product with a smooth
germ.  Consequently a specified analytic disc in the local
deformation base lifts to the global base, without replacing
the prescribed local deformation by a different smoothing
component.
The global family is proper by the compact deformation
construction.  Away from $p$ its central fibre is smooth,
and smoothness is open; compactness permits one common
shrinking of the disc.  A local smoothing at $p$ therefore
gives smooth global noncentral fibres, as asserted.
\end{proof}

\begin{proposition}[B.4.12 (Proved): a proper quadric--parabolic surface smoothing]
Let $S=S_8(a,\kappa)$ be as in B.4.11, with elliptic axis $E$
and normal line $N$ of degree $-8$.
There is an embedding $E\hookrightarrow Q=\PP^1\times\PP^1$
as a smooth divisor of bidegree $(2,2)$ with
$N_{E/Q}\simeq N^{-1}$.
There is a proper flat map from a smooth complex threefold
\[
 \varpi:\mathcal V\longrightarrow\Delta
\]
with smooth noncentral fibres and reduced simple normal-crossings
central fibre exactly $S\cup_E Q$.  Its two normal lines are
\[
 N_{S/\mathcal V}=\OO_S(-E),\qquad
 N_{Q/\mathcal V}=K_Q=\OO_Q(-2,-2).
\]
The family is obtained by globalizing a specified cone smoothing
and blowing up its vertex; no finite base change is needed.

For the elliptic-quartic configuration of B.4.10, the three
pairwise normal-line products, including the triple-curve factors,
can consequently all be made trivial using
\[
 A=\operatorname{Bl}_E\PP^3,\qquad
 B=S\times\PP^1,\qquad C=\mathcal V.
\]
Here $C$ is a proper family over an open disc and is not a compact
threefold.  This assertion supplies the indicated surface smoothing
and normal-line products only.  It does not construct a compact
three-component central fibre, a compatible global smoothing of
such a fibre, or a comparison with the prescribed punctured
$S^6$ family.
\end{proposition}
\begin{proof}
Put $L=N^{-1}$, of degree eight.  Multiplication by two on the
elliptic Picard variety is surjective, so choose a degree-four
line bundle $M$ with $M^{\otimes2}=L$.
Choose degree-two line bundles $M_1,M_2$ with
$M_1\otimes M_2=M$ and $M_1\not\simeq M_2$.
Each complete linear system is a basepoint-free double cover of
$\PP^1$.  The product of the two maps is an embedding in $Q$.
Indeed two distinct points identified by both maps would give
$M_1\simeq\OO_E(p+q)\simeq M_2$.
A common ramification point would similarly give
$M_1\simeq\OO_E(2p)\simeq M_2$.
Thus the product is injective and immersive; properness makes it a
closed embedding.  Its image has bidegree $(2,2)$, and
\[
 \OO_Q(E)|_E=(M_1\otimes M_2)^{\otimes2}=L.
\]
This also allows us to start with the fixed quadric and elliptic
quartic in B.4.10: the normal line there was chosen to be the
inverse of this same $L$.

Use the anticanonical embedding of $Q$ in $\PP^8$ and write
\[
 R_Q=\bigoplus_{d\ge0}H^0(Q,\OO_Q(2d,2d)),\qquad
 Z=\operatorname{Spec}R_Q.
\]
The ring is generated in degree one and is normal.
One direct verification uses the bihomogeneous monomials:
every monomial of bidegree $(2d,2d)$ is a product of $d$
monomials of bidegree $(2,2)$, and the associated monoid is the
intersection of a rational polyhedral cone with its full lattice.
The punctured cone is the complement of the zero section in
$\OO_Q(-2,-2)$, hence is smooth.
Choose the section $e\in H^0(Q,\OO_Q(2,2))$ cutting out $E$.
The degree-one function $t=e$ on $Z$ defines a flat map
$Z\to\mathbb A^1$: its coordinate ring is torsion-free over
$\C[t]$, a principal ideal domain.
For every $d\ge0$, the sequence
\[
 0\longrightarrow\OO_Q(2d-2,2d-2)
 \xrightarrow{\ e\ }\OO_Q(2d,2d)
 \longrightarrow L^{\otimes d}\longrightarrow0
\]
has vanishing first cohomology in its left term.
For $d=0$ this follows from the Kunneth formula for
$\OO_Q(-2,-2)$; for $d\ge1$ it follows from the same formula
and nonnegative degrees.
Taking sections identifies
\[
 R_Q/(e)=\bigoplus_{d\ge0}H^0(E,L^{\otimes d}).
\]
The central surface is therefore exactly the affine cone over
$(E,L)$, including its scheme structure.
For $t\ne0$, projection from the cone vertex identifies the
fibre with $Q\setminus E$: over each point where $e$ is nonzero
there is exactly one radial coordinate on which $e$ has value $t$.
In particular this is an explicit smoothing of the isolated
singularity of B.4.11.
It is the affine version of Pinkham's cone construction
\citep[Remark~7.6(iii), p.~46; Section~9, pp.~63--65]{Pinkham1974Root}.
In degree eight the quadric component is distinguished from the
other smoothing component; no replacement of it by that other
component is made here.

Blowing up the vertex of $Z$ gives
$\operatorname{Tot}(\OO_Q(-2,-2))=\operatorname{Tot}(K_Q)$.
This can be checked on the blow-up charts given by ratios of the
degree-one homogeneous coordinates; these recover a point of
$Q$, with the remaining common factor the tautological fibre
coordinate $\xi$.
The map to the parameter line is precisely
\[
 t=\xi e.
\]
Its central divisor is the union of the zero section $Q$ and
the restriction $\operatorname{Tot}(K_Q|_E)$, each with
multiplicity one.  Near $E$ the equation is $t=\xi z$, where
$z$ is a local defining function of $E$ in $Q$.
Thus the total space is smooth and the central divisor has
simple normal crossings.

Now contract $E$ in the compact surface $S$ as in B.4.11.
The regular global-to-local deformation map in that proposition
lifts the specified analytic germ $Z\to\mathbb A^1$ to a proper
deformation $\mathcal Y\to\Delta$ of $Y=S/E$.
The lifted deformation agrees analytically with this cone
smoothing in a neighbourhood of its singular point.
Away from that neighbourhood the central surface is smooth;
openness of smoothness and properness allow a common shrinking
of $\Delta$ for which the family there is a submersion.
Blow up the single cone vertex in $\mathcal Y$.
The resulting map $\mathcal V\to\Delta$ remains proper.
Near the exceptional divisor its local model is the smooth
space just computed; elsewhere it is unchanged and smooth.
Its central strict transform is exactly $S$, because the
maximal-ideal blow-up of the degree-eight elliptic cone is its
given resolution $\operatorname{Tot}(N)$ near $E$.
The other component is $Q$, and their intersection is the original
marked elliptic curve.  The local equation $t=\xi z$ proves
flatness and reducedness near the intersection; at a smooth
point of either component $t$ is a local coordinate.
The noncentral fibres are the smooth compact fibres of
$\mathcal Y$.
This proves every asserted properness and smoothness property
for the surface family.

The exceptional zero section has normal $K_Q$ in its local
total-space model.  Since $S+Q$ is the principal divisor of $t$,
restriction to $S$ gives
\[
 \OO_{\mathcal V}(S)|_S
   =\OO_{\mathcal V}(-Q)|_S=\OO_S(-E).
\]
These are the two asserted normal lines.  They also show directly
that the two branch normals along $E$ multiply to the trivial line.

We record the precise extent of the three-piece matching.
Realize $Q$ as a smooth quadric in $\PP^3$ and $E$ as its
intersection with a second quadric.  Put $L=\OO_E(2)$.
Then $N_{E/\PP^3}=L\oplus L$, so the exceptional divisor
of $A=\operatorname{Bl}_E\PP^3$ is
$D_{AB}=E\times\PP^1$, with normal
$L\boxtimes\OO_{\PP^1}(-1)$.
The strict transform of $Q$ is still $Q$, because the centre is
Cartier on it, and its normal in $A$ is
\[
 N_{Q/A}=\OO_Q(2)\otimes\OO_Q(-E)=\OO_Q.
\]
It meets $D_{AB}$ in the constant section $E\times\{p\}$
corresponding to that member of the pencil of quadrics.
In $B=S\times\PP^1$, use the divisors
$D_{BA}=E\times\PP^1$ and $D_{BC}=S\times\{p\}$.
Their normals are $L^{-1}\boxtimes\OO$ and $\OO_S$.
In $C=\mathcal V$, use its central divisors $Q$ and $S$.
The three normal products with the triple contribution are then
\[
\begin{array}{c|c}
 \text{double surface}&
 N_{D/\text{first branch}}\otimes
 N_{D/\text{second branch}}\otimes\OO_D(\text{triple curve})\\ \hline
 E\times\PP^1&
 (L\boxtimes\OO(-1))\otimes(L^{-1}\boxtimes\OO)
       \otimes(\OO_E\boxtimes\OO(1))=\OO\\
 Q&\OO_Q\otimes K_Q\otimes\OO_Q(E)=\OO_Q\\
 S&\OO_S\otimes\OO_S(-E)\otimes\OO_S(E)=\OO_S.
\end{array}
\]
These are pairwise line-bundle identities.  They do not supply
a compact third component or a compatible global smoothing
section for a compact normal-crossings threefold.
In fact the third space constructed here maps surjectively and
properly to an open disc, so it cannot itself be compact.

For a useful topological check on the local model,
the affine smooth fibre $Q\setminus E$ has
\[
 H_1(Q\setminus E;\Z)=\Z/2,\qquad
 H_2(Q\setminus E;\Z)=\Z^3.
\]
Indeed the Thom sequence of the pair $(Q,Q\setminus E)$ has
intersection map $H_2(Q;\Z)\to H_0(E;\Z)$ given by
$(a,b)\mapsto2a+2b$.  Its cokernel is $\Z/2$ and its
kernel is $\Z$; the preceding term $H_1(E;\Z)=\Z^2$
injects into $H_2(Q\setminus E;\Z)$.
This exact extension of free groups splits abstractly.
The map $H_4(Q;\Z)\to H_2(E;\Z)$ is an isomorphism,
so $H_3(Q\setminus E;\Z)=0$, and its Euler number is four.
The Euler number is also obtained independently by projection
to either ruling of $Q$: away from the four branch values of
$E\to\PP^1$, the fibre is $\C^*$ and has Euler number zero;
at each branch value it is $\C$ and has Euler number one.
In particular no simply connected local Milnor fibre is being
assumed.  The compact noncentral surface has Euler number
$\chi(S)+\chi(Q)-2\chi(E)=8+4=12$ by the local
$\xi z=t$ retraction, not the Euler number of an $S^6$ fibre.
Exact polynomial, Hilbert-function, normal-product and Smith-form
checks appear in \texttt{code/root\_quadric\_smoothing.py}.
\end{proof}

\begin{proposition}[B.4.13 (Proved): the second period does not lift to a finite marked Hopf modification]
Fix $0<|\alpha|,|a|<1$, and put
\[
 H=(\C^2\setminus\{0\})/
       \langle(x,v)\mapsto(\alpha x,av)\rangle,
       \qquad E=\{x=0\}\simeq\C^*/a^\Z.
\]
Let $h$ be the automorphism induced by
$(x,v)\mapsto(c x,bv)$, where $b,c\in\C^*$.
Suppose $p:Y\to H$ is a finite composition of point blow-ups,
allowing infinitely near centres, and at least one point of $E$
is an image of its exceptional locus. If $[b]\in\C^*/a^\Z$ has
infinite order, there is no holomorphic automorphism of $Y$
which agrees, on the dense unmodified core, with the marked
automorphism $h$.

The same obstruction holds for a finite collection of such marked
modifications if the proposed period permutes the collection and
acts as $h$ on each marked core.
More generally, if a compact smooth surface contains a preserved
marked elliptic curve $E$ of nonzero normal degree $d$, every
holomorphic automorphism restricting to translation by $b$ on $E$
must satisfy $d\,b=0$ in $E$.
The latter assertion also excludes a finite cycle of such preserved
marked cap axes when the cycle composite is an infinite-order
translation. These assertions apply to the direct attempt
to lift the second actual cusp period to the finite Hopf
modifications in B.4.10, for generic nonzero parameters, after every
finite base change. It concerns a holomorphic period action on
these full modifications. It does not exclude a meromorphic
correspondence, a different connected unfolding made from
noncompact pieces, or a cap present only over the central parameter.
\end{proposition}
\begin{proof}
Let $S\subset H$ be the image of the exceptional curves of $p$.
It is finite. Repeated infinitely near blow-ups change the curves
over a point but do not make $S$ infinite.
If $\widetilde h$ were a lift, the maps
\[
 p\widetilde h,\ hp:Y\longrightarrow H
\]
would agree on the dense core and hence everywhere, by the identity
principle for holomorphic maps to the separated complex surface $H$.
The image under $\widetilde h$ of a curve contracted by $p$ is
therefore again contracted by $p$. Applying the same argument to
$\widetilde h^{-1}$ proves that the total exceptional locus is
preserved. Thus $h(S)=S$.

The automorphism $h$ preserves $E$, where it is translation by
$[b]$. Consequently the nonempty finite set $S\cap E$ would be
invariant under that translation. The orbit of any one of its
points is finite, so for some $k>0$,
\[
 [b^k v]=[v],\qquad\text{equivalently } b^k\in a^\Z.
\]
This says exactly that $[b]$ is torsion, a contradiction.
In particular, no choice of multiplicities or infinitely near
centres changes the obstruction: only their nonempty finite
support on $E$ is needed.

For finitely many modifications $p_i:Y_i\to H$, a proposed deck
map has biholomorphisms
$Y_i\to Y_{\sigma(i)}$ satisfying the same core comparison.
The identity principle again gives equivariance everywhere.
The union of the finitely many exceptional image sets is
$h$-invariant; its intersection with $E$ gives the same contradiction.
Equivalently, around a cycle of length $\ell$ the map $h^\ell$
would lift to a single modification, while $[b^\ell]$ still has
infinite order.

For the more general assertion no blowdown is needed.
If an automorphism preserves $E$, its derivative in the normal
direction gives
\[
 t_b^*N_{E/S}\simeq N_{E/S}.
\]
For a line bundle $N$ of degree $d$ on an elliptic curve, the
homomorphism $b\mapsto t_b^*N\otimes N^{-1}$ is multiplication
by $-d$ under the usual identification $\operatorname{Pic}^0(E)=E$.
Indeed write $N=\OO_E(\sum_j n_j[p_j])$ with
$\sum_j n_j=d$; the Abel--Jacobi class of its translated divisor
minus the original one is $-d b$.
Thus preservation of $N$ forces $d b=0$.
Taking a power proves the finite-cycle version, since its composite
preserves one axis and its nonzero-degree normal line.
In particular a nonzero-degree parabolic axis cannot carry the
marked generic period translation even if no finite Hopf
modification is specified. The original finite-support argument
also covers modifications where that normal-degree test alone
would give no information.

We check the actual cusp rather than assume that torsion fails.
The periods of A.1.2 have the form
\[
 g_1:(Z,W)\longmapsto(q c_1 Z,c_2W),\qquad
 g_2:(Z,W)\longmapsto(c_3Z,r^{12}q c_4 W),
\]
where the $c_i$ are nonvanishing holomorphic units near the cusp.
Choose $W$ or $W^{-1}$ so that
$a=c_2^{\varepsilon}$ has modulus less than one,
$\varepsilon\in\{1,-1\}$.
This is possible at the actual marked cusp because its limiting
unit has modulus different from one, as computed in B.4.8.
For $q\ne0$ small, the partial $g_1$ quotient of $(\C^*)^2$
is the dense core of the diagonal Hopf surface with multipliers
$(q c_1,a)$. The second period induces the above $h$, with
\[
 b(q,r)=(r^{12}q c_4(q,r))^{\varepsilon}.
\]
Fix a sufficiently small nonzero $q$. For each $k>0$ and
$j\in\Z$, the equation $b(q,r)^k=a(q)^j$ has a discrete
solution set in the punctured $r$-disc: its left-hand side has
a zero or pole of nonzero order at $r=0$, while the right-hand
side is a nonzero constant in $r$.
Hence outside a countable union of discrete sets, $[b(q,r)]$
has infinite order. Under a finite base change in $r$ the
nonzero order is multiplied by its ramification index and
the same conclusion holds. A finite base change in $q$ does
not affect this fixed-nonzero-$q$ argument.

To compare with B.4.10, write $q c_1=\kappa(s)s^n$ after
the chosen finite cover in the cusp variable. Its generic Kato
fibre is precisely a finite Hopf modification with nonempty
axis centre $[-s]$. A holomorphic extension of the actual second
period to that whole fibre would therefore contradict the result
just proved. Allowing finitely many translated copies does not
remove the contradiction, by the finite-collection assertion.
\end{proof}

\begin{remark}[What an infinite unfolding must additionally prove]
There is no corresponding obstruction to merely indexing translated
modifications by all integers and letting the deck map shift the
index. That construction alone has the wrong core. Indeed suppose
$p_j:Y_j\to H$ and
$\widetilde h_j:Y_j\to Y_{j+1}$ satisfy
$p_{j+1}\widetilde h_j=hp_j$.
On the disjoint union of their unmodified cores, the quotient by
the shift is identified with the single core $H^\circ$ by
\[
 [(j,y)]\longmapsto h^{-j}p_j(y).
\]
It is not the required quotient $H^\circ/\langle h\rangle$.
Thus a useful unfolding must have additional connected overlaps
and must recover the original second-period identifications there.

Nor can one simply glue several full compact Hopf modifications
along the same dense core and assume a Hausdorff space results.
If two embedded compact spaces in a Hausdorff space share an open
subset dense in both, their images, being closed, both equal its
closure. They coincide. A construction by genuinely distinct
compact copies with only that common core consequently cannot be
such a Hausdorff gluing.
These observations leave open connected unfoldings of noncompact
pieces with a new boundary geometry. Their analytic transition
maps, proper discontinuity, correct punctured quotient and compact
central quotient would all have to be proved; none follows from
the auxiliary compact family in B.4.10.
\end{remark}

% Owner: boundary_d.  Reserved B.4.14; integration after independent referee.
\begin{proposition}[B.4.14 (Proved): tangent weights and the compact-bundle obstruction reduction]
For the explicit parabolic surface $S=S_8(a,\kappa)$ and its
elliptic contraction $\pi:S\to Y$ of B.4.11, the genuine action
$\lambda\cdot(x,v)=(\lambda x,v)$ induces trivial representations
on the following cohomology groups:
\[
 H^0(S,\Theta_S(-\log E))=\C,\qquad
 H^1(S,\Theta_S(-\log E))=\C^9,
\]
\[
 H^0(Y,\Theta_Y)=\C,\qquad H^1(Y,\Theta_Y)=\C^8.
\]
The two second cohomology groups vanish.

Let $p:Z_L\to\PP^1$ be the compact holomorphic $Y$-bundle
associated to any holomorphic line bundle $L$ and this
$\C^*$ action.  Put
$\mathcal T=\Theta_{Z_L/\PP^1}$.
For every $b\in\PP^1$ one has
\[
 R^0p_*\mathcal T=\mathcal O_{\PP^1},\qquad
 R^1p_*\mathcal T=\mathcal O_{\PP^1}^{\oplus8},\qquad
 R^2p_*\mathcal T=0,
\]
and
\[
 H^2(Z_L,\mathcal T)=0,
 \qquad
 H^2\bigl(Z_L,\mathcal T\otimes
                   p^*\mathcal O_{\PP^1}(-b)\bigr)=0.
\]
These are cohomological reductions.  They do not, by themselves,
globalize a prescribed deformation along the nonisolated
singular section of $Z_L$ or construct a compact third component.
\end{proposition}
\begin{proof}
The action commutes with the central Kato gluing and hence
acts on the whole compact surface and its contraction.
Its toric coordinate description is also given in
\citet[Section~2, equations~(1)--(8)]{Fujiki2009D}.
For any $\lambda\ne1$, its fixed locus consists of the elliptic
curve $E$ and the eight nodes of the rational cycle.
To check that there are no further points, on the dense torus
an equality with a Kato deck iterate would require $a^jv=v$;
since $0<|a|<1$, this forces $j=0$, and then $\lambda=1$.
The action on each rational component is nontrivial and fixes
only its two endpoints.  In the unfolded fan the consecutive
rays are $(j,1)$ and $(j+1,1)$.
Their dual characters evaluate on the scaling cocharacter
$(1,0)$ as $-1,+1$.
Thus the two tangent weights at every cycle node are
$\lambda^{-1},\lambda$.

We also record a direct check of the space of global vector
fields, agreeing with \citet[Theorem~1.1]{Fujiki2009D}.
Write the contracted Kato germ as
\[
 F(x,v)=(Kxv^8,av),\qquad K=\kappa a^8.
\]
A global vector field restricts on the outer spherical annulus
to $P(x,v)\partial_x+Q(x,v)\partial_v$.
Hartogs extension gives holomorphic coefficients on the ball.
Compatibility with the collar gluing, followed by analytic
continuation, gives the exact germ identity
\[
 V(F)=DF\cdot V.
\]
Expand $Q(x,v)=\sum_{i\ge0}x^iQ_i(v)$.
The second component says
\[
 K^iv^{8i}Q_i(av)=aQ_i(v).
\]
For $i>0$ the lowest nonzero $v$-order on the two sides would
differ by $8i$, so $Q_i=0$.
For $i=0$, power-series coefficients give $Q_0(v)=cv$.
The constant-in-$x$ coefficient of the first component also
vanishes by the same lowest-order argument.
Its linear coefficient satisfies
$P_1(av)=P_1(v)+8c$.
Evaluation at zero gives $c=0$, and then $P_1$ is constant.
For $i>1$ the remaining equations are
\[
 P_i(v)=K^{i-1}v^{8(i-1)}P_i(av),
\]
so again $P_i=0$.
Consequently $H^0(S,\Theta_S)=\C\,x\partial_x$.
This Euler field is tangent to $E$ and has weight zero.

Set $\mathcal F=\Theta_S(-\log E)$.
B.4.11 proved $H^2(S,\mathcal F)=0$.
Apply the holomorphic Lefschetz formula to a root of unity
$\lambda\ne1$, using its finite cyclic group action.
The applicable original statement is
\citet[Theorem~(4.6), p.~566]{AtiyahSinger1968IIIE}; it requires
a compact complex manifold, not a K\"ahler manifold.
On $E$, the Atiyah extension $\mathcal F|_E$ has two
weight-zero terms and Chern character $2$; the normal line
has weight $\lambda$ and degree $-8$.
Since the tangent bundle of the elliptic curve has trivial
Todd class, the contribution of $E$ is
\[
 \int_E\frac{2}{1-\lambda^{-1}\exp(-c_1(N))}
   =-\frac{2\lambda^{-1}\deg N}
           {(1-\lambda^{-1})^2}
   =\frac{16\lambda}{(\lambda-1)^2}.
\]
At a cycle node $\mathcal F$ equals the ordinary tangent
bundle.  Its contribution is
\[
 \frac{\lambda+\lambda^{-1}}
      {(1-\lambda)(1-\lambda^{-1})}
   =-\frac{\lambda^2+1}{(\lambda-1)^2}.
\]
Adding the eight nodes gives the constant character
\[
 \sum_q(-1)^q\operatorname{Tr}
       (\lambda\mid H^q(S,\mathcal F))
 =\frac{16\lambda-8(\lambda^2+1)}{(\lambda-1)^2}
 =-8.
\]
Since $H^0$ is one weight-zero copy and $H^2=0$,
$H^1(S,\mathcal F)$ consists of nine trivial copies.
Equality on infinitely many roots of unity determines the
finite Laurent-polynomial character of the $\C^*$ representation.

Every map in the Leray sequence of B.4.11 is equivariant.
Its surjection
\[
 H^1(S,\mathcal F)\longrightarrow
 H^0(Y,R^1\pi_*\mathcal F)
       \simeq H^1(E,\Theta_E)=\C
\]
has a trivial one-dimensional target.
The kernel is $H^1(Y,\Theta_Y)$, so it is eight-dimensional
and trivial.  The equality in degree zero follows from
$\pi_*\mathcal F=\Theta_Y$; the second cohomology vanishing
was already proved in B.4.11.

Finally form $Z_L$ by choosing line-bundle transition functions
and gluing products $Y\times U_i$ by their action on $Y$.
This is a holomorphic locally trivial bundle with compact
fibre over the compact curve $\PP^1$, hence is proper and
compact.  It has the section of isolated cone points as its
singular locus; that locus is one dimensional in the total
space.  On each product chart, coherent cohomology and base
change identify $R^qp_*\mathcal T$ with
$H^q(Y,\Theta_Y)\otimes\mathcal O_{U_i}$.
The gluing is exactly by the representations computed above.
Their triviality proves the displayed direct-image formulas.
The projection formula and the Leray sequence then give
\[
 H^2\bigl(Z_L,\mathcal T\otimes p^*M\bigr)
      \simeq H^1(\PP^1,M^{\oplus8})
\]
for $M=\mathcal O$ or $\mathcal O(-b)$, both of which have
zero first cohomology on $\PP^1$.
This proves the asserted reduction, including the twist
appropriate to retaining one whole fibre.

The singular locus here is not isolated, so the isolated
criterion of B.4.11 is not applicable to $Z_L$.
In particular its $H^2$-vanishing is not a proof of analytic
gluing of a prescribed local smoothing along the whole
singular section.  Compatibility of local lifts and their
analytic effectivity remain separate requirements.
The exact character and ray calculations are recorded in
\texttt{code/ab\_push\_kato\_weights.py}.
\end{proof}


\begin{problem}[B.5.1; Open]\label{push:B51}
For an actual corrected filling in B.4.1,
compute the integral topology of its
central fibre and of a descent resolution,
the nearby-cycle Euler count, and the
applicable monodromy or Hodge comparison.
These invariants are not assigned to the
incompatible assembly refuted in B.3.1.
\end{problem}
